%% 大人用『物理を理解するための数学　ガイドブック　保護者・先生のための本』親ファイル（第16回編集委員会、組 G-A）
%% vol=0（巻の色 bkVol0＝紺灰色）。ルビなし。大人向けの通常の漢字使い（10月30日 10:15 依頼者の指示）。表紙に「大人用」の札（四角い形）。
%% 表1・表2（全9巻の構成）・とびら（p.1）・もくじ（p.2〜3）・本文（p.4〜）・奥付・表4。
%% 本文：guide-a1（G-B、第1章）→ guide-p2a（G-B、第2〜10章）→ guide-p2b（G-C、付録A・B）→ guide-p2a-src（G-B、出典）。
%% 第2章以降は前の章に続けて組む（改ページは各ファイルの側で）。第1章の頭だけ、ここで偶数ページにそろえる。
%% もくじは \pageref で拾う：kgd-ch1、kgb-ch2〜kgb-ch10、kgc-a・kgc-b、kgb-src。共通の部品は guide-kga.tex。
%% 組み方：cd tex && lualatex guide-adult.tex（2回）
\documentclass[12pt]{ltjsarticle}
\def\bkfontpath{fonts/}
\usepackage[vol=0]{bk-b5}
\usepackage{bk-matter}
\usepackage{bk-honbun}
\usepackage{bk-volume}
\usepackage{bk-pict1}
%% bk-motifs.tex  表紙の図柄（各巻の数学の考え方を絵にしたもの）
%% 描く範囲：x 0〜150, y 0〜106（mm）。色 mA＝主（白か濃い色）, mB＝従（淡い色か白）, mC＝影
\makeatletter
\tikzset{mot/.style={line cap=round, line join=round}}

%% 第1巻：10のまとまり・ものさし・時計
\newcommand\bk@apple[2]{\bkapple{(#1,#2)}{4.3}{bkAppleRed}}% リンゴ（へた・葉つき）
\newcommand\bkMotifi{%
  % 10のわく：白いマスにリンゴを1こずつ
  \foreach \i in {0,...,4} \foreach \j in {0,1} {
    \fill[white, rounded corners=1.6mm] (2.6+\i*14.4,48.6+\j*14.4) rectangle ++(13.2,13.2);
    \bk@apple{9.2+\i*14.4}{54.6+\j*14.4}}
  % のこりの3こは、白いおさらの上に
  \bkplate{(101,56)}{38}
  \foreach \i in {0,1,2} \bk@apple{89+\i*12}{61.5}
  \node[mA, font=\sffamily\bfseries\fontsize{22}{22}\selectfont] at (38,39) {10};
  \node[mA, font=\sffamily\bfseries\fontsize{22}{22}\selectfont] at (101,42) {3};
  % 時計（3時）
  \fill[mB] (130,88) circle (15.5);
  \draw[mA, line width=1.6pt] (130,88) circle (15.5);
  \foreach \a in {0,30,...,330} \draw[mC, line width=1pt] (130,88) ++(\a:12.6) -- ++(\a:2.2);
  \draw[mC, line width=2pt, line cap=round] (130,88) -- ++(0:8);
  \draw[mC, line width=1.4pt, line cap=round] (130,88) -- ++(90:11.5);
  \fill[mC] (130,88) circle (1.3);
  % ものさし（0〜10）
  \fill[mB, rounded corners=1.6mm] (2,6) rectangle (122,24);
  \foreach \k in {0,...,20} {\pgfmathsetmacro\bkx{7+\k*5.5}
    \ifodd\k \draw[mC, line width=0.7pt] (\bkx,24) -- (\bkx,20.5);
    \else \draw[mC, line width=0.9pt] (\bkx,24) -- (\bkx,18.5);\fi}
  \foreach \k in {0,...,10} \node[mC, font=\sffamily\bfseries\fontsize{7.5}{8}\selectfont] at (7+\k*11,13.5) {\k};
}

%% 第2巻：落ちるボール（同じ時間ごと）・倍々・単位・数直線
\newcommand\bkMotifii{%
  \draw[mA, line width=1.4pt] (5,2) -- (5,104);
  \foreach \y in {2,12.2,...,104} \draw[mA, line width=0.9pt] (5,\y) -- (8.5,\y);
  \foreach \t/\y in {0/100, 1/93.6, 2/74.5, 3/42.7, 4/-1.9} {
    \fill[mA] (16,{\y+3}) circle (3.0);
    \node[anchor=west, mA, font=\sffamily\bfseries\fontsize{8.5}{10}\selectfont] at (20.5,{\y+3}) {\t\R{秒}{びょう}};}
  % 単位のふだ
  \foreach \x/\t in {92/m, 112/kg, 132/s} {
    \fill[mB, rounded corners=2mm] (\x,88) rectangle ++(17,13);
    \node[mC, font=\rmfamily\bfseries\fontsize{17}{17}\selectfont] at (\x+8.5,94.5) {\t};}
  % 倍々のブロック（1, 2, 4, 8）
  \begin{scope}[shift={(41,44)}]
    \foreach \bx/\bw/\bh/\bn in {0/1/1/1, 19.5/2/1/2, 47.5/2/2/4, 75.5/4/2/8} {
      \fill[mA] (\bx,0) rectangle ++(\bw*8.5,\bh*8.5);
      \draw[mC, line width=0.7pt, shift={(\bx,0)}] (0,0) grid[step=8.5] ({\bw*8.5},{\bh*8.5});
      \node[mA, font=\sffamily\bfseries\fontsize{16}{16}\selectfont] at (\bx+\bw*4.25,\bh*8.5+6.5) {\bn};}
    \foreach \xm in {14, 42, 70} \node[mB, font=\sffamily\bfseries\fontsize{11}{11}\selectfont] at (\xm,4.25) {×2};
  \end{scope}
  % 数直線（−5〜5）
  \draw[mA, line width=1.6pt, -{Stealth[length=8pt]}] (40,14) -- (150,14);
  \foreach \k in {-5,...,5} {\pgfmathsetmacro\bkx{93+\k*10}
    \draw[mA, line width=1.1pt] (\bkx,11.5) -- (\bkx,16.5);
    \node[mA, font=\sffamily\bfseries\fontsize{9.5}{9.5}\selectfont] at (\bkx,6.5) {\ifnum\k<0 −\number-\k\else\k\fi};}% 第7回：数字をゴシックに（マイナスは記号の字）
  \draw[mB, line width=3.4pt, -{Stealth[length=9pt]}] (93,24) -- (123,24);
  \draw[mB, line width=3.4pt, -{Stealth[length=9pt]}] (93,24) -- (73,24);
  \fill[mA] (93,24) circle (2.4);
}

%% 第3巻：データに直線・投げたボールの道・ヒストグラム
\newcommand\bkMotifiii{%
  \draw[mA, line width=1.3pt, -{Stealth[length=7pt]}] (4,6) -- (98,6);
  \draw[mA, line width=1.3pt, -{Stealth[length=7pt]}] (6,4) -- (6,104);
  \draw[mB, line width=4pt] (10,12) -- (94,86);
  \foreach \x/\e in {15/3, 24/-4, 33/2.5, 42/-2, 51/4, 60/-3, 69/1.5, 78/-1.5, 87/2.5}
    \fill[mA] (\x,{12+(\x-10)*74/84+\e}) circle (2.3);
  \draw[mA, line width=1.3pt, dash pattern=on 4pt off 3pt] plot[domain=104:148, samples=40] (\x,{8+74*(1-((\x-126)/22)^2)});
  \foreach \x in {104,111.33,118.67,126,133.33,140.67,148} \fill[mA] (\x,{8+74*(1-((\x-126)/22)^2)}) circle (2.4);
  \foreach \h [count=\i] in {2,5,10,15,18,15,10,5,2}
    \fill[mB] (99+\i*5.4,88) rectangle ++(4.8,\h);
  \draw[mA, line width=1pt] (103,88) -- (150,88);
}

%% 第4巻：円運動の影＝波・ベクトルの合成
\newcommand\bkMotifiv{%
  \def\bkcy{64}
  \draw[mB, line width=0.9pt] (2,\bkcy) -- (66,\bkcy) (34,\bkcy-32) -- (34,\bkcy+32);
  \draw[mA, line width=1.6pt] (34,\bkcy) circle (28);
  \fill[mB] (34,\bkcy) -- (42,\bkcy) arc (0:50:8) -- cycle;
  \draw[mA, line width=2pt] (34,\bkcy) -- ++(50:28);
  \draw[mC, line width=1.4pt, dash pattern=on 3pt off 2pt] (52,{\bkcy+28*sin(50)}) -- (52,\bkcy);
  \draw[mC, line width=3pt] (34,\bkcy) -- (52,\bkcy);
  \fill[mA] (52,{\bkcy+28*sin(50)}) circle (2.4);
  \draw[mA, line width=1pt, dash pattern=on 3pt off 2pt] (54,{\bkcy+28*sin(50)}) -- (72,{\bkcy+28*sin(50)});
  \draw[mB, line width=0.9pt] (72,\bkcy) -- (150,\bkcy);
  \draw[mA, line width=2.2pt] plot[domain=72:150, samples=120, smooth] (\x,{\bkcy+28*sin(50+(\x-72)*5.142857)});
  % ベクトル
  \begin{scope}[shift={(92,4)}]
    \draw[mB, line width=1pt, dash pattern=on 3pt off 2pt] (32,7) -- (44,22) -- (12,15);
    \draw[mA, line width=2pt, -{Stealth[length=8pt]}] (0,0) -- (32,7);
    \draw[mA, line width=2pt, -{Stealth[length=8pt]}] (0,0) -- (12,15);
    \draw[mC, line width=2.6pt, -{Stealth[length=9pt]}] (0,0) -- (44,22);
  \end{scope}
}

%% 第5巻：接線（微分）と細い長方形（積分）
\newcommand\bkMotifv{%
  \draw[mA, line width=1.3pt, -{Stealth[length=7pt]}] (4,6) -- (150,6);
  \draw[mA, line width=1.3pt, -{Stealth[length=7pt]}] (6,4) -- (6,104);
  \foreach \k in {0,...,7} {\pgfmathsetmacro\bkx{16+\k*10}
    \pgfmathsetmacro\bkh{16+0.35*((\bkx-6)/10)^2+9*sin(3*(\bkx-6))}
    \fill[mB] (\bkx,6) rectangle (\bkx+10,\bkh);
    \draw[mC, line width=0.6pt] (\bkx,6) rectangle (\bkx+10,\bkh);}
  \draw[mA, line width=2.4pt] plot[domain=8:146, samples=120, smooth] (\x,{16+0.35*((\x-6)/10)^2+9*sin(3*(\x-6))});
  \pgfmathsetmacro\bky{16+0.35*11.4^2+9*sin(3*114)}
  \pgfmathsetmacro\bks{0.007*114+9*3*pi/180*cos(3*114)}
  \draw[mA, line width=1.2pt, dash pattern=on 3pt off 2pt] (120,\bky) -- (134,\bky) -- (134,\bky+14*\bks);% 第6回：緑地で見えるように白に
  \draw[mA, line width=1.8pt] (100,{\bky-20*\bks}) -- (146,{\bky+26*\bks});
  \fill[mA] (120,\bky) circle (2.6);
  \node[mA, anchor=north, font=\sffamily\bfseries\fontsize{10}{10}\selectfont] at (127,\bky-1.8) {$\Delta x$};% 第6回：横の辺の中央の下
  \node[mA, anchor=west, font=\sffamily\bfseries\fontsize{10}{10}\selectfont] at (135,\bky+8) {$\Delta y$};
}

%% 第6巻：ゆれながら止まる（減衰振動）・回る矢印（複素数）
\newcommand\bkMotifvi{%
  \draw[mB, line width=0.9pt] (6,56) -- (150,56);
  \draw[mA, line width=1.3pt, -{Stealth[length=7pt]}] (6,6) -- (6,106);
  \draw[mB, line width=1.2pt, dash pattern=on 4pt off 3pt] plot[domain=8:148, samples=80] (\x,{56+40*exp(-0.022*(\x-8))});
  \draw[mB, line width=1.2pt, dash pattern=on 4pt off 3pt] plot[domain=8:148, samples=80] (\x,{56-40*exp(-0.022*(\x-8))});
  \draw[mA, line width=2.4pt] plot[domain=8:148, samples=200, smooth] (\x,{56+40*exp(-0.022*(\x-8))*cos(9.17*(\x-8))});
  \begin{scope}[shift={(132,90)}]
    \draw[mB, line width=0.8pt] (-15,0) -- (15,0) (0,-14) -- (0,15);
    \draw[mA, line width=1.2pt] (0,0) circle (11);
    \draw[mA, line width=2pt, -{Stealth[length=7pt]}] (0,0) -- (40:11);
    \draw[mB, line width=1.4pt, -{Stealth[length=6pt]}] (55:13.5) arc (55:120:13.5);
  \end{scope}
}

%% 第7巻：行列で変わる格子
\newcommand\bkMotifvii{%
  \begin{scope}
    \clip (0,0) rectangle (150,106);
    \begin{scope}[shift={(72,50)}]
      \foreach \i in {-7,...,7} {
        \draw[mB, line width=0.5pt] (\i*12,-60) -- (\i*12,60);
        \draw[mB, line width=0.5pt] (-90,\i*12) -- (90,\i*12);}
      \begin{scope}[cm={1,0.28,0.48,1,(0,0)}]
        \foreach \i in {-8,...,8} {
          \draw[mA, line width=0.9pt] (\i*12,-80) -- (\i*12,80);
          \draw[mA, line width=0.9pt] (-110,\i*12) -- (110,\i*12);}
        \fill[mC] (0,0) rectangle (12,12);
        \draw[mA, line width=2.4pt, -{Stealth[length=9pt]}] (0,0) -- (24,0);
        \draw[mA, line width=2.4pt, -{Stealth[length=9pt]}] (0,0) -- (0,24);
      \end{scope}
      \fill[mA] (0,0) circle (2);
    \end{scope}
  \end{scope}
}

%% 第8巻：わき出し（div）と うず（rot）
\newcommand\bkMotifviii{%
  % うすい等高線と、いちばん急な向き（grad）
  \foreach \r in {6,12,...,30} \draw[mB, line width=0.5pt] (75,94) ellipse ({\r*1.5} and {\r*0.3});
  % わき出し
  \begin{scope}[shift={(36,50)}]
    \fill[mB] (0,0) circle (5);\fill[mA] (0,0) circle (2.4);
    \foreach \a in {0,30,...,330} \draw[mA, line width=1.6pt, -{Stealth[length=5pt]}] (\a:9) -- (\a:16);
    \foreach \a in {15,45,...,345} \draw[mA, line width=1.3pt, -{Stealth[length=5pt]}] (\a:21) -- (\a:28);
    \node[mB, font=\sffamily\bfseries\fontsize{14}{14}\selectfont] at (0,-38) {div};
  \end{scope}
  % うず
  \begin{scope}[shift={(112,50)}]
    \fill[mB] (0,0) circle (2.4);
    \foreach \a in {0,45,...,315} \draw[mA, line width=1.6pt, -{Stealth[length=5pt]}] (\a:11) ++(\a-90:2.6) -- ++(\a+90:5.2);
    \foreach \a in {15,45,...,345} \draw[mA, line width=1.3pt, -{Stealth[length=5pt]}] (\a:24.5) ++(\a-90:3.6) -- ++(\a+90:7.2);
    \node[mB, font=\sffamily\bfseries\fontsize{14}{14}\selectfont] at (0,-38) {rot};
  \end{scope}
}

%% 第9巻：sinをたして四角い波・スペクトル
\newcommand\bkMotifix{%
  \def\bkmid{72}\def\bkamp{15}
  \draw[mB, line width=1.2pt, dash pattern=on 4pt off 3pt] (8,\bkmid-\bkamp) -- (8,\bkmid+\bkamp) -- (43,\bkmid+\bkamp) -- (43,\bkmid-\bkamp) -- (78,\bkmid-\bkamp) -- (78,\bkmid+\bkamp) -- (113,\bkmid+\bkamp) -- (113,\bkmid-\bkamp) -- (148,\bkmid-\bkamp) -- (148,\bkmid);
  \draw[mB, line width=1.1pt] plot[domain=8:148, samples=120, smooth] (\x,{\bkmid+\bkamp*4/pi*sin((\x-8)*5.142857)});
  \draw[mA, line width=2.2pt] plot[domain=8:148, samples=320, smooth] (\x,{\bkmid+\bkamp*4/pi*(sin((\x-8)*5.142857)+sin(3*(\x-8)*5.142857)/3+sin(5*(\x-8)*5.142857)/5+sin(7*(\x-8)*5.142857)/7+sin(9*(\x-8)*5.142857)/9)});
  \draw[mA, line width=1.3pt] (8,8) -- (148,8);
  \foreach \k [count=\i] in {1,3,5,7,9,11} {
    \fill[mA] ({4+\i*18},8) rectangle ++(9,{36/\k});
    \node[mB, font=\sffamily\bfseries\fontsize{8}{8}\selectfont] at ({8.5+\i*18},4) {\k};}
}
\makeatother

%% 自動生成（gen.py）。直接編集しないこと
\makeatletter
\bkdefvol{1}{かぞえる・はかる}{}{かずで くらべる、かずで つたえる}{小学1年〜3年}{1}{1}
\bkdefvol{2}{\R{数}{かず}と\R{量}{りょう}のことば}{}{「5」ではなく「5\,m」。\R{数}{かず}と\R{量}{りょう}のことばを\R{身}{み}につけよう}{小学4年〜中学1年}{1}{2}
\bkdefvol{3}{\R{関数}{かん|すう}とデータ}{きまりを\R{見}{み}つける}{\R{表}{ひょう}・グラフ・\R{式}{しき}で、\R{変化}{へん|か}のきまりをつかもう}{中学1年〜3年}{1}{3}
\bkdefvol{4}{\R{形}{かたち}と\R{向}{む}き}{\R{図形}{ず|けい}・\R{三角関数}{さん|かく|かん|すう}・ベクトル}{長さと角度で、形と向きを数にしよう}{中学2年〜高校2年}{1}{5}
\bkdefvol{5}{\R{変化}{へん|か}をつかむ}{\R{微分}{び|ぶん}と\R{積分}{せき|ぶん}}{変わる速さと、\mbox{たまった量}を計算しよう}{高校1年〜2年}{2}{4}
\bkdefvol{6}{\R{未来}{み|らい}を\R{計算}{けい|さん}する}{\R{微分方程式}{び|ぶん|ほう|てい|しき}と\R{複素数}{ふく|そ|すう}}{「次の瞬間のルール」から、未来を計算しよう}{高校2年〜3年}{2}{5}
\bkdefvol{7}{まとめて\R{計算}{けい|さん}する}{\R{行列}{ぎょう|れつ}}{たくさんの数を、表のまま計算しよう}{高校3年（発展）〜大学1年}{3}{5}
\bkdefvol{8}{\R{空間}{くう|かん}に\R{広}{ひろ}がる\R{量}{りょう}}{ベクトル\R{解析}{かい|せき}}{空間に広がる量を、矢印で読みとこう}{大学1年〜2年}{3}{5}
\bkdefvol{9}{\R{波}{なみ}と\R{変化}{へん|か}を\R{読}{よ}みとく}{フーリエ\R{解析}{かい|せき}とラプラス\R{変換}{へん|かん}}{波を成分に分け、微分方程式をかけ算で解こう}{大学2年}{3}{5}
\@namedef{bk@rname@1}{\R{変化}{へん|か}}\@namedef{bk@rsub@1}{ひき\R{算}{ざん}→\R{傾}{かたむ}き→\R{微分}{び|ぶん}}\@namedef{bk@rend@1}{\R{流体}{りゅう|たい}・\R{熱}{ねつ}}
\@namedef{bk@rname@2}{たし\R{合}{あ}わせ}\@namedef{bk@rsub@2}{かけ\R{算}{ざん}→\R{面積}{めん|せき}→\R{積分}{せき|ぶん}}\@namedef{bk@rend@2}{\R{電磁気}{でん|じ|き}}
\@namedef{bk@rname@3}{\R{倍々}{ばい|ばい}}\@namedef{bk@rsub@3}{\R{累乗}{るい|じょう}→\R{指数}{し|すう}→$e$}\@namedef{bk@rend@3}{\R{制御}{せい|ぎょ}}
\@namedef{bk@rname@4}{\R{回転}{かい|てん}と\R{波}{なみ}}\@namedef{bk@rsub@4}{\R{角}{かく}→sin→\R{複素数}{ふく|そ|すう}}\@namedef{bk@rend@4}{\R{音}{おと}・\R{交流}{こう|りゅう}}
\@namedef{bk@rname@5}{\R{表}{ひょう}}\@namedef{bk@rsub@5}{\R{九九}{く|く}の\R{表}{ひょう}→\R{行列}{ぎょう|れつ}}\@namedef{bk@rend@5}{\R{回路網}{かい|ろ|もう}・\R{制御}{せい|ぎょ}}
\@namedef{bk@rname@6}{\R{矢印}{や|じるし}}\@namedef{bk@rsub@6}{\R{向}{む}き→ベクトル→\R{場}{ば}}\@namedef{bk@rend@6}{ロボット・\\\R{電磁気}{でん|じ|き}}
\@namedef{bk@rname@7}{はかる}\@namedef{bk@rsub@7}{\R{目}{め}もり→\R{誤差}{ご|さ}→\R{統計}{とう|けい}}\@namedef{bk@rend@7}{\R{計測}{けい|そく}・\R{診断}{しん|だん}}
\makeatother

%% 自動生成（gen.py）
\makeatletter
\def\bk@npart{3}
\def\bk@routepage{4}
\bkdefpart{1}{第1部}{フーリエ解析}{（9-1〜9-4）}
\bkdefpart{2}{第2部}{ラプラス変換}{（9-5〜9-9）}
\bkdefpart{3}{第3部}{行列と微分方程式}{（9-10〜9-12）}
\@namedef{bk@rhere@1}{（この\R{巻}{かん}には\R{駅}{えき}がありません）}
\@namedef{bk@rhere@2}{（この\R{巻}{かん}には\R{駅}{えき}がありません）}
\@namedef{bk@rhere@3}{9-5、9-6、9-7、9-8、9-9}
\@namedef{bk@rhere@4}{9-1、9-2、9-3}
\@namedef{bk@rhere@5}{9-10、9-11}
\@namedef{bk@rhere@6}{（この\R{巻}{かん}には\R{駅}{えき}がありません）}
\@namedef{bk@rhere@7}{9-4、9-12}
\def\bk@stations{\bkstation{1}{1}{1-5}\bkstation{1}{2}{2-3}\bkstation{1}{3}{3-5}\bkstation{1}{5}{5-1}\bkstation{1}{6}{6-1}\bkstation{1}{8}{8-4}\bkstation{2}{1}{1-21}\bkstation{2}{2}{2-3}\bkstation{2}{4}{4-5}\bkstation{2}{5}{5-10}\bkstation{2}{6}{6-3}\bkstation{2}{8}{8-11}\bkstation{3}{2}{2-7}\bkstation{3}{3}{3-11}\bkstation{3}{5}{5-14}\bkstation{3}{6}{6-4}\bkstation{3}{9}{9-5}\bkstation{4}{1}{1-20}\bkstation{4}{4}{4-8}\bkstation{4}{5}{5-15}\bkstation{4}{6}{6-10}\bkstation{4}{9}{9-1}\bkstation{5}{1}{1-13}\bkstation{5}{2}{2-18}\bkstation{5}{3}{3-6}\bkstation{5}{7}{7-1}\bkstation{5}{9}{9-10}\bkstation{6}{2}{2-2}\bkstation{6}{3}{3-2}\bkstation{6}{4}{4-10}\bkstation{6}{7}{7-7}\bkstation{6}{8}{8-5}\bkstation{7}{1}{1-11}\bkstation{7}{2}{2-6}\bkstation{7}{3}{3-14}\bkstation{7}{9}{9-4}}
\@namedef{bk@rtop@1}{182}
\@namedef{bk@rtop@2}{182}
\@namedef{bk@rtop@3}{165}
\@namedef{bk@rtop@4}{182}
\@namedef{bk@rtop@5}{182}
\@namedef{bk@rtop@6}{165}
\@namedef{bk@rtop@7}{182}
\makeatother

%% ルビを付けない（親字だけ）。数の区切りは第4巻以降と同じ（整数部分5けた以上）
\protected\def\R#1#2{#1}
\sisetup{group-digits=integer, group-minimum-digits=5}
\hypersetup{pdftitle={物理を理解するための数学　ガイドブック　保護者・先生のための本（大人用）}, pdfauthor={「物理を理解するための数学」編集委員会}}
\newcommand\kgaBookTitle{ガイドブック}
\newcommand\kgaBookSub{保護者・先生のための本}
\newcommand\kgaCatch{シリーズの概要と使い方、学習指導要領との対応}
\newcommand\kgaBandLabel{ガイドブック\hspace{0.6\zw}保護者・先生のための本}
%% guide-kga.tex　ガイドブック2冊（子ども用 guide-kids.tex・大人用 guide-adult.tex）で共通の部品（組 G-A、頭文字 kga）
%% 第16回編集委員会（10月30日 10:14 の依頼者の指示で2冊に分けた）。親ファイルの前付けと奥付・子ども用の本文で使う。
\makeatletter
%% 色：kgaPA＝子ども用の札の色（れんが色。白い地で約5.9:1）、kgaPB＝大人用の札の色（巻の色 bkVol0）
%% 札は色だけでなく、字（子ども用／大人用）と形（丸／四角）でも区別する
\definecolor{kgaPA}{cmyk}{0.05,0.65,0.95,0.28}
\colorlet{kgaPB}{bkVol0}
\newcommand\kgaNoRuby{\def\R##1##2{##1}}% ルビを外す
%% 本ごとに親ファイルで決めるもの（書名・副題・キャッチ・帯の小さな字・偶数ページの柱）
\providecommand\kgaBookTitle{ガイドブック}
\providecommand\kgaBookSub{}
\providecommand\kgaCatch{}
\providecommand\kgaBandLabel{}
%% ラベルのページ番号（まだ組んでいない所は「—」）  \kgaPage{ラベル}
\newcommand\kgaPage[1]{\@ifundefined{r@#1}{—}{\pageref{#1}}}

%% --------------------------------------------------------------------
%% 9冊の本のたなの絵  \kgaShelf{x}{y}{倍率}{棚板の色}：本の高さはページ数に比例（400〜640ページ）
%% --------------------------------------------------------------------
\newcommand\kgaShelf[4]{%
  \begin{scope}[shift={(#1,#2)}, scale=#3]
    \foreach \v/\p in {1/400,2/456,3/564,4/560,5/492,6/600,7/492,8/632,9/640} {
      \pgfmathsetmacro\kga@bx{(\v-1)*16.2}\pgfmathsetmacro\kga@bh{\p*0.13}
      \fill[bkVol\v, rounded corners=0.8mm] (\kga@bx,0) rectangle ++(14.4,\kga@bh);
      \draw[bkBandOn\v, line width=0.5pt, opacity=0.55] (\kga@bx+1.6,\kga@bh-6) -- ++(11.2,0) (\kga@bx+1.6,6) -- ++(11.2,0);
      \node[text=bkBandOn\v, font=\sffamily\bfseries\fontsize{15}{15}\selectfont] at (\kga@bx+7.2,\kga@bh-14) {\v};
    }
    \fill[#4, rounded corners=0.6mm] (-3,-3.2) rectangle (148.2,0);
  \end{scope}}

%% --------------------------------------------------------------------
%% 表紙の札  \kgaBadge{kids|adult}{中心x}{中心y}{大きさ（1＝表紙）}
%%   子ども用：白い丸に、れんが色の太い輪と「子ども用」。大人用：白い四角に、紺灰色の二重の枠と「大人用」
%% --------------------------------------------------------------------
\newcommand\kgaBadge[4]{%
  \begin{scope}[shift={(#2,#3)}, scale=#4, transform shape]
  \def\kga@b{#1}\def\kga@kids{kids}%
  \ifx\kga@b\kga@kids
    \fill[white] (0,0) circle (14);
    \draw[kgaPA, line width=2.4pt] (0,0) circle (12.6);
    \node[text=kgaPA, font=\sffamily\bfseries\fontsize{15}{17}\selectfont, inner sep=0pt] at (0,0) {\R{子}{こ}ども\R{用}{よう}};
  \else
    \fill[white] (-17,-11) rectangle (17,11);
    \draw[kgaPB, line width=1.6pt] (-15.4,-9.4) rectangle (15.4,9.4);
    \draw[kgaPB, line width=0.6pt] (-13.6,-7.6) rectangle (13.6,7.6);
    \node[text=kgaPB, font=\sffamily\bfseries\fontsize{14}{16}\selectfont, inner sep=0pt] at (0,0) {大人用};
  \fi
  \end{scope}}

%% --------------------------------------------------------------------
%% 表1（表紙）：\bkFrontArt と同じ配置。巻の番号の札のかわりに「子ども用」「大人用」の札
%%   \kgaFrontArt{kids|adult}
%% --------------------------------------------------------------------
\newcommand\kgaFrontArt[1]{%
  \fill[bkVol] (-3,64) rectangle (185,260);
  \fill[white] (-3,-3) rectangle (185,64);
  \node[anchor=base west, inner sep=0pt, text=bkOn, font=\sffamily\bfseries\fontsize{12.5}{15}\selectfont] at (16,242) {\bkSeries};
  \draw[bkOn, line width=0.6pt] (16,237.5) -- (112,237.5);
  \kgaBadge{#1}{147}{236}{1}
  \node[anchor=base west, inner sep=0pt, text=bkOn, font=\bk@titlefont] at (15,206) {\kgaBookTitle};
  \node[anchor=base west, inner sep=0pt, text=bkOn, font=\sffamily\bfseries\fontsize{17}{20}\selectfont] at (16,192) {\kgaBookSub};
  \node[anchor=base west, inner sep=0pt, text=bkOn, font=\sffamily\fontsize{11.5}{14}\selectfont] at (16,181.5) {\kgaCatch};
  \foreach \r in {1,...,7} {\node[inner sep=0pt] at ({20+(\r-1)*11.5},164) {\bkrouteicon[7.5mm]{\r}};}
  \kgaShelf{18}{74}{1}{white}
  \bk@lineup{16}{36}
  \node[anchor=base, text=bkGray, font=\sffamily\fontsize{8}{10}\selectfont] at (91,16) {\R{全}{ぜん}9\R{巻}{かん}　\bkSeries};
  \node[anchor=base, text=bkGray, font=\sffamily\fontsize{7.5}{9}\selectfont] at (91,8.5) {\bkEditor\hspace{1pt}\R{編}{へん}};
}

%% --------------------------------------------------------------------
%% 表4（裏表紙）  \kgaBackArt{見出し}{中身（TikZ。y=200 から下へ）}{ふきだしの文}
%% --------------------------------------------------------------------
\newcommand\kgaBackArt[3]{%
  \fill[bkVol!9] (-3,-3) rectangle (185,260);
  \fill[bkVol] (-3,226) rectangle (185,260);
  \node[anchor=base west, inner sep=0pt, text=bkOn, font=\sffamily\fontsize{9.5}{12}\selectfont] at (16,246) {\bkSeries};
  \node[anchor=base west, inner sep=0pt, text=bkOn, font=\sffamily\bfseries\fontsize{19}{24}\selectfont] at (16,233)
    {\kgaBookTitle\hspace{0.5em}{\fontsize{12}{14}\selectfont\kgaBookSub}};
  \node[anchor=base west, inner sep=0pt, text=bkDeep, font=\sffamily\bfseries\fontsize{13}{16}\selectfont] at (16,210) {#1};
  \draw[bkVol, line width=0.8pt] (16,206) -- (166,206);
  #2
  \node[anchor=base west, inner sep=0pt, text=bkDeep, font=\sffamily\bfseries\fontsize{13}{16}\selectfont] at (16,140) {\R{全}{ぜん}9\R{巻}{かん}のならび};
  \draw[bkVol, line width=0.8pt] (16,136) -- (98,136);
  \foreach \i in {1,...,9} {
    \pgfmathsetmacro\bk@ly{127-(\i-1)*8.6}
    \fill[bkVol\i, rounded corners=0.8mm] (16,\bk@ly-2.8) rectangle ++(5.6,5.6);
    \node[text=bkBandOn\i, font=\sffamily\bfseries\fontsize{8}{8}\selectfont] at (18.8,\bk@ly) {\i};
    \node[anchor=west, inner sep=0pt, text=bkInk, font=\sffamily\fontsize{9.5}{11}\selectfont] at (24.5,\bk@ly-0.3)
      {\ltjsetruby{stretch=001}\bkTitle{\i}};
  }
  \node[rectangle callout, callout absolute pointer={(136.5,91)}, anchor=south, draw=bkVol, line width=0.8pt, fill=white,
    rounded corners=3mm, inner xsep=3.5mm, inner ysep=2.5mm, text width=50mm, align=flush center, text=bkDeep,
    font=\sffamily\bfseries\fontsize{10.5}{16}\selectfont] at (136.5,97) {#3};
  \fill[bkVol!22, rounded corners=2mm] (107,50) rectangle (166,54);
  \bkChar{sora}{(117,53.5)}{31mm}\bkChar{mina}{(136.5,53.5)}{31mm}\bkChar{delta}{(156,53.5)}{19mm}
  \draw[bkGray, line width=0.6pt, dash pattern=on 2pt off 1.5pt, fill=white] (114,12) rectangle (166,38);
  \node[text=bkGray, align=flush center, font=\sffamily\fontsize{7}{9}\selectfont] at (140,25) {ISBN・バーコード\\（入稿時に配置）};
  \node[anchor=base west, inner sep=0pt, text=bkGray, font=\sffamily\fontsize{7.5}{10}\selectfont] at (16,25) {\bkSeries　\kgaBookTitle};
  \node[anchor=base west, inner sep=0pt, text=bkGray, font=\sffamily\fontsize{7.5}{10}\selectfont] at (16,16) {\bkEditor\hspace{1pt}\R{編}{へん}};
}
% 表4の中身の1行（札つき）  \kgaBackItem{y}{札の文字}{文}
\newcommand\kgaBackItem[3]{%
  \node[anchor=west, fill=bkVol, text=bkOn, rounded corners=1.2mm, inner xsep=1.6mm, inner ysep=1mm, minimum width=7mm,
    font=\sffamily\bfseries\fontsize{8.5}{10}\selectfont] (kgabi) at (16,#1) {#2};
  \node[anchor=west, inner sep=0pt, text=bkInk, font=\sffamily\fontsize{10}{13}\selectfont] at ([xshift=2.5mm]kgabi.east) {#3};}

%% --------------------------------------------------------------------
%% とびら（p.1）  \kgaTitlePage{kids|adult}
%% --------------------------------------------------------------------
\newcommand\kgaTitleArt[1]{%
  \fill[bkVol] (-3,246) rectangle (185,260);
  \node[anchor=base, text=bkDeep, font=\sffamily\bfseries\fontsize{12}{15}\selectfont] at (91,206) {\bkSeries};
  \kgaBadge{#1}{91}{188}{0.85}
  \node[anchor=base, text=bkDeep, font=\sffamily\bfseries\fontsize{32}{40}\selectfont] at (91,158) {\kgaBookTitle};
  \node[anchor=base, text=bkDeep, font=\sffamily\bfseries\fontsize{15}{18}\selectfont] at (91,145) {\kgaBookSub};
  \node[anchor=base, text=bkGray, font=\fontsize{11}{15}\selectfont] at (91,130) {\kgaCatch};
  \kgaShelf{52.5}{78}{0.5}{bkVol}
  \node[anchor=base, text=bkGray, font=\sffamily\fontsize{9}{12}\selectfont] at (91,34) {\bkEditor\hspace{1pt}\R{編}{へん}};
  \fill[bkVol] (-3,-3) rectangle (185,4);
}
\newcommand\kgaTitlePage[1]{%
  \clearpage\setcounter{page}{1}%
  \begingroup\global\bk@decofalse\bkfolio{false}%
  \AddToHookNext{shipout/background}{\put(0,-\paperheight){%
    \begin{tikzpicture}[overlay, x=1mm, y=1mm]\kgaTitleArt{#1}\end{tikzpicture}}}%
  \null\clearpage\endgroup\global\bk@decotrue\bkfolio{true}}

%% --------------------------------------------------------------------
%% もくじ
%% --------------------------------------------------------------------
% 行  \kgaTocRow{色}{番号}{題名}{ひとこと（空なら1行）}{ラベル}
\newcommand\kgaTocRow[5]{%
  \par\noindent
  \parbox[t]{12mm}{\leavevmode\tikz[baseline=(p.base)]\node[fill=#1, text=white, rounded corners=1.3mm,
      minimum width=9mm, inner xsep=1mm, inner ysep=1.1mm, font=\sffamily\bfseries\fontsize{9.5}{10}\selectfont] (p) {#2};}%
  \parbox[t]{\dimexpr\linewidth-12mm-16mm\relax}{\setlength\parskip{0pt}%
    \leavevmode{\sffamily\bfseries\fontsize{11}{17}\selectfont\color{bkInk}#3}\if\relax\detokenize{#4}\relax\else\par\nobreak\vspace{-0.6mm}
    {\rmfamily\fontsize{9.5}{15}\selectfont\color{bkInk}#4}\fi}%
  \parbox[t]{16mm}{\raggedleft\leavevmode{\sffamily\bfseries\fontsize{11}{17}\selectfont\color{bkInk}\bkpageno{\kgaPage{#5}}}}%
  \par\vspace{1.2mm}{\color{bkLine}\hrule height 0.35pt}\vspace{1.2mm}}
% 見出し  \kgaTocHead
\newcommand\kgaTocHead{%
  \noindent\begin{tikzpicture}
    \node[anchor=base west, inner sep=0pt, text=bkDeep, font=\sffamily\bfseries\fontsize{30}{36}\selectfont] (h) at (0,0) {もくじ};
    \node[anchor=base west, inner sep=0pt, text=bkDeep, font=\sffamily\bfseries\fontsize{12}{15}\selectfont] at ([xshift=6mm]h.base east) {\kgaBookTitle\hspace{0.6em}\kgaBookSub};
    \draw[bkVol, line width=1.6pt, overlay] (0,-3.5mm) -- (\linewidth,-3.5mm);
  \end{tikzpicture}\par\vspace{6mm}}
% \kgaTocPages{1ページめ}{2ページめ（空ならもくじは1ページ）}
\newcommand\kgaTocPages[2]{%
  \clearpage\setcounter{page}{2}%
  \global\bk@decofalse\bkfolio{true}\bksetrunning{}{}{}%
  \kgaTocHead #1\par\clearpage
  \if\relax\detokenize{#2}\relax\else #2\par\clearpage\fi
  \global\bk@decotrue}

%% --------------------------------------------------------------------
%% 色の面の上に問いを置いたページ（子ども用の「はじめに」、p.3）。飾り・ノンブルなし
%%   \kgaArtPage{大きな字}{ラベル}{題名}{問い y/文,…}{絵}{人物}{ふきだしの文}
%% --------------------------------------------------------------------
\newcommand\kgaArtPage[7]{%
  \clearpage
  \begingroup\global\bk@decofalse\bkfolio{false}%
  \AddToHookNext{shipout/background}{\put(0,-\paperheight){%
  \begin{tikzpicture}[overlay, x=1mm, y=1mm]
    \fill[bkVol] (-3,104) rectangle (185,260);
    \begin{scope}\clip (-3,104) rectangle (185,260);
      \node[text=white, opacity=0.12, font=\sffamily\bfseries\fontsize{200}{200}\selectfont] at (140,190) {#1};
    \end{scope}
    \node[anchor=base west, text=white, font=\sffamily\bfseries\fontsize{16}{20}\selectfont] at (22,230) {#2};
    \node[anchor=north west, text=white, font=\sffamily\bfseries\fontsize{30}{40}\selectfont, text width=140mm, align=left, inner sep=0pt] at (21,222) {#3};
    \foreach \y/\t in {#4} {
      \node[anchor=west, inner sep=0pt] at (22,\y) {\bkicon[1.6em]{hatena}{bkYellow}};
      \node[anchor=west, text=white, font=\sffamily\bfseries\fontsize{12.5}{16}\selectfont] at (31,\y) {\t};
    }
    #5
    \bkChar{#6}{(28,62)}{30mm}
    \node[rectangle callout, callout absolute pointer={(43,80)}, anchor=west, draw=bkVol, line width=0.8pt, fill=white,
      rounded corners=3mm, inner xsep=4mm, inner ysep=3mm,
      text width=94mm, align=left, font=\fontsize{10.5}{18}\selectfont, text=bkInk] at (49,80) {#7};
  \end{tikzpicture}}}%
  \null\clearpage\endgroup\global\bk@decotrue\bkfolio{true}}

%% --------------------------------------------------------------------
%% 子ども用の本文の節の見出し（偶数ページから）  \kgasec{番号}{題名}{ひとこと}{しおり用の文字}
%%   ラベル kga-sec番号 を置く（もくじが拾う）。柱：偶数ページ＝\kgaBandLabel、奇数ページ＝番号と題名
%% --------------------------------------------------------------------
\newcommand\kgasec[4]{%
  \bkevenpage
  \global\bk@decotrue
  \bksetrunning{}{\kgaBandLabel}{#1\hspace{0.6\zw}#2}%
  \pdfbookmark[0]{#1　#4}{kgasec#1}\typeout{BKPAGE kga #1 \thepage}%
  \label{kga-sec#1}%
  \noindent\begin{tikzpicture}
    \useasboundingbox (0,0) rectangle (\linewidth,3.1);
    \begin{scope}\clip[rounded corners=2.5mm] (0,0) rectangle (\linewidth,3.1);
      \fill[bkTheme] (0,0) rectangle (\linewidth,3.1);
      \fill[white, opacity=0.13] (\linewidth-2.2,3.1) circle (2.0);
      \fill[white, opacity=0.10] (\linewidth-0.6,0) circle (1.3);
    \end{scope}
    \node[anchor=west, text=bkOn, font=\sffamily\small] at (0.5,2.6) {\kgaBandLabel};
    \node[anchor=west, circle, fill=white, text=\bkTT, minimum size=1.5cm, font=\sffamily\bfseries\Huge, inner sep=0pt] (n) at (0.45,1.25) {#1};
    \node[anchor=west, text=bkOn, font=\sffamily\bfseries\LARGE, text width=\linewidth-3.6cm, align=left, inner sep=0pt] at ([xshift=0.35cm]n.east) {#2};
  \end{tikzpicture}\par\vspace{2.5mm}
  \noindent{\sffamily\fontsize{10.5}{19}\selectfont\color{bkInk}#3}\par
  \vspace{1.5mm}{\color{bkLine}\hrule height 0.5pt}\vspace{3mm}}

%% --------------------------------------------------------------------
%% 奥付  \kgaColophon{書名の行}{副題の行}{版の行}{いちばん下の注}
%% --------------------------------------------------------------------
\newcommand\kgaColophon[4]{%
  \bkpadfour
  \global\bk@decofalse\bkfolio{false}%
  \null\vfill
  \noindent\begin{minipage}{\linewidth}\small
  {\sffamily\bfseries\Large\color{bkVol0}#1}\par\vspace{1mm}
  {\sffamily\bfseries\large #2}\par\vspace{3mm}
  {\color{bkVol0}\hrule height 1pt}\vspace{3mm}
  \begin{tabular}{@{}l@{\hspace{6mm}}l@{}}
  \R{編者}{へん|しゃ} & 「\bkSeries」\R{編集委員会}{へん|しゅう|い|いん|かい} \\
   & \R{委員長}{い|いん|ちょう}　\R{小学校長}{しょう|がっ|こう|ちょう}（まとめ\R{役}{やく}） \\
   & \R{委員}{い|いん}　\R{数学}{すう|がく}・\R{電気工学}{でん|き|こう|がく}・\R{機械工学}{き|かい|こう|がく}・\R{教育学}{きょう|いく|がく}・\R{量子力学}{りょう|し|りき|がく}・\R{建築工学}{けん|ちく|こう|がく}・\\
   & \phantom{委員　}\R{半導体工学}{はん|どう|たい|こう|がく}の\R{大学}{だい|がく}\R{教授}{きょう|じゅ}、\R{幼稚園}{よう|ち|えん}\R{教諭}{きょう|ゆ}、\R{塾}{じゅく}\R{講師}{こう|し}、\\
   & \phantom{委員　}\R{特別支援学校}{とく|べつ|し|えん|がっ|こう}\R{教諭}{きょう|ゆ}、\R{心理士}{しん|り|し}、\R{科学}{か|がく}\R{解説者}{かい|せつ|しゃ}、\R{小学校}{しょう|がっ|こう}\R{教諭}{きょう|ゆ} \\
  デザイン & ユニバーサルデザイン\R{担当}{たん|とう}\R{委員}{い|いん} \\
  \R{挿絵}{さし|え} & イラスト\R{担当}{たん|とう}\R{委員}{い|いん} \\
  \R{版}{はん} & #3 \\
  \end{tabular}\par\vspace{3mm}
  {\color{bkVol0}\hrule height 1pt}\vspace{2mm}
  {\footnotesize #4}
  \end{minipage}
  \clearpage}
\makeatother


\makeatletter
%% 表2：全9巻の構成（書名・副題・学年の目安・ページ・★の範囲）
% ★の書き方（付録Bと同じ約束）：塗りの★＝下の高さ、白抜きの☆＝上の高さまで  \kgastars{下}{上}
\newcommand\kgastars[2]{\mbox{\sffamily\@tempcnta=0\relax
  \loop\advance\@tempcnta\@ne
    \ifnum\@tempcnta>#1\relax\textcolor{bkStar!70!black}{☆}\else\textcolor{bkStar}{★}\fi
  \ifnum\@tempcnta<#2\repeat}}
\newcommand\kgaAdultInsideArt{%
  \fill[bkVol!7] (-3,-3) rectangle (185,260);
  \node[anchor=base west, text=bkDeep, font=\sffamily\bfseries\fontsize{20}{26}\selectfont] at (16,228) {全9巻の構成};
  \draw[bkVol, line width=1pt] (16,222) -- (166,222);
  \node[anchor=north west, text=bkInk, font=\sffamily\fontsize{9.5}{14}\selectfont, text width=150mm, align=left, inner sep=0pt] at (16,218)
    {学年の目安は編集上の目安です（各巻の紙面と子ども用のガイドブックには書いていません）。詳しくは第4章・第9章をご覧ください。};
  \foreach \t/\x in {巻/16, 書名・副題/30, 学年の目安/84, ページ/131, ★/146} {
    \node[anchor=base west, text=bkInk, font=\sffamily\bfseries\fontsize{9.5}{11}\selectfont] at (\x,200) {\t};}
  \draw[bkLine, line width=0.6pt] (16,197) -- (166,197);
  \foreach \v/\p in {1/400,2/456,3/564,4/560,5/492,6/600,7/492,8/632,9/640} {
    \pgfmathsetmacro\kga@ay{188-(\v-1)*17.2}
    \fill[bkVol\v, rounded corners=1mm] (16,\kga@ay-4.5) rectangle ++(10,10);
    \node[text=bkBandOn\v, font=\sffamily\bfseries\fontsize{12}{12}\selectfont] at (21,\kga@ay+0.5) {\v};
    \node[anchor=base west, inner sep=0pt, text=bkInk, font=\sffamily\bfseries\fontsize{11}{13}\selectfont] at (30,\kga@ay+1.4) {\bkTitle{\v}};
    \node[anchor=base west, inner sep=0pt, text=bkInk, font=\sffamily\fontsize{9.5}{11}\selectfont] at (30,\kga@ay-4) {\setbox0\hbox{\bkSub{\v}}\ifdim\wd0>0pt \bkSub{\v}\else ―\fi};
    \node[anchor=base west, inner sep=0pt, text=bkInk, font=\sffamily\fontsize{9.5}{11}\selectfont] at (84,\kga@ay-1) {\bkGrade{\v}};
    \node[anchor=base east, inner sep=0pt, text=bkInk, font=\sffamily\fontsize{9.5}{11}\selectfont] at (142,\kga@ay-1) {\p};
    \node[anchor=base west, inner sep=0pt, font=\fontsize{9.5}{11}\selectfont] at (146,\kga@ay-1) {\kgastars{\bkSmin{\v}}{\bkSmax{\v}}};
    \draw[bkLine, line width=0.4pt] (16,\kga@ay-7.5) -- (166,\kga@ay-7.5);}
  \node[anchor=base west, text=bkInk, font=\sffamily\fontsize{9.5}{11}\selectfont] at (16,25) {計　152章（序章9・章139・読みもの3・終章1）、4836ページ};
  \node[anchor=north west, text=bkInk, font=\sffamily\fontsize{9.5}{14}\selectfont, text width=150mm, align=left, inner sep=0pt] at (16,20)
    {★は章の内容の高さの目安です。塗りの★はその巻の下の高さ、白抜きの☆は上の高さまでを表します（例：\kgastars{2}{3}＝★★〜★★★）。★は指導要領の学年と一対一ではありません。};
}
\makeatother

\begin{document}
\makeatletter
%% ===================== 前付け =====================
\bkartpage{表1}{\kgaFrontArt{adult}}
\bkartpage{表2}{\kgaAdultInsideArt}
\kgaTitlePage{adult}
\kgaTocPages{%
  \noindent{\sffamily\fontsize{10}{17}\selectfont\color{bkGray}本書は、シリーズ『物理を理解するための数学』（全9巻）を使う保護者・先生のための案内です。第1章でシリーズの概要と使い方を、第2章以降で編集方針と、現行の学習指導要領（小学校・中学校は平成29年告示、高等学校は平成30年告示）との対応を説明します。}\par\vspace{4mm}
  \kgaTocRow{bkVol0}{1}{シリーズの概要と使い方}{目的、全9巻の構成、7本の路線、章の構成、★と3段、どこから始めるか}{kgd-ch1}
  \kgaTocRow{bkVol0}{2}{編集方針と書き方の原則}{編集方針10項目と、子どもに伝わる書き方の原則}{kgb-ch2}
  \kgaTocRow{bkVol0}{3}{学習指導要領との関係のあらまし}{現行の学習指導要領と、このシリーズの立場}{kgb-ch3}
  \kgaTocRow{bkVol0}{4}{学年と巻の対応}{小学1年から大学までの学年と、9巻の対応}{kgb-ch4}
  \kgaTocRow{bkVol0}{5}{巻ごとの対応}{各巻が指導要領のどこに当たるか、範囲外の内容は何か}{kgb-ch5}
  \kgaTocRow{bkVol0}{6}{追加した内容と、その理由}{微分方程式・行列・ベクトル解析・フーリエ解析・ラプラス変換など}{kgb-ch6}
  \kgaTocRow{bkVol0}{7}{含めていない・軽く扱う内容と、補い方}{学校の教科書・問題集で補う箇所と、その時期}{kgb-ch7}
}{%
  \kgaTocRow{bkVol0}{8}{大学の工学部の数学との対応}{第6〜9巻と、工学部1〜2年の数学の科目}{kgb-ch8}
  \kgaTocRow{bkVol0}{9}{★と学年の目安の読み方}{★・学年の目安・指導要領の学年は、それぞれ別のものです}{kgb-ch9}
  \kgaTocRow{bkVol0}{10}{授業・家庭での使い方の例}{今後刊行する問題集・指導書についても}{kgb-ch10}
  \kgaTocRow{bkVol0}{A}{付録A　学習指導要領の項目からシリーズの章へ}{小学1年〜高校の各科目の項目が、どの章にあるか}{kgc-a}
  \kgaTocRow{bkVol0}{B}{付録B　全152章の一覧}{章名・★・路線・相当する学年・範囲外の内容}{kgc-b}
  \kgaTocRow{bkVol0}{＊}{出典}{学習指導要領と解説、参考にした資料}{kgb-src}
  \vfill
  \noindent\begin{tikzpicture}
    \node[anchor=north west, fill=bkVol0!8, rounded corners=3mm, inner sep=4mm, text width=\linewidth-8mm,
      align=left, font=\fontsize{9.5}{17}\selectfont, text=bkInk] at (0,0)
      {{\sffamily\bfseries\color{bkVol0}本書の使い方}\par
       シリーズの全体像を知りたいときは第1章から、お子さんの学年からシリーズの巻を探すときは第4章と付録A、シリーズの章から学校の学習内容を確かめるときは第5章と付録Bをご覧ください。お子さん向けには、ルビつきの子ども用『ガイドブック　はじめに読む本』があります。};
  \end{tikzpicture}
}
\makeatother

%% ===================== 本文 =====================
\bkevenpage
\IfFileExists{guide-a1.tex}{%% guide-a1.tex　ガイドブック　保護者・先生のための本（大人用、vol=0）第1章「シリーズの概要と使い方」
%% （第16回編集委員会の追加の決定、組 G-B、頭文字 kgd）。ルビなし。大人向けの漢字の使い方（10月30日 10:15 依頼者）。
%% 材料：tools/guide-series.md（RG-2）。数・事実は材料のとおり。
%% このファイルの先頭で、大人用の本文（guide-a1・guide-p2a・guide-p2a-src）と付録（guide-p2b、G-C）が共通に使う部品を定義する。
%% 親ファイル（tex/guide-adult.tex、G-A）は guide-a1 → guide-p2a → guide-p2b → guide-p2a-src の順に読む。
\protected\def\R#1#2{#1}% 大人用の本はルビなし（親ファイルでも同じ定義をする）
% 文中の「→」のあとで改行しない（行末に→を残さない。大人用の本の全体に効く）
\ltjsetparameter{postbreakpenalty={`→,10000}}
% ★の並び（★★〜★★★ など）が行末で切れないようにする（委員確認 rA #22。付録にも効く）
\ltjsetparameter{postbreakpenalty={`★,10000}}\ltjsetparameter{postbreakpenalty={`☆,10000}}
\ltjsetparameter{prebreakpenalty={`〜,10000}}\ltjsetparameter{postbreakpenalty={`〜,10000}}
\makeatletter
%% ===== 章見出し・小見出し（guide-p2a・guide-p2b・guide-p2a-src と共通。G-C は \providecommand の写しを持つ） =====
% \kgbsec[柱・しおり用の題（省略可）]{番号}{題}。改ページはしない（前の章に続けて組む。見出しだけがページの下に残らないよう 8 行分あける）
\providecommand\kgbsec[3][]{%
  \par\addvspace{18pt}\Needspace{8\baselineskip}\global\bk@decotrue
  \setcounter{bkfig}{0}\setcounter{bktab}{0}\gdef\kgbsecno{#2}%
  \if\relax\detokenize{#1}\relax\def\kgb@st{#3}\else\def\kgb@st{#1}\fi
  \bksetrunning{}{保護者・先生のための本}{\if\relax\detokenize{#2}\relax\else#2\hspace{0.6\zw}\fi\kgb@st}%
  \pdfbookmark[1]{#2　\kgb@st}{kgbsec-#2}%
  \kgb@secband{#2}{#3}}
\newcommand\kgb@secband[2]{%
  \noindent\begin{tikzpicture}
    \if\relax\detokenize{#1}\relax
      \node[anchor=west, text=white, font=\sffamily\bfseries\LARGE, inner sep=0pt] (t) at (0.6,0) {#2};
    \else
      \node[anchor=west, circle, fill=white, text=bkVol0, minimum size=1.35cm, font=\sffamily\bfseries\Large, inner sep=0pt] (n) at (0.35,0) {#1};
      \node[anchor=west, text=white, font=\sffamily\bfseries\LARGE, text width=\linewidth-2.4cm, align=left, inner sep=0pt] (t) at ([xshift=0.4cm]n.east) {#2};
    \fi
    \begin{pgfonlayer}{background}
      \path let \p1=(t.north), \p2=(t.south) in coordinate (kgbT) at (0,{max(\y1+0.4cm,1cm)}) coordinate (kgbB) at (0,{min(\y2-0.4cm,-1cm)});
      \fill[bkVol0] (kgbB) rectangle (\linewidth,0 |- kgbT);
    \end{pgfonlayer}
  \end{tikzpicture}\par\vspace{10pt}}
\providecommand\kgbsub[2]{%
  \par\addvspace{12pt}\Needspace{5\baselineskip}\noindent
  \begin{tikzpicture}[baseline=(t.base)]
    \node[anchor=west, fill=bkVol0, text=white, rounded corners=1mm, font=\sffamily\bfseries\small, inner xsep=1.6mm, inner ysep=1mm] (i) at (0,0) {#1};
    \node[anchor=west, font=\sffamily\bfseries\large, text=bkVol0, inner sep=0pt] (t) at ([xshift=0.6\zw]i.east) {#2};
    \draw[bkVol0, line width=0.8pt] ([yshift=-4pt]i.south west) -- ++(\linewidth-1pt,0);
  \end{tikzpicture}\par\nobreak\vspace{6pt}}
% 図・表の番号は章ごとに「章-番号」（図4-1、表6-2 など）
\providecommand\kgbsecno{}
\renewcommand\thebkfig{\kgbsecno-\arabic{bkfig}}\renewcommand\thebktab{\kgbsecno-\arabic{bktab}}
\renewcommand\theHbkfig{kgb.\kgbsecno.\arabic{bkfig}}\renewcommand\theHbktab{kgb.\kgbsecno.\arabic{bktab}}% リンクの名前がぶつからないように
%% ===== 大人用の本文（第1〜10章）で使う部品（頭文字 kgb。名前は第2部のときのまま） =====
% 巻の札  \kgbvol{1〜9}（巻の色の地に、その巻の帯の字の色）
\newcommand\kgbvol[1]{\tikz[baseline=(v.base)]\node[fill=bkVol#1, text=bkOn#1, rounded corners=0.8mm, inner xsep=1.2mm, inner ysep=0.7mm, font=\sffamily\bfseries\footnotesize] (v) {第#1巻};}
% 箇条書き（1字下げのぶら下がり）  \kgbitem{文}
\newcommand\kgbitem[1]{\par\noindent\hangindent=1\zw\hangafter=1 ・#1\par}
% 注（表・図の下の小さな字の説明）  \kgbnote{文}
\newcommand\kgbnote[1]{\par\vspace{1pt}{\footnotesize\noindent\hangindent=1\zw\hangafter=1 #1\par}\vspace{2pt}}
% ★の範囲（表の中で）  \kgbst{最小}{最大}
\newcommand\kgbst[2]{\mbox{\footnotesize\bkstars{#1}\ifnum#1<#2\relax{\color{bkInk}〜}\bkstars{#2}\fi}}
% 長い表（ページをまたげる表。続きのページでも1行めの見出しをくり返す）
% 表の名前は、表の最後の行（線なし）に置き、すぐ上の行とページで離れないようにする（\\*）。
% 使い方：\kgbcap{表の名前}{ラベル} を \begin{bktable}[long, theme=kgb] の直前に置く
\NewTblrTheme{kgb}{\DefTblrTemplate{caption}{default}{}\DefTblrTemplate{conthead}{default}{}%
  \DefTblrTemplate{capcont}{default}{}\DefTblrTemplate{contfoot}{default}{}%
  \DefTblrTemplate{note}{default}{}\DefTblrTemplate{remark}{default}{}}
\def\kgb@capt{}
\newcommand\kgbcap[2]{\refstepcounter{bktab}\if\relax\detokenize{#2}\relax\else\label{#2}\fi\gdef\kgb@capt{#1}}
\newcommand\kgbcapline{{\sffamily\small{\color{\bkTT}\bfseries 表\thebktab}\hspace{0.8\zw}\kgb@capt}}
% 学年と巻の対応の帯の図：左の巻の欄 24 mm、小1〜小6・中1〜中3・高校の6科目は 7.2 mm、いちばん右の「大学」の列は 8 mm
\newcommand\kgbgridcell[3]{% {y}{巻}{列} 内容がある学年・科目のます
  \pgfmathsetmacro\kgb@x{ifthenelse(#3<16,24+(#3-1)*7.2,132)}\pgfmathsetmacro\kgb@w{ifthenelse(#3<16,7.2,8)}%
  \fill[bkVol#2, draw=bkInk!55, line width=0.4pt] ({\kgb@x+(\kgb@w-5.4)/2},{#1-2.7}) rectangle ++(5.4,5.4);}
\newcommand\kgbgridband[4]{% {y}{巻}{はじめの列}{終わりの列} 学年の目安の太い枠
  \pgfmathsetmacro\kgb@xa{ifthenelse(#3<16,24+(#3-1)*7.2,132)}\pgfmathsetmacro\kgb@xb{ifthenelse(#4<16,24+#4*7.2,140)}%
  \draw[bkInk, line width=1.3pt, rounded corners=1.3mm] ({\kgb@xa+0.35},{#1-3.6}) rectangle ({\kgb@xb-0.35},{#1+3.6});}
\newcommand\kgbgridempty[1]{% {y} 空のますのうすい枠（16列）
  \foreach \kgb@c in {1,...,16} {\pgfmathsetmacro\kgb@x{ifthenelse(\kgb@c<16,24+(\kgb@c-1)*7.2,132)}\pgfmathsetmacro\kgb@w{ifthenelse(\kgb@c<16,7.2,8)}%
    \draw[bkLine, line width=0.4pt] ({\kgb@x+(\kgb@w-5.4)/2},{#1-2.7}) rectangle ++(5.4,5.4);}}
% 1行  \kgbgridrow{y}{巻}{内容がある列のならび}{目安のはじめの列}{目安の終わりの列}
\newcommand\kgbgridrow[5]{%
  \fill[bkPaper] (132,{#1-4.1}) rectangle (140,{#1+4.1});
  \kgbgridempty{#1}%
  \foreach \kgb@c in {#3} {\kgbgridcell{#1}{#2}{\kgb@c}}%
  \kgbgridband{#1}{#2}{#4}{#5}%
  \node[anchor=west, inner sep=0pt] at (0.5,#1) {\kgbvol{#2}};}
% 見出し（上から y=0 まで）。#1＝高校の上に学年の例の行をつけるか（1/0）
\newcommand\kgbgridhead[1]{%
  \tikzset{kgb gh/.style={font=\sffamily\bfseries\footnotesize, text=bkInk, inner sep=0.5pt}}
  \ifnum#1=1 \def\kgb@hy{12}\else\def\kgb@hy{4.8}\fi
  \fill[bkPaper] (132,0) rectangle (140,\kgb@hy+7);
  \node[kgb gh] at ({24+3*7.2},\kgb@hy+3.5) {小学校（学年）};
  \node[kgb gh] at ({24+7.5*7.2},\kgb@hy+3.5) {中学校};
  \node[kgb gh] at ({24+12*7.2},\kgb@hy+3.5) {高校（数学の科目）};
  \node[kgb gh] at (136,\kgb@hy+3.5) {大学};
  \foreach \kgb@c/\kgb@t in {1/1,2/2,3/3,4/4,5/5,6/6,7/1,8/2,9/3,10/Ⅰ,11/A,12/Ⅱ,13/B,14/Ⅲ,15/C}
    {\node[kgb gh] at ({24+(\kgb@c-0.5)*7.2},3) {\kgb@t};}
  \ifnum#1=1
    \foreach \kgb@c/\kgb@t in {10/1年の例,12/2年の例,14/3年の例}
      {\node[kgb gh, font=\sffamily\footnotesize] at ({24+\kgb@c*7.2},8.6) {\kgb@t};
       \draw[bkGray, line width=0.5pt] ({24+(\kgb@c-1)*7.2+0.8},6.6) -- ({24+(\kgb@c+1)*7.2-0.8},6.6);}
  \fi
  \foreach \kgb@x in {24,{24+6*7.2},{24+9*7.2},132,140} {\draw[bkGray, line width=0.6pt] (\kgb@x,\kgb@hy+7) -- (\kgb@x,0);}
  \draw[bkGray, line width=0.6pt] (24,0) -- (140,0);}
% 縦の区切り線を下まで  \kgbgridlines{いちばん下の y}
\newcommand\kgbgridlines[1]{\foreach \kgb@x in {24,{24+6*7.2},{24+9*7.2},132,140} {\draw[bkGray, line width=0.6pt] (\kgb@x,0) -- (\kgb@x,#1);}}
% 各巻のカード（5章）の1行の帯  \kgbstrip{巻}{内容がある列}{目安のはじめ}{目安の終わり}
\newcommand\kgbstrip[4]{\begin{tikzpicture}[x=0.935mm, y=1mm]
  \kgbgridhead{0}\kgbgridrow{-5}{#1}{#2}{#3}{#4}\kgbgridlines{-9.1}\end{tikzpicture}}
% 巻のカード（5章）  \begin{kgbcard}{巻}{題}
\newtcolorbox{kgbcard}[2]{enhanced, breakable=false, arc=2mm, boxrule=1pt, colframe=bkVol#1, colback=white,
  left=3mm, right=3mm, top=1.2mm, bottom=1.6mm, before skip=6pt, after skip=6pt,
  title={\kgbvol{#1}\hspace{0.6\zw}{\color{bkInk}#2}}, colbacktitle=bkVol#1!14, coltitle=bkInk,
  fonttitle=\sffamily\bfseries, toptitle=0.6mm, bottomtitle=0.6mm, titlerule=0pt}
% カードの中の1行（見出し＋文）  \kgbrow{見出し}{文}
\newcommand\kgbrow[2]{\par\addvspace{1pt}{\footnotesize\baselineskip=15.2pt\noindent\hangindent=6.6\zw\hangafter=1
  \makebox[6.6\zw][l]{\sffamily\bfseries\color{bkVol0}#1}#2\par}}
% 図の箱
\tikzset{
  kgb box/.style={draw=bkVol0, line width=0.9pt, fill=white, rounded corners=1.5mm, align=center,
    font=\sffamily\footnotesize, text=bkInk, inner xsep=2mm, inner ysep=1.6mm},
  kgb gbox/.style={draw=bkGray, line width=0.8pt, fill=white, rounded corners=1.5mm, align=left,
    font=\footnotesize, text=bkInk, inner xsep=1.8mm, inner ysep=1.4mm},
  kgb tbox/.style={draw=bkVol0, line width=0.9pt, fill=bkVol0!8, rounded corners=1.5mm, align=left,
    font=\footnotesize, text=bkInk, inner xsep=1.8mm, inner ysep=1.4mm},
  kgb arr/.style={-{Stealth[length=2.4mm]}, line width=1pt, bkInk!75, shorten <=1.2mm, shorten >=1.2mm},
  kgb lab/.style={font=\sffamily\footnotesize, text=bkInk, inner sep=0.8pt},
}
\makeatother

%% ===== 第1章だけで使う部品（頭文字 kgd） =====
\makeatletter
% 本棚の1冊（背）  \kgdspine{x}{巻}{はば}{題}{ページ}
\newcommand\kgdspine[5]{%
  \fill[bkVol#2, draw=bkInk!60, line width=0.5pt, rounded corners=0.6mm] (#1,0) rectangle ++(#3,64);
  \node[text=bkOn#2, font=\sffamily\bfseries\footnotesize, inner sep=0pt] at ({#1+#3/2},58.5) {#2};
  \node[text=bkOn#2, font=\sffamily\bfseries\footnotesize, inner sep=0pt, anchor=north, align=center, text width=#3 mm, execute at begin node={\baselineskip=10.6pt}] at ({#1+#3/2},53.5) {#4};% 題は1字ずつ縦に積む（\\ で区切る）
  \node[font=\sffamily\footnotesize, text=bkInk, inner sep=0pt, anchor=north] at ({#1+#3/2},-1.6) {#5};}
% 章の流れの1つの札  \kgdpart{名前}{アイコン}{x}{y}{ノード名}
\newcommand\kgdpart[5]{\node[kgb box, anchor=west, minimum width=43.5mm, text width=39.5mm, minimum height=8.6mm, align=left, inner xsep=1.5mm, inner ysep=1mm] (#5) at (#3,#4)
  {\makebox[5mm][l]{\if\relax\detokenize{#2}\relax\else\ifcsname kgd@ibox@#2\endcsname\usebox{\csname kgd@ibox@#2\endcsname}\else\bkicon[1.1em]{#2}{bkVol0}\fi\fi}\hspace{0.5mm}#1};}
% 中に \node を持つアイコン（はてな）は、図の外で箱に入れておく（図の中で入れ子にすると位置がずれるため。委員確認 rA #20）
\newsavebox\kgd@ibox@hatena
\sbox\kgd@ibox@hatena{\bkicon[1.1em]{hatena}{bkVol0}}
\makeatother

%% ====================================================================
\kgbsec{1}{シリーズの概要と使い方}\label{kgd-ch1}

この本は、『物理を理解するための数学』全9巻を使う保護者・先生のための案内です。第1章でシリーズの概要と使い方を、第2章以降で編集方針と学習指導要領との対応を説明します。子ども向けの『はじめに読む本』を読まなくても、この1冊でシリーズの全体と使い方が分かるように書きました。

\kgbsub{1.1}{このシリーズの目的}

このシリーズは、工学の研究を発展させるために、小学校から高校までの算数・数学の学び方を見直す、という考えのもとで作られました。小学1年の「リンゴは何個？」から始めて、工学部で使う行列・微分方程式・ベクトル解析・フーリエ解析・ラプラス変換までを、一続きの9巻で学べるようにしています。

狙いは、入試を解くための発想力の数学ではなく、物理や工学の現象を説明するための数学を、子どもの「なぜ？」「何がうれしいの？」に答えながら教えることです。このシリーズを終えてから大学の専門書を開くと、書いてあることがすんなり頭に入る状態を目標にしています。

\begin{bkfigure}
\begin{tikzpicture}[x=1mm, y=1mm]
  \node[kgb box, text width=30mm] (a) at (17,0) {第1巻 1-1\\[1pt]{\bfseries いくつ あるかな？}\\（数える・比べる）};
  \node[kgb box, text width=36mm, minimum height=19mm] (b) at (70,0) {7本の路線\\[1pt]{\mdseries 変化・たし合わせ・倍々・\\回転と波・表・矢印・はかる}};
  \node[kgb box, text width=32mm, fill=bkVol0!8] (c) at (120.5,0) {工学部の専門書\\[1pt]{\mdseries 行列・微分方程式・\\ベクトル解析・\\フーリエ解析・\\ラプラス変換}};
  \draw[kgb arr] (a) -- (b);
  \draw[kgb arr] (b) -- (c);
  \node[kgb lab, anchor=north] at (70,-21) {全9巻・152章（4836ページ）を、7本の路線でつなぐ};
\end{tikzpicture}
\bkfigcaption{このシリーズの出発点と到達点}
\end{bkfigure}

本文の表記は、読む子どもの学年に合わせています。第1〜3巻は全ての漢字にルビを付け（第1巻は分かち書きも付けます）、第4巻からはルビを付けずに常用漢字までを使います。

\kgbsub{1.2}{9巻の並び}

全9巻は、数えることから始めて、関数とデータ、図形とベクトル、微分と積分を経て、大学の工学の数学へ進みます。

\begin{bkfigure}
\begin{tikzpicture}[x=1mm, y=1mm]
  \draw[bkGray, line width=1.2pt] (0.6,-0.3) -- (139.4,-0.3);
  \kgdspine{0.10}{1}{11.11}{か\\ぞ\\え\\る\\・\\は\\か\\る}{400}
  \kgdspine{11.89}{2}{12.67}{数\\と\\量\\の\\こ\\と\\ば}{456}
  \kgdspine{25.24}{3}{15.67}{関\\数\\と\\デ\\\rotatebox[origin=c]{-90}{ー}\\タ}{564}
  \kgdspine{41.59}{4}{15.56}{形\\と\\向\\き}{560}
  \kgdspine{57.83}{5}{13.67}{変\\化\\を\\つ\\か\\む}{492}
  \kgdspine{72.18}{6}{16.67}{未\\来\\を\\計\\算\\す\\る}{600}
  \kgdspine{89.53}{7}{13.67}{ま\\と\\め\\て\\計\\算\\す\\る}{492}
  \kgdspine{103.88}{8}{17.56}{空\\間\\に\\広\\が\\る\\量}{632}
  \kgdspine{122.12}{9}{17.78}{波\\と\\変\\化\\を\\読\\み\\と\\く}{640}
  \node[kgb lab, anchor=north west, text=bkGray] at (0,-6) {本の下の数はページ数。背の幅はページ数に比例させた（1 mm＝36ページ）};
\end{tikzpicture}
\bkfigcaption{全9巻の本棚（色は各巻の表紙の色）}
\end{bkfigure}

\par\addvspace{4pt}
\kgbcap{9巻の一覧（学年の目安・★の範囲・章の数・ページ数）}{kgd-t-vols}
\begin{bktable}[long, theme=kgb]{hline{Z}={0pt,white}, vline{1-Z}={Z}{0pt,white}, row{Z}={abovesep=3pt, belowsep=1pt}, hspan=minimal, rowhead=1, colspec={Q[c,m,wd=13mm] Q[l,m,wd=35mm] Q[l,m,wd=26mm] Q[c,m,wd=31mm] Q[c,m,wd=8mm] Q[r,m,wd=11mm]}, rowsep=1.5pt, colsep=3pt, cells={font=\small}, row{1}={font=\sffamily\bfseries\small}}
巻 & 書名（副題） & 学年の目安 & ★の範囲 & 章 & ページ \\
\kgbvol{1} & かぞえる・はかる & 小学1年〜3年 & \kgbst{1}{1} & 22 & 400 \\
\kgbvol{2} & 数と量のことば & 小学4年〜中学1年 & \kgbst{1}{2} & 20 & 456 \\
\kgbvol{3} & 関数とデータ（きまりを見つける） & 中学1年〜3年 & \kgbst{1}{3} & 18 & 564 \\
\kgbvol{4} & 形と向き（図形・三角関数・ベクトル） & 中学2年〜高校2年 & \kgbst{1}{5} & 16 & 560 \\
\kgbvol{5} & 変化をつかむ（微分と積分） & 高校1年〜2年 & \kgbst{2}{4} & 19 & 492 \\
\kgbvol{6} & 未来を計算する（微分方程式と複素数） & 高校2年〜3年 & \kgbst{2}{5} & 16 & 600 \\
\kgbvol{7} & まとめて計算する（行列） & 高校3年（発展）〜大学1年 & \kgbst{3}{5} & 12 & 492 \\
\kgbvol{8} & 空間に広がる量（ベクトル解析） & 大学1年〜2年 & \kgbst{3}{5} & 15 & 632 \\
\kgbvol{9} & 波と変化を読みとく（フーリエ解析とラプラス変換） & 大学2年 & \kgbst{3}{5} & 14 & 640 \\
計 & & & & 152 & 4836 \\*
\SetCell[c=6]{c} \kgbcapline & & & & & \\
\end{bktable}
\par\addvspace{4pt}

\kgbnote{※ 章の数は序章を含みます。第7・8巻は読みもの1章、第9巻は読みもの1章と終章を含みます（全152章＝序章9・章139・読みもの3・終章1）。ページ数は奥付までのページです。}
\kgbnote{※ 学年の目安は、編集のときに決めた目安で、生徒用の紙面には書いていません（理由と読み方は 4.2・9.3）。実際には、目安の学年より前や後に学校で習う内容も含まれます（図\ref{kgb-f-grid}）。}

各巻は、いくつかの「部」に分かれています。

\par\addvspace{4pt}
\kgbcap{各巻の部の構成}{}
\begin{bktable}[long, theme=kgb]{hline{Z}={0pt,white}, vline{1-Z}={Z}{0pt,white}, row{Z}={abovesep=3pt, belowsep=1pt}, hspan=minimal, rowhead=1, colspec={Q[c,m,wd=13mm] Q[l,m,wd=113mm]}, rowsep=1.2pt, colsep=4pt, cells={font=\small}, row{1}={font=\sffamily\bfseries\small}}
巻 & 部の構成（序章の後の部の名前） \\
\kgbvol{1} & かずの はじまり／はかる／かけ算と わり算／表・グラフ・かたち \\
\kgbvol{2} & 量と数／大きな数・小さな数・単位／文字と式／コンピュータで計算 \\
\kgbvol{3} & 関数とグラフ／曲がったグラフ／倍々とけた／データと統計 \\
\kgbvol{4} & 図形／三角比と三角関数／ベクトル \\
\kgbvol{5} & 微分／積分／特別な関数 \\
\kgbvol{6} & 微分方程式の入口／複素数と双曲線関数／式で解く \\
\kgbvol{7} & 行列の計算／回転と座標の変換／固有値と行列の使いみち \\
\kgbvol{8} & たくさんの変数の微分／ベクトル場と積分／発散と回転／偏微分方程式 \\
\kgbvol{9} & フーリエ解析／ラプラス変換／行列と微分方程式／終章 \\*
\SetCell[c=2]{c} \kgbcapline & \\
\end{bktable}
\par\addvspace{4pt}

\kgbsub{1.3}{7本の路線}

全9巻の内容は、7本の「路線」でつながっています。どの章も、どれかの路線の駅になっています。①〜④の路線は、第6巻の微分方程式で乗り換えます。各巻の前付けには全9巻の路線図があり（図\ref{kgd-f-map}）、第9巻の終章（9-13）では、7本の路線がどこで合流したかをたどります。路線の色には、必ず形の記号と番号を添えています（第1巻の紙面では路線名を「たしあわせ」「やじるし」と平仮名で書きます）。

\par\addvspace{4pt}\Needspace{80mm}% 表1-3 を1ページに
\kgbcap{7本の路線（主な駅は巻-章の番号）}{kgd-t-route}
\begin{bktable}[long, theme=kgb]{hline{Z}={0pt,white}, vline{1-Z}={Z}{0pt,white}, row{Z}={abovesep=3pt, belowsep=1pt}, hspan=minimal, rowhead=1, colspec={Q[l,m,wd=27mm] Q[l,m,wd=31mm] Q[l,m,wd=49mm] Q[l,m,wd=22mm]}, rowsep=1.2pt, colsep=3pt, cells={font=\footnotesize}, row{1}={font=\sffamily\bfseries\footnotesize}}
路線 & 流れ & 主な駅（巻-章） & 終点 \\
\bkrouteicon[1.1em]{1}\ ① 変化 & ひき算→傾き→微分 & 1-5・2-3・3-5・5-1・6-1・8-4 & 流体・熱 \\
\bkrouteicon[1.1em]{2}\ ② たし合わせ & かけ算→面積→積分 & 1-21・2-3・4-5・5-10・6-3・8-11 & 電磁気 \\
\bkrouteicon[1.1em]{3}\ ③ 倍々 & 累乗→指数→$e$ & 2-7・3-11・5-14・6-4・9-5 & 制御 \\
\bkrouteicon[1.1em]{4}\ ④ 回転と波 & 角→sin→複素数 & 1-20・4-8・5-15・6-10・9-1 & 音・交流 \\
\bkrouteicon[1.1em]{5}\ ⑤ 表 & 九九の表→行列 & 1-13・2-18・3-6・7-1・9-10 & 回路網・制御 \\
\bkrouteicon[1.1em]{6}\ ⑥ 矢印 & 向き→ベクトル→場 & 2-2・3-2・4-10・7-7・8-5 & ロボット・電磁気 \\
\bkrouteicon[1.1em]{7}\ ⑦ はかる & 目もり→誤差→統計 & 1-11・2-6・3-14・9-4 & 計測・診断 \\*
\SetCell[c=4]{c} \kgbcapline & & & \\
\end{bktable}
\par\addvspace{4pt}


\kgbsub{1.4}{1つの章の構成と紙面の部品}

第2巻からは、どの章も同じ順番で進みます（第1巻も、言葉を易しくした同じ形です）。章の長さは章ごとに違います（序章を除くと、短い章で14ページ、長い章で54ページ）。章と部扉は、必ず偶数ページ（見開きの左）から始まります。

\begin{bkfigure}
\begin{tikzpicture}[x=1mm, y=1mm]
  \foreach \kgdn/\kgdi/\kgdx/\kgdy/\kgdk in {
    章の帯・道しるべ//0/0/p1, できるようになること/goal/0/-11/p2, はてな？/hatena/0/-22/p3, やってみよう/yattemiyou/0/-33/p4,
    思い出そう/omoidasu/0/-44/p5, 数学の道具/doguu/0/-55/p6,
    もっと知りたい・\\せいかくに言うと/motto/47.5/0/p7, たいせつ/taisetsu/47.5/-11/p8, ことば・記号カード/card/47.5/-22/p9,
    例題・じぶんでとこう/reidai/47.5/-33/p10, まちがえやすいところ/tsumazuki/47.5/-44/p11, 物理・工学では\\こう使う/butsuri/47.5/-55/p12,
    章末問題/jibun/95/0/p13, ふりかえり/check/95/-11/p14} {\kgdpart{\kgdn}{\kgdi}{\kgdx}{\kgdy}{\kgdk}}
  \foreach \kgda/\kgdb in {p1/p2,p2/p3,p3/p4,p4/p5,p5/p6,p7/p8,p8/p9,p9/p10,p10/p11,p11/p12,p13/p14} {\draw[kgb arr, shorten <=0.4mm, shorten >=0.4mm] (\kgda) -- (\kgdb);}
  \draw[kgb arr, rounded corners=2mm] (p6.east) -- ++(2.5,0) |- (p7.west);
  \draw[kgb arr, rounded corners=2mm] (p12.east) -- ++(2.5,0) |- (p13.west);
  \node[kgb gbox, anchor=north west, text width=39.5mm] at (95,-20) {{\sffamily\bfseries 専門書をのぞいてみよう}\\★★★★以上の章。専門書での書き方を1行ずつ訳す};
  \node[kgb gbox, anchor=north west, text width=39.5mm] at (95,-40) {{\sffamily\bfseries ふきだし}\\案内役が次の章へつなぐ（1見開きに1回まで）};
\end{tikzpicture}
\bkfigcaption{1つの章の流れ（左の列から順に。右下の2つは、章の途中に出てくる部品）}\label{kgd-f-flow}
\end{bkfigure}

\par\addvspace{4pt}
\kgbcap{紙面の部品と役割}{kgd-t-parts}
\begin{bktable}[long, theme=kgb]{hline{Z}={0pt,white}, vline{1-Z}={Z}{0pt,white}, row{Z}={abovesep=3pt, belowsep=1pt}, hspan=minimal, rowhead=1, colspec={Q[l,m,wd=36mm] Q[l,m,wd=92mm]}, rowsep=1.2pt, colsep=4pt, cells={font=\footnotesize}, row{1}={font=\sffamily\bfseries\footnotesize}}
部品 & 役割 \\
章の帯 & 章番号・題名・部の名前・★。題名は子どもに問いかける形を基本にし、目次にもそのまま載せる \\
道しるべ & 帯のすぐ下。「この章の前に」（2〜3章）と「この章の先に」 \\
できるようになること & その章の目標（普通は3つ）。各項目の頭にその★ \\
はてな？ & 章の問い（身近な不思議）。「数学の道具」の中で必ず答える \\
やってみよう & 実験・観察・電卓・表計算。手順に番号 \\
思い出そう & 前の巻・前の章の要点を2〜4行、章番号付きで。前の巻にあることは繰り返して教えない \\
数学の道具 & 小見出しごとに1つの考え \\
もっと知りたい & 発展。★付き \\
せいかくに言うと & たとえ話のすぐ後に、どこまで正しいかと正確な定義 \\
たいせつ & 1行に1つのこと。定義を置く \\
この章で出てきたことば & 新しい言葉と短い説明 \\
記号カード & 新しい記号・単位の読み方・意味・英語・初めて出た章。1章に2枚まで \\
例題・じぶんでとこう & 「考え方」「式」「答え」の札（第1〜3巻は「しき」「こたえ」）。考え方と途中の式を1行ずつ書く。「じぶんでとこう」は例題と同じ形で同じ数 \\
まちがえやすいところ & ×と○を1行ずつ、もう1行で「なぜ間違いか」 \\
物理・工学ではこう使う & 本物の機械・仕組みの例（第1巻は「つかわれている ところ」）。「大学への種」を添える章もある \\
コラム & 話題の広がり \\
章末問題 & 基本とチャレンジ。問題の形は全て、例題か「じぶんでとこう」で一度見せてある \\
ふりかえり & 「できるようになること」と同じ項目をチェックの形で。各項目に戻るページ \\
ふきだし & ソラ・ミナ・デルタが次の章へつなぐ。1見開きに1回まで \\
専門書をのぞいてみよう & ★★★★以上の章。専門書での書き方と、1行ずつの訳 \\
「ここまでで十分」の旗 & 3段の章の「手でわかる」の終わり（1.5） \\*
\SetCell[c=2]{c} \kgbcapline & \\
\end{bktable}
\par\addvspace{4pt}

\begin{figure}[p]
\noindent\begin{minipage}{\linewidth}\centering
\bkSeriesMap
\bkfigcaption{全9巻の路線図（縦が巻、横が路線。各巻の前付けにあるものと同じ図）}\label{kgd-f-map}
\end{minipage}
\end{figure}

1冊の本は、次の順に並んでいます。

\begin{bkfigure}
\begin{tikzpicture}[x=1mm, y=1mm]
  \foreach \kgdt/\kgdw/\kgdf/\kgdx in {表紙/12/bkVol0!25/0.5, 扉/9/white/13, 目次/11/white/23, 路線図/13/white/35, この本の\\使い方/16/white/49, 序章・部扉・\\章（本文）/31/bkVol0!12/66, 巻末/15/white/98, 奥付/11/white/114, 裏表紙/12.6/bkVol0!25/126} {
    \node[draw=bkVol0, line width=0.8pt, fill=\kgdf, minimum width=\kgdw mm, minimum height=13mm, anchor=west, align=center, font=\sffamily\footnotesize, text=bkInk, inner sep=0.5mm] at (\kgdx,0) {\kgdt};}
  \node[kgb lab, anchor=north west, text=bkGray, text width=138mm, align=left] at (0,-8) {扉は1ページ、目次は2〜3ページ（第1巻は2〜5ページ）、路線図は4〜5ページ（第1巻は6〜7ページ）、「この本の使い方」は6〜7ページ（第1巻は8〜9ページ）。巻末の中身は 1.9。};
\end{tikzpicture}
\bkfigcaption{1冊の並び（表紙の裏には案内役3人の紹介、裏表紙には全巻の一覧）}
\end{bkfigure}

\kgbsub{1.5}{★と3つの段、「名前」と「実際の操作」の型}

各章と「できるようになること」の各項目には、★を1〜5個付けています。★は、このシリーズの中での難しさ、つまり、その章でどの道具を使うかの目安です（★の数と内容の高さのおおよその対応は、9.1 の表）。範囲のある章は「★〜★★」のように書き、「この本の使い方」では★を「どこまで登るか」の山の図で説明しています。子どもには「自分の読める★まで読んで、先へ進んでもだいじょうぶ」と伝えています。

大学の内容を扱う第6〜9巻では、紙面に2つの型を使っています。1つめは章の3段です。大学の内容も中学生の言葉で書くため（編集方針8）、6-10〜6-14、7-7〜7-10、8-1〜8-13、9-1〜9-11 の各章を次の3段に分けています（第6〜9巻の★★★★★の章は、読みものと終章を除いてすべて含まれます）。

\begin{bkfigure}
\begin{tikzpicture}[x=1mm, y=1mm]
  \foreach \kgbi/\kgbic/\kgbn/\kgbd/\kgbs/\kgbf in {0/me/目でわかる/図と実験/3/bkVol0!6,1/te/手でわかる/表計算・プログラム/4/bkVol0!6,2/shiki/式でわかる/専門書の式/5/white} {
    \node[kgb box, minimum width=38mm, minimum height=21mm, text width=36mm, anchor=north west, fill=\kgbf] (s\kgbi) at ({\kgbi*48},0) {};
    \node[anchor=north west, inner sep=0pt] at ({\kgbi*48+2.5},-2.5) {\bkicon[1.6em]{\kgbic}{bkVol0}};
    \node[kgb lab, anchor=north west, font=\sffamily\bfseries\small] at ({\kgbi*48+11},-3.2) {\kgbn};
    \node[kgb lab, anchor=north west] at ({\kgbi*48+2.5},-11) {\kgbd};
    \node[anchor=north west, inner sep=0pt, font=\sffamily\footnotesize] at ({\kgbi*48+2.5},-16) {項目の★：\bkstars{\kgbs}};}
  \draw[kgb arr] (s0) -- (s1);
  \draw[kgb arr] (s1) -- (s2);
  \draw[bkVol0, line width=1pt] (88.5,0.5) -- (88.5,9);
  \fill[bkVol0] (88.5,9) -- (100,6.6) -- (88.5,4.2) -- cycle;
  \node[kgb lab, anchor=west, font=\sffamily\bfseries\footnotesize] at (101,6.6) {ここまでで十分};
  \draw[decorate, decoration={brace, mirror, amplitude=2mm}, bkGray, line width=0.8pt] (0,-23) -- (86,-23);
  \node[kgb lab, anchor=north] at (43,-26.5) {使う道具は★★★まで};
\end{tikzpicture}
\bkfigcaption{章の3段（6-10〜6-14、7-7〜7-10、8-1〜8-13（8-5 は2段）、9-1〜9-11）}\label{kgd-f-dan}
\end{bkfigure}

「目でわかる」（図と実験）と「手でわかる」（表計算・プログラム）は★★★以下の道具だけでたどれ、「手でわかる」の終わりに「ここまでで十分」の旗を立てています。その後「式でわかる」で★★★★★の式を作ります。「できるようになること」の3つの項目は、普通、段ごとに★★★・★★★★・★★★★★と並びます。図\ref{kgd-f-dan}の「項目の★」は、その段で到達する内容の高さ、「使う道具は★★★まで」は、そこまでたどるのに使う計算の高さを表します。

2つめは、「名前」と「実際の操作」の型です。$\iint$ や $\mathcal{L}[f(t)]$ のような記号は「操作の名前」で、記号を見ただけではまだ計算できません。実際の値は、第8巻では「ふつうの定積分（第5巻）の繰り返し」に、第9巻では4つの道具（ふつうの定積分・表を引く・部分分数・行列の計算）のどれかに直して求めます。紙面では「名前の段→直す段→楽をする段」の3つの段、場合分けの表、流れ図を使い、札の形（点線の枠・実線の枠・塗りつぶし）で段を見分けられるようにしています。これは、依頼者からの「操作の名前と実際の操作の違いは相当わかりにくく、どのような操作をすれば求められるのかも場合によりけりなので、この場合はどうすれば良いのかを丁寧にする必要がある」という指示（10月29日）を紙面にしたもので、章の3段とは別の、その1つ下の段に当たります。

\par\addvspace{4pt}
\kgbcap{第6巻からの2つの型（「名前」と「実際の操作」の型は、第6・7巻にはありません）}{kgd-t-kata}
\begin{bktable}[long, theme=kgb]{hline{Z}={0pt,white}, vline{1-Z}={Z}{0pt,white}, row{Z}={abovesep=3pt, belowsep=1pt}, hspan=minimal, rowhead=1, colspec={Q[l,m,wd=34mm] Q[l,m,wd=56mm] Q[l,m,wd=36mm]}, rowsep=2pt, colsep=4pt, cells={font=\small}, row{1}={font=\sffamily\bfseries\small}}
型 & 中身 & 使っている巻・章 \\
章の3段 & 「目でわかる」→「手でわかる」→「式でわかる」。「手でわかる」の終わりに「ここまでで十分」の旗 & 第6〜9巻：6-10〜6-14、7-7〜7-10、8-1〜8-13（8-5 は2段）、9-1〜9-11 \\
「名前」と「実際の操作」 & 名前の段→直す段→楽をする段の3つの段、場合分けの表、流れ図 & 第8巻（8-0 予告・8-6〜8-11・8-14・巻末の「積分の地図」）、第9巻（全章と巻末の「変換の地図」） \\*
\SetCell[c=3]{c} \kgbcapline & & \\
\end{bktable}
\par\addvspace{4pt}

\kgbsub{1.6}{案内役：ソラ・ミナ・デルタ}

どの巻にも、3人の案内役が登場します。表紙の裏での紹介の他、序章・部扉と、各章のふきだし（1見開きに1回まで）に出てきて、問いかけたり、次の章へつないだりします。

\begin{bkfigure}
\begin{tikzpicture}[x=1mm, y=1mm]
  \bkChar{sora}{(18,0)}{30mm}\bkChar{mina}{(62,0)}{30mm}\bkChar{delta}{(108,0)}{19mm}
  \node[kgb lab, anchor=north, font=\sffamily\bfseries\small] at (18,-1.5) {ソラ};
  \node[kgb lab, anchor=north, font=\sffamily\bfseries\small] at (62,-1.5) {ミナ};
  \node[kgb lab, anchor=north, font=\sffamily\bfseries\small] at (108,-1.5) {デルタ};
  \node[kgb lab, anchor=north, text width=40mm, align=center] at (40,-7.5) {子ども。巻が進むごとに成長し、\\第7巻からは大学生};
  \node[kgb lab, anchor=north, text width=40mm, align=center] at (108,-7.5) {三角形のロボット。\\名前は $\Delta$（デルタ）から};
\end{tikzpicture}
\bkfigcaption{3人の案内役}
\end{bkfigure}

\kgbsub{1.7}{どこから始めるか}

学年の目安（表\ref{kgd-t-vols}）は、読む時期を決めるものではありません。始める所は、次の決まりで選びます。

\begin{taisetsu}
いま学校で習っている学年の内容が含まれる、一番前の巻・章から始める
\end{taisetsu}

学年の内容がどの巻・章にあるかは、表\ref{kgb-t-gakunen}（4.3）によります。高校の数学は、1つの科目の内容がいくつもの巻に分かれているので、いま習っている単元が含まれる一番前の章から始めます。始めた章より前の章にも、まだ習っていない内容が少し含まれることがあります。分からない所は、章の「道しるべ」の「この章の前に」の章へ戻ってください。巻の順に読むのが基本ですが、自分の読める★まで読んで先へ進んでも構いません。子ども用の『はじめに読む本』の流れ図「はじめる巻を決める」も、同じ決まりで作っています。

\par\addvspace{4pt}
\kgbcap{いま習っていることと、始める巻・章}{kgd-t-start}
\begin{bktable}[long, theme=kgb]{hline{Z}={0pt,white}, vline{1-Z}={Z}{0pt,white}, row{Z}={abovesep=3pt, belowsep=1pt}, hspan=minimal, rowhead=1, colspec={Q[l,m,wd=27mm] Q[l,m,wd=33mm] Q[l,m,wd=66mm]}, rowsep=1.6pt, colsep=4pt, cells={font=\footnotesize}, row{1}={font=\sffamily\bfseries\footnotesize}}
いま習っていること & 始める巻・章 & その学年の内容がある主な章と、進め方 \\
小学1年 & 第1巻の序章（1-0） & 小1の内容は第1巻だけ。第1巻は本に直接書き込んで使います。大人が一緒に読み、「やってみよう」の実験を手伝ってください \\
小学2年 & 第1巻 1-2 & 小2の内容（4けたの数、筆算、九九、長さとかさ）は 1-2 からの第1巻 \\
小学3年 & 第1巻 1-2 & 小3の内容（万の位、3けたの筆算、割り算、小数と分数の入口）は 1-2 からの第1巻。掛け算の筆算は学校の教科書で（7.2） \\
小学4年 & 第1巻 1-15 & 小4の内容は 1-15（余りのある割り算の確かめの式）・1-18（折れ線グラフ）・1-20（角）・1-21（面積）と、第2巻。割り算の筆算は学校の教科書で（7.2） \\
小学5年 & 第2巻の序章（2-0） & 小5の内容は第2巻（2-1 から）と第4巻（図形）。第2巻からは問題をノートに解きます \\
小学6年 & 第2巻 2-4 & 小6の内容（分数の掛け算・割り算、文字、比、比例と反比例）は 2-4 から。2-1〜2-3 にも、まだ習っていない内容（負の数 2-2 など）があるので、道しるべで戻って読みます \\
中学1年 & 第2巻 2-2 & 中1の内容は第2巻（正負の数 2-2 から、文字式・一次方程式・比例と反比例）、第3巻（座標・関数・データ）、第4巻（平面図形・空間図形） \\
中学2年 & 第2巻 2-12 & 中2の内容は第2巻（式の計算 2-12・等式の変形 2-14）、第3巻（連立方程式・一次関数・箱ひげ図）、第4巻（平行線と角・合同） \\
中学3年 & 第2巻 2-6 & 中3の内容は第2巻（誤差 2-6・平方根 2-8・無理数 2-9）、第3巻（展開・二次方程式・$y=ax^2$）、第4巻（相似・三平方の定理） \\
高校の数学Ⅰ & 第2巻 2-9 & 数学Ⅰの内容は第2巻（実数 2-9）、第3巻（二次関数・データの分析）、第4巻（三角比）。第2巻の表計算（2-18）・プログラム（2-19）は、後の巻で使います \\
高校の数学Ⅱ & 習っている単元の最初の章（指数・対数は 2-7、三角関数は 4-2、微分・積分は第5巻の序章 5-0） & 数学Ⅱの内容は第2〜6巻と第9巻に分かれています（表\ref{kgb-t-gakunen}）。数学Ⅱの微分・積分を習っていて、数学Ⅲはまだの人は、第5巻から（数学Ⅲの内容も含みます） \\
高校の数学Ⅲ & 第5巻 5-3 & 極限・微分法・積分法は第5巻（5-3 から）。合成関数・逆関数は第3巻（3-3・3-12） \\
大学生・社会人の学び直し & 学び直したい科目の巻（表\ref{kgb-t-univ}） & 微分・積分は第5巻、微分方程式は第6巻、行列は第7巻。前提になる章は、各章の道しるべでたどれます \\
工学の専門書を読むため & 第6巻（全巻の乗り換え駅） & 第6〜9巻の3段の章は、「ここまでで十分」の旗まで読んでから「式でわかる」へ。第8巻・第9巻の「名前」と「実際の操作」の型で、専門書の記号を計算に直せるようになります \\*
\SetCell[c=3]{c} \kgbcapline & & \\
\end{bktable}
\par\addvspace{4pt}

\kgbsub{1.8}{電卓・表計算・プログラム（Python）}

このシリーズは「計算の量より意味」を大切にしています（編集方針5）。手で計算する練習は第1巻（九九と、足し算・引き算の筆算）で行い、第2巻からは電卓・表計算・プログラムに任せてよいことにしています。掛け算・割り算の筆算は今の版にないので、学校の教科書で練習してください（7.2）。各巻の「この本の使い方」にある「3つのやくそく」は、次の3つです。

\kgbitem{数にはいつも単位}
\kgbitem{式を書く前に表や図}
\kgbitem{くり返しの計算は電卓やコンピュータにまかせてよい}

\par\addvspace{4pt}
\kgbcap{電卓・表計算・プログラムを使い始める章}{kgd-t-tools}
\begin{bktable}[long, theme=kgb]{hline{Z}={0pt,white}, vline{1-Z}={Z}{0pt,white}, row{Z}={abovesep=3pt, belowsep=1pt}, hspan=minimal, rowhead=1, colspec={Q[l,m,wd=24mm] Q[l,m,wd=20mm] Q[l,m,wd=82mm]}, rowsep=2pt, colsep=4pt, cells={font=\footnotesize}, row{1}={font=\sffamily\bfseries\footnotesize}}
道具 & 使い始め & 内容 \\
電卓 & 2-1 から & 四則と $\sqrt{\phantom{2}}$ のキー（2-8）までは普通の電卓でできます。整数でない指数の \texttt{x\^{}y} キー（3-11）、\texttt{log} キー（3-12）、三角比と角の単位の切りかえ（4-2・4-6）からは関数電卓（または表計算）を使います \\
表計算 & 2-18 で学ぶ & セルと番地、「＝」で始まる式、式のコピー、\$ で固定する番地、上の行の結果で下の行を計算する繰り返しの計算、表からグラフ。以後の全巻で使い、第6巻のオイラー法（6-2）もこの形で解きます \\
プログラム & 2-19 で学ぶ & ブロックをつないで書く方法（Scratch など）と文字で書く方法（Python など）を並べて紹介。変数と代入、繰り返し、条件分岐、小数のずれ \\
Python & 第7巻・第8巻の本文など & 行列の計算（7-1：表計算の MMULT と Python）、2変数の関数の地図（8-1：表計算と Python）など。第6〜9巻の巻末には「プログラムのひな形」 \\*
\SetCell[c=3]{c} \kgbcapline & & \\
\end{bktable}
\par\addvspace{4pt}

第1巻は本に直接書き込んで使います。第2巻からは、問題をノートに解きます。ノートを1冊用意してください。

\kgbsub{1.9}{答えと巻末}

巻末の「こたえ」は、「じぶんでとこう」と章末問題の答えだけです（考え方は、今後作る問題集と指導書に載せる予定です）。最後に「答え合わせのしかた」があります。巻末には、他に次のものがあります。

\par\addvspace{4pt}
\kgbcap{巻末の中身}{kgd-t-back}
\begin{bktable}[long, theme=kgb]{hline{Z}={0pt,white}, vline{1-Z}={Z}{0pt,white}, row{Z}={abovesep=3pt, belowsep=1pt}, hspan=minimal, rowhead=1, colspec={Q[l,m,wd=38mm] Q[l,m,wd=88mm]}, rowsep=1.5pt, colsep=4pt, cells={font=\footnotesize}, row{1}={font=\sffamily\bfseries\footnotesize}}
巻末の部品 & 中身 \\
こたえ & 「じぶんでとこう」と章末問題の答え。最後に「答え合わせのしかた」 \\
記号カード & 各章の記号カードを集めた一覧（第4巻からはギリシャ文字の読み方、記号を書くときの注意なども） \\
よく使う定数・単位・公式 & 第1巻は「たんいのひょう」。第6巻からは「この巻でくり返し使う場面と数」と「プログラムのひな形」、第8巻は初めに見開きの「積分の地図」、第9巻は初めに「変換の地図」 \\
チェック表 & 各章の「ふりかえり」を集めた表 \\
つぎの巻へ & 次の巻の予告と「この巻の道具は、この先でこう使う」。第9巻は代わりに「物理シリーズへ」 \\
さくいん & その巻の索引。第9巻には第1〜9巻を通した「シリーズ全体のさくいん」 \\
ふろく（第1巻だけ） & コピーして使う10のわく・かずカード・ものさし・とけい・方眼 \\*
\SetCell[c=2]{c} \kgbcapline & \\
\end{bktable}
\par\addvspace{4pt}

\kgbsub{1.10}{問題集と指導書（今後の予定）}

このシリーズは教科書を優先して作りました。問題集と、教員向けの指導書は、今後作る予定で、まだ刊行していません。巻末の「こたえ」に載せていない問題の考え方や、章ごとの学年の対応は、指導書に載せる予定です。それまでの間、学年との対応は、この本の第4章と付録A・付録Bをご覧ください。

\kgbsub{1.11}{物理シリーズへ}

このシリーズの後には、物理シリーズ（全5巻）が続きます。書き方は数学シリーズの第4巻以降と同じ（ルビなし・常用漢字まで）で、各章の道しるべに、前に学ぶ数学の章の番号（6-2 など）を書きます。数学シリーズの第9巻の巻末「物理シリーズへ」では、物理の「なぜ？」とこのシリーズの章の対応表と、物理5巻の紹介を載せています。

\par\addvspace{4pt}
\kgbcap{物理シリーズ（全5巻）と、前に学ぶ数学の巻}{kgd-t-phys}
\begin{bktable}[long, theme=kgb]{hline{Z}={0pt,white}, vline{1-Z}={Z}{0pt,white}, row{Z}={abovesep=3pt, belowsep=1pt}, hspan=minimal, rowhead=1, colspec={Q[c,m,wd=16mm] Q[l,m,wd=30mm] Q[l,m,wd=78mm]}, rowsep=1.5pt, colsep=4pt, cells={font=\small}, row{1}={font=\sffamily\bfseries\small}}
物理の巻 & 巻名 & 前に学ぶ数学の巻（章ごとの指定から集計） \\
1 & 力学 & 第2〜8巻 \\
2 & 熱学 & 第2・3・5〜9巻 \\
3 & 波動 & 第3〜9巻 \\
4 & 電磁気学 & 第3〜9巻 \\
5 & 原子 & 第2〜9巻 \\*
\SetCell[c=3]{c} \kgbcapline & & \\
\end{bktable}
\par\addvspace{4pt}
}{\typeout{BKPAGE missing guide-a1}}
\IfFileExists{guide-p2a.tex}{%% guide-p2a.tex　ガイドブック　保護者・先生のための本（大人用、vol=0）第2〜10章
%% （第16回編集委員会、組 G-B、頭文字 kgb）。ルビなし。大人向けの漢字の使い方（10月30日 10:15 依頼者）。
%% 材料：tools/guide-shidou.md（RG-1）・tools/guide-series.md（RG-2）・tools/guide-taiou.md（RG-3）。
%% 数・事実は材料のとおり。材料で「未確認」とされたことは書かない（または断定しない）。
%% 部品（\kgbsec・\kgbsub・\kgbvol・表の型 kgb など）は guide-a1.tex（第1章、先に読む）の先頭で定義している。
%% 出典は guide-p2a-src.tex（親ファイルで付録 guide-p2b のあとに読む）。
\ifdefined\kgbsec\else\errmessage{guide-p2a.tex の前に guide-a1.tex を読んでください（部品の定義がそこにあります）}\fi

%% ====================================================================
\kgbsec{2}{編集方針と書き方の原則}\label{kgb-ch2}

この章では、編集委員会が第1回の委員会で決めた編集方針と、全巻で守っている書き方の原則を紹介します。シリーズの目的・巻の並び・紙面の部品は、第1章にまとめました。

\kgbsub{2.1}{編集方針の10項目}

次の10項目が編集方針です。どの巻も、この方針に沿って書かれています。

\par\addvspace{4pt}
\kgbcap{編集方針の10項目と、紙面での表れ}{}
\begin{bktable}[long, theme=kgb]{hline{Z}={0pt,white}, vline{1-Z}={Z}{0pt,white}, row{Z}={abovesep=3pt, belowsep=1pt}, hspan=minimal, rowhead=1, colspec={Q[c,m,wd=6mm] Q[l,m,wd=62mm] Q[l,m,wd=58mm]}, rowsep=2pt, colsep=4pt, cells={font=\small}, row{1}={font=\sffamily\bfseries\small}}
 & 編集方針 & 紙面では \\
1 & 量＝数値×単位から出発する。「5」ではなく「5 m」 & 2-1・2-10。各巻の「3つのやくそく」の1つめ「数にはいつも単位」 \\
2 & 具体→表→グラフ→式の順。式は「表とグラフを短く書く方法」として最後に出す & 3-0 の「はかる→表→点→線→式」の5つの段 \\
3 & 大学まで続く種を早くまく（ひき算＝変化量、かけ算＝面積、九九の表＝行列）。7本の路線でつなぐ & 1-5「かわった量」、1-13 九九の表（行列の種）、7本の路線（1.3） \\
4 & 微分方程式は「次の瞬間を決めるルール」。式で解く前に表計算で解く & 第6巻を全巻の乗りかえ駅にした（6-2 オイラー法） \\
5 & 計算の量より意味。九九と筆算は第1巻で手でできるまで練習し、第2巻からは電卓・表計算・プログラムに任せる & 電卓は 2-1 から、表計算は 2-18 から、プログラムは 2-19 から使う。今の版の第1巻の筆算は足し算・引き算だけで、掛け算・割り算の筆算はない（7.2） \\
6 & どの章も実験や現象とセット。大学の章も、家や教室で確かめられる現象から入る & 「やってみよう」「物理・工学ではこう使う」 \\
7 & 物理や工学で使わない単元は省くか軽くする（統計は戻した） & 7章にまとめた \\
8 & 大学の内容も中学生の言葉で。「目でわかる→手でわかる→式でわかる」の3段で、前の2段は★★★以下の道具だけ & 6-10〜6-14、7-7〜7-10、8-1〜8-13、9-1〜9-11（1.5） \\
9 & 専門書の書き方に慣れさせる。新しい記号は1章に2つまでで、読み方・意味・英語を記号カードで示す & 記号カード、★★★★以上の章の「専門書をのぞいてみよう」 \\
10 & 学年ではなく、★と道しるべで進む & 章の帯の★、帯の下の「道しるべ」（前に読む章・この先の章） \\*
\SetCell[c=3]{c} \kgbcapline & & \\
\end{bktable}
\par\addvspace{4pt}

\kgbsub{2.2}{書き方の原則}

新しい考えは、いつも次の順で書いています。新しい言葉は体験の後で「このことを〇〇といいます」と1回だけ紹介し、名前を最後にします。

\begin{bkfigure}
\begin{tikzpicture}[x=1mm, y=1mm]
  \foreach \kgbi/\kgbt/\kgbf [count=\kgbn] in {0/場面/white,1/ふしぎ・\\こまったこと/white,2/手を動かす/white,3/気付く/white,4/言葉・\\記号・式/bkVol0!10} {
    \node[kgb box, minimum width=22mm, minimum height=13mm, text width=20mm, fill=\kgbf] (k\kgbn) at ({\kgbi*28.5+12},0) {\kgbt};}
  \foreach \kgbn [evaluate=\kgbn as \kgbm using int(\kgbn+1)] in {1,...,4} {\draw[kgb arr] (k\kgbn) -- (k\kgbm);}
  \node[kgb lab, anchor=north, text width=24mm, align=center] at (69,-9) {測る・数える・\\表にする・図をかく};
  \node[kgb lab, anchor=north, text width=22mm, align=center] at (126,-9) {名前は最後};
\end{tikzpicture}
\bkfigcaption{新しい考えを紹介する順番（名前は最後）}
\end{bkfigure}

このほか、次のことを全巻で守っています。

\kgbitem{公式は作ってから使う。表や図で、いくつかの数で同じことが起きるのを確かめてから、決まりにします。}
\kgbitem{数字と文章だけで説明しない。新しい考えには必ず図を付けます。例題は、考え方と途中の式を1行ずつ全部書きます。}
\kgbitem{式が出てくるたびに「なぜその式になるのか」、数が出てくるたびに「なぜその数なのか」を書きます。初めての考え方は2通り以上で説明します。}
\kgbitem{例題で見せていない形を、「じぶんでとこう」や章末問題に出しません。}
\kgbitem{数学として後で困る言い方をしません（「かけ算すると大きくなる」と言い切らない、「0でわる」ことはできないと書く、など）。たとえ話には「せいかくに言うと」を添えます。}
\kgbitem{工学につなげます。本物の機械やしくみで、その章の数学がどう働いているかを見せます。}

%% ====================================================================
\kgbsec{3}{学習指導要領との関係のあらまし}\label{kgb-ch3}

\kgbsub{3.1}{現行の学習指導要領}

この本で照らし合わせたのは、次の現行の学習指導要領（算数・数学）です。授業時数は学校教育法施行規則の別表、高校の単位数は指導要領の標準単位数です。

\par\addvspace{4pt}
\kgbcap{現行の学習指導要領（算数・数学）}{}
\begin{bktable}[long, theme=kgb]{hline{Z}={0pt,white}, vline{1-Z}={Z}{0pt,white}, row{Z}={abovesep=3pt, belowsep=1pt}, hspan=minimal, rowhead=1, colspec={Q[l,m,wd=17mm] Q[l,m,wd=27mm] Q[l,m,wd=36mm] Q[l,m,wd=46mm]}, rowsep=2pt, colsep=4pt, cells={font=\small}, row{1}={font=\sffamily\bfseries\small}}
学校 & 告示と実施 & 授業時数・単位数 & 領域・科目 \\
小学校（算数） & 平成29年告示、令和2年度から全面実施 & 計1011時間（1年136、2〜6年は各175。1単位時間45分） & A 数と計算／B 図形／C 測定（1〜3年）・C 変化と関係（4〜6年）／D データの活用 \\
中学校（数学） & 平成29年告示、令和3年度から全面実施 & 計385時間（1年140、2年105、3年140。1単位時間50分） & A 数と式／B 図形／C 関数／D データの活用 \\
高等学校（数学） & 平成30年告示、令和4年度から学年進行で実施 & 数学Ⅰ 3単位（必履修）、数学Ⅱ 4、数学Ⅲ 3、数学A 2、数学B 2、数学C 2 & 数学Ⅰ・Ⅱ・Ⅲ・A・B・C（「数学活用」をなくし、数学Cを新しく置いた） \\*
\SetCell[c=4]{c} \kgbcapline & & & \\
\end{bktable}
\par\addvspace{4pt}

高校の6科目は、次の図のように組み立てられています。数学A・B・Cは、それぞれ3つの項目から適宜選んで学ぶ科目です。

\begin{bkfigure}
\begin{tikzpicture}[x=1mm, y=1mm]
  \tikzset{kgb sub/.style={kgb gbox, text width=38mm, anchor=north}}
  \node[kgb sub, draw=bkVol0, line width=1pt, minimum height=31mm] (s1) at (21,0) {{\sffamily\bfseries 数学Ⅰ}（必履修）\\(1) 数と式\\(2) 図形と計量\\(3) 二次関数\\(4) データの分析};
  \node[kgb sub, minimum height=31mm] (s2) at (70,0) {{\sffamily\bfseries 数学Ⅱ}\\(1) いろいろな式\\(2) 図形と方程式\\(3) 指数関数・対数関数\\(4) 三角関数\\(5) 微分・積分の考え};
  \node[kgb sub, minimum height=31mm] (s3) at (119,0) {{\sffamily\bfseries 数学Ⅲ}\\(1) 極限\\(2) 微分法\\(3) 積分法};
  \draw[kgb arr] (s1.east |- 0,-15.5) -- (s2.west |- 0,-15.5);
  \draw[kgb arr] (s2.east |- 0,-15.5) -- (s3.west |- 0,-15.5);
  \foreach \kgbx/\kgbn/\kgba/\kgbb/\kgbc/\kgbo in {21/数学A/図形の性質/場合の数と確率/数学と人間の活動/数学Ⅰと並行して、または後に,
    70/数学B/数列/統計的な推測/数学と社会生活/数学Ⅰの後が原則,
    119/数学C/ベクトル/平面上の曲線と複素数平面/数学的な表現の工夫/数学Ⅰの後が原則} {
    \node[kgb sub, minimum height=29mm] at (\kgbx,-45) {{\sffamily\bfseries \kgbn}\\(1) \kgba\\(2) \kgbb\\(3) \kgbc\\[2pt]{\sffamily\color{bkGray}\kgbo}};}
  \node[kgb lab, anchor=west, text=bkGray] at (2,-41) {数学A・B・C：それぞれ3つの項目から適宜選ぶ科目};
\end{tikzpicture}
\bkfigcaption{高校の数学の6科目と、学ぶ順序（矢印：数学Ⅰ→Ⅱ→Ⅲの順が原則）}
\end{bkfigure}

\kgbnote{※ どの学年でどの科目を取るかは、指導要領では決まっていません。一般的な例では、1年で数学Ⅰ・数学A、2年で数学Ⅱ・数学B、3年で数学Ⅲ・数学Cを取ります（数学Cの一部を2年に置く学校もあります）。この本で高校の「学年」を書くときは、この一般的な例によっています。}
\kgbnote{※ 令和7年度の大学入学共通テストでは、数学A・B・C のうち (3) の項目（数学と人間の活動・数学と社会生活・数学的な表現の工夫）は出題範囲外です。行列は、このうち数学Cの (3)「数学的な表現の工夫」に入っています（6.2を参照）。}

\kgbsub{3.2}{このシリーズの立場}

このシリーズは、学習指導要領の学年の並びに沿うのではなく、物理・工学で使う道筋（7本の路線）に沿って内容を並べています。そのため、指導要領の学年とは、次の3つの点で一致しない所があります。

\kgbitem{{\sffamily\bfseries 先に出る内容}　指導要領では上の学年の内容が、前の巻に出てきます。例：小4の折れ線グラフ・角・面積（第1巻 1-18・1-20・1-21）、中3の平方根（第2巻 2-8）、数学Ⅱの指数関数・対数（第3巻 3-11・3-12）。}
\kgbitem{{\sffamily\bfseries 後に出る内容}　指導要領では下の学年の内容が、後の巻に出てきます。例：小3〜小6の図形の面積・体積・合同・対称（第4巻 4-1・4-3・4-5）、小6の代表値（第3巻 3-13）。}
\kgbitem{{\sffamily\bfseries 分かれている内容}　指導要領の1つの単元が、いくつかの巻・章に分かれています（表\ref{kgb-t-tangen}）。}

\par\addvspace{4pt}
\kgbcap{1つの単元が、いくつかの章に分かれている例}{kgb-t-tangen}
\begin{bktable}[long, theme=kgb]{hline{Z}={0pt,white}, vline{1-Z}={Z}{0pt,white}, row{Z}={abovesep=3pt, belowsep=1pt}, hspan=minimal, rowhead=1, colspec={Q[l,m,wd=30mm] Q[l,m,wd=98mm]}, rowsep=1.5pt, colsep=4pt, cells={font=\small}, row{1}={font=\sffamily\bfseries\small}}
単元 & このシリーズで学ぶ章（学ぶ順） \\
面積 & 1-21（長方形）→ 4-5（平行四辺形・三角形・台形・円・柱・すい・球）。円周率は 2-9、おうぎ形は 4-2 \\
平均・データ & 2-6（測定値の平均）→ 3-13（度数分布・代表値・箱ひげ図）→ 3-14・3-15 \\
比例・関数 & 2-15（比例・反比例）→ 3-3（関数）→ 3-5（1次関数）→ 3-7（反比例のグラフ）→ 3-9（$y=ax^2$） \\
座標 & 3-2（平面）→ 4-13（空間）→ 4-15（極座標・円柱座標・球座標）。極座標は 4-2 でも紹介 \\
平方完成 & 3-8（式）・3-9（頂点）・3-10（解の公式）・3-16（最小二乗法） \\
指数 & 2-7（累乗・指数のきまり）→ 3-11（指数関数・分数の指数）→ 5-14（$e$） \\
三角比・三角関数 & 4-6 → 4-7 → 4-8 → 4-9 → 5-15 \\
数列にあたるもの & $\Sigma$ は 5-10、漸化式と等比数列・等比級数は 6-2（オイラー法）と 7-10 \\
確率 & 3-17（場合の数は確率に必要な分だけ） \\*
\SetCell[c=2]{c} \kgbcapline & \\
\end{bktable}
\par\addvspace{4pt}

\kgbsub{3.3}{用語と書き方の違い}

工学の習慣に合わせたり、子どもに分かりやすい言葉を選んだりしたため、学校の教科書と用語・書き方が違う所があります。

\par\addvspace{4pt}
\kgbcap{学校の教科書と、用語・書き方が違う所}{kgb-t-yougo}
\begin{bktable}[long, theme=kgb]{hline{Z}={0pt,white}, vline{1-Z}={Z}{0pt,white}, row{Z}={abovesep=3pt, belowsep=1pt}, hspan=minimal, rowhead=1, colspec={Q[l,m,wd=40mm] Q[l,m,wd=50mm] Q[l,m,wd=36mm]}, rowsep=1.5pt, colsep=4pt, cells={font=\small}, row{1}={font=\sffamily\bfseries\small}}
学校の教科書 & このシリーズ & 章 \\
虚数単位 $i$ & $j$（$i$ は電流に使うため） & 6-8 から \\
複素数平面 & 複素平面 & 6-9 \\
楕円 & だ円 & 4-8・4-15 ほか \\
概数 & がい数 & 2-6 \\
ドットプロット & 点を並べた図（名前は付けていない） & 2-6 \\
素因数分解 & 素数のかけ算に分ける（語はない） & 2-4 \\
最大公約数・最小公倍数 & 語はない（「公倍数で一番小さい数」など） & 2-4 \\
二項定理 & 語はない（パスカルの三角形） & 3-8・5-4（もっと知りたい） \\
二項分布・確率変数・確率分布 & 語はない（コインの表の枚数の分布） & 3-17 \\
漸化式 & 6-2 では語を使わず、7-10 で使う & 6-2・7-10 \\
変曲点 & 語はない（2回微分で山と谷を見分ける） & 5-7 \\
2進数の書き方 & $1011_{(2)}$ & 2-17 \\*
\SetCell[c=3]{c} \kgbcapline & & \\
\end{bktable}
\par\addvspace{4pt}

%% ====================================================================
\kgbsec{4}{学年と巻の対応}\label{kgb-ch4}

\kgbsub{4.1}{学年と巻の対応の図}

次の図は、指導要領の学年・科目の内容が、どの巻にあるかを示したものです。塗りつぶしたマスは、その学年・科目の内容（一部だけの場合も含みます）がその巻にあることを、太い枠は、その巻の学年の目安（4.2）を表します。高校の欄は科目ごとに分け、上に一般的な学年の例を添えました（学年の目安の枠も、この例で置いています）。一番右の「大学の内容」は、指導要領のどこにもなく、大学の工学部で学ぶ内容です。

\begin{bkfigure}
\begin{tikzpicture}[x=1mm, y=1mm]
  \kgbgridhead{1}
  \kgbgridrow{-5.0}{1}{1,2,3,4}{1}{3}
  \kgbgridrow{-13.2}{2}{2,3,4,5,6,7,8,9,10,11,12}{4}{7}
  \kgbgridrow{-21.4}{3}{4,6,7,8,9,10,11,12,13,14,15,16}{7}{9}
  \kgbgridrow{-29.6}{4}{3,4,5,6,7,8,9,10,11,12,15,16}{8}{13}
  \kgbgridrow{-37.8}{5}{11,12,13,14,16}{10}{13}
  \kgbgridrow{-46.0}{6}{12,13,14,15,16}{12}{15}
  \kgbgridrow{-54.2}{7}{13,14,15,16}{14}{16}
  \kgbgridrow{-62.4}{8}{14,15,16}{16}{16}
  \kgbgridrow{-70.6}{9}{12,14,16}{16}{16}
  \kgbgridlines{-75.2}
  % 凡例
  \fill[bkVol0!70, draw=bkInk!55, line width=0.4pt] (2,-82.7) rectangle ++(5.4,5.4);
  \node[kgb lab, anchor=west] at (9,-80) {内容がある（一部だけの場合も含む）};
  \draw[bkInk, line width=1.3pt, rounded corners=1.3mm] (78,-83.6) rectangle ++(13.8,7.2);
  \node[kgb lab, anchor=west] at (93.5,-80) {学年の目安};
\end{tikzpicture}
\bkfigcaption{学年・科目と巻の対応（マスは表\ref{kgb-t-gakunen} を基に作った）}\label{kgb-f-grid}
\end{bkfigure}

図を見ると、次のことが分かります。

\kgbitem{第1〜3巻は、ほぼ目安の学年の内容です。ただし第1巻には小4の内容（折れ線グラフ・角・面積）が、第2・3巻には中3〜高校の内容が入っています。}
\kgbitem{第4巻は、小3の図形から大学の内容（外積、円柱座標・球座標）までを含む、一番幅の広い巻です。第2巻は小2〜高校、第3巻は小4〜高校の内容を含みます。}
\kgbitem{第5巻の目安は「高校1年〜2年」ですが、数学Ⅲ（一般的な例では3年）の微分法・積分法も入っています。}
\kgbitem{第6巻からは、ほとんどが指導要領の外の内容です。高校の欄にあるますは、複素数（6-8・6-9）や、行列の章の中の等比数列・無限等比級数（7-10）などです。}

\kgbsub{4.2}{9巻の学年の目安}

各巻には、編集のときに決めた「学年の目安」があります。生徒用の紙面には学年を書かない決まりにしています（上の学年の子は恥ずかしく、下の学年の子は手を出さなくなるため）。学年の目安は、この本と、今後作る指導書にだけ書きます。各巻の学年の目安は、表\ref{kgd-t-vols}（1.2）の通りです。

\kgbsub{4.3}{学年・科目ごとの、主な巻と章}

図\ref{kgb-f-grid} のマスを、章まで示しました（章ごとの対応は付録A・付録B）。

\par\addvspace{4pt}
\kgbcap{学年・科目の内容が置かれている巻と、主な章}{kgb-t-gakunen}
\begin{bktable}[long, theme=kgb]{hline{Z}={0pt,white}, vline{1-Z}={Z}{0pt,white}, row{Z}={abovesep=3pt, belowsep=1pt}, hspan=minimal, colspec={Q[l,m,wd=12mm] Q[l,m,wd=116mm]}, rowhead=1, rowsep=1.5pt, colsep=4pt, cells={font=\footnotesize}, row{1}={font=\sffamily\bfseries\footnotesize}}
学年 & 置かれている巻（主な章） \\
小1 & 第1巻のみ \\
小2 & 第1巻（＞・＜だけ第2巻 2-2） \\
小3 & 第1巻（中心）、第2巻（□の式 2-11、k・t 2-10、万→億 2-7）、第4巻（二等辺三角形・正三角形 4-1、球 4-5） \\
小4 & 第1巻（折れ線グラフ 1-18、角 1-20、面積 1-21、余りの式 1-15）、第2巻（億・兆、概数、小数・分数、式、割合）、第3巻（ものの位置 3-2）、第4巻（平行・垂直 4-1、立体 4-5・4-13） \\
小5 & 第2巻（数と計算・速さ・割合・比例・平均・円周率）、第4巻（合同・多角形・面積・体積・角柱と円柱） \\
小6 & 第2巻（分数の乗除・文字・比・比例と反比例）、第3巻（代表値・度数分布・起こり得る場合）、第4巻（拡大図と縮図・対称・円の面積・体積） \\
中1 & 第2巻（正負の数・文字式・一次方程式・比例と反比例・素数）、第3巻（座標・関数・比例式・ヒストグラム・確率）、第4巻（平面図形・空間図形） \\
中2 & 第2巻（式の計算 2-12、等式の変形 2-14）、第3巻（連立方程式・一次関数・箱ひげ図・確率）、第4巻（平行線と角・合同・証明） \\
中3 & 第2巻（平方根・無理数・誤差・$a\times10^n$）、第3巻（展開・因数分解・二次方程式・$y=ax^2$・標本調査）、第4巻（相似・円周角・三平方の定理） \\
数Ⅰ & 第2巻（実数 2-9）、第3巻（数と式・二次関数・データの分析）、第4巻（図形と計量） \\
数A & 第2巻（2-4・2-17）、第3巻（3-17）、第4巻（4-1・4-2・4-7・4-11・4-13）、第5巻（$n!$ 5-17） \\
数Ⅱ & 第2巻（指数の拡張 2-7・2-12）、第3巻（3次の展開・恒等式・判別式・直線・指数と対数）、第4巻（距離・円の式・内分点・三角関数）、第5巻（微分・積分）、第6巻（複素数・解と係数の関係）、第9巻（分数式・3次方程式の一部） \\
数B & 第3巻（3-14・3-16・3-17）、第5巻（$\Sigma$ 5-10）、第6巻（漸化式の形 6-2）、第7巻（7-10） \\
数Ⅲ & 第3巻（合成関数・逆関数・$y=a/x$）、第5巻（極限・微分法・積分法）、第6巻（曲線の長さ：せいかくに言うと）、第7巻（$r^n$・無限等比級数）、第8巻（長さでたす線積分）、第9巻（無限級数としての $e^{At}$） \\
数C & 第3巻（双曲線の名 3-7、行列の種 3-6）、第4巻（ベクトル・極座標・媒介変数・だ円）、第6巻（複素平面 6-9）、第7巻（行列）、第8巻（だ円の式 8-1、媒介変数 8-7） \\*
\SetCell[c=2]{c} \kgbcapline & \\
\end{bktable}
\par\addvspace{4pt}

%% ====================================================================
\kgbsec{5}{巻ごとの対応}\label{kgb-ch5}

巻ごとに、指導要領のどの学年・科目に当たるかをまとめます（カードの帯は図\ref{kgb-f-grid} のその巻の行。章ごとの対応は付録B）。

\begin{kgbcard}{1}{かぞえる・はかる}
\kgbstrip{1}{1,2,3,4}{1}{3}
\kgbrow{学年の目安}{小学1年〜3年　{\sffamily ★の範囲} \kgbst{1}{1}　22章・400ページ}
\kgbrow{当たる範囲}{小1〜小3の数と計算（足し算・引き算、掛け算の意味と九九、割り算、小数と分数の入口）、測定（長さ・かさ・重さ・時刻と時間）、データ（表・棒グラフ）、図形（三角形・四角形・円）。小4の折れ線グラフ（1-18）・角（1-20）・面積（1-21）}
\kgbrow{目安とのずれ}{小4の内容が3章（★の基準「小4まで」の内）。小3の一部（そろばん、掛け算の筆算、二等辺三角形、km、□の式）は入っていないか、後の巻にある}
\kgbrow{指導要領の外}{引き算を「かわった量」（変化量）として見る見方（1-5）は、このシリーズ独自です}
\end{kgbcard}

\begin{kgbcard}{2}{数と量のことば}
\kgbstrip{2}{2,3,4,5,6,7,8,9,10,11,12}{4}{7}
\kgbrow{学年の目安}{小学4年〜中学1年　{\sffamily ★の範囲} \kgbst{1}{2}　20章・456ページ}
\kgbrow{当たる範囲}{小4〜小6の数と計算・変化と関係（割合・速さ・比・比例と反比例）・平均と概数、中1の数と式（正負の数・文字式・一次方程式）。中3の平方根・無理数（2-8・2-9）、誤差と $a\times10^n$（2-6・2-7）、数学Ⅰの実数（2-9）、数学Ⅱの0乗・負の指数（2-7・2-12）、数学Aの内容の取扱いで題材の例に挙がる二進法（2-17の2進数）}
\kgbrow{目安とのずれ}{中3〜高校の内容が5章。小4〜小6の図形は第4巻へ、データ（代表値・度数分布）は第3巻へ移した}
\kgbrow{指導要領の外}{量＝数値×単位、とびとびの数と切れ目のない数（2-1）、変化量 $\Delta$（2-3）、SI単位・組立単位・接頭語 μ（2-10）、16進数（2-17）、表計算（2-18）、プログラム（2-19）}
\end{kgbcard}

\begin{kgbcard}{3}{関数とデータ（きまりを見つける）}
\kgbstrip{3}{4,6,7,8,9,10,11,12,13,14,15,16}{7}{9}
\kgbrow{学年の目安}{中学1年〜3年　{\sffamily ★の範囲} \kgbst{1}{3}　18章・564ページ}
\kgbrow{当たる範囲}{中1〜中3の関数（比例と反比例・一次関数・$y=ax^2$）、データの活用（ヒストグラム・箱ひげ図・確率・標本調査）、数と式の一部（連立方程式・展開と因数分解・二次方程式）。数学Ⅰ（二次関数・二次不等式・データの分析）、数学Ⅱ（3次の展開・恒等式・判別式・直線の式・指数関数・対数）、数学Ⅲ（合成関数・逆関数）、数学A（組合せ・期待値）、数学B（正規分布・標本の平均のずれ）}
\kgbrow{目安とのずれ}{高校の内容が多い（3-3・3-5・3-8・3-10〜3-12・3-14〜3-17）。中学校の図形は第4巻にある}
\kgbrow{指導要領の外}{ブロック線図（3-3）、内挿・外挿（3-4）、片対数グラフ・デシベル（3-12）、最小二乗法・回帰直線（3-16。数学Bの「散布図に表したデータを関数とみなして処理する」は入口だけ）}
\end{kgbcard}

\begin{kgbcard}{4}{形と向き（図形・三角関数・ベクトル）}
\kgbstrip{4}{3,4,5,6,7,8,9,10,11,12,15,16}{8}{13}
\kgbrow{学年の目安}{中学2年〜高校2年　{\sffamily ★の範囲} \kgbst{1}{5}　16章・560ページ}
\kgbrow{当たる範囲}{小3〜小6の図形（面積・体積・合同・対称・拡大図と縮図・角柱と円柱）、中1〜中3の図形（作図をのぞく）、数学Ⅰの図形と計量（三角比・正弦定理・余弦定理）、数学Ⅱの三角関数と図形と方程式の一部（距離・円の式・内分点）、数学Aの図形の性質の一部、数学Cのベクトルと、平面上の曲線の一部（極座標・媒介変数・だ円）}
\kgbrow{目安とのずれ}{小3〜中1の図形（4-1・4-2・4-3・4-5、★〜★★）と、大学の内容（4-14・4-15、★★★★★）がある}
\kgbrow{指導要領の外}{外積（4-14）、円柱座標・球座標（4-15）、sin$^{-1}$ などで角を求めること（4-6）、ロボットの腕の順運動学・逆運動学（4-0・4-7）}
\end{kgbcard}

\begin{kgbcard}{5}{変化をつかむ（微分と積分）}
\kgbstrip{5}{11,12,13,14,16}{10}{13}
\kgbrow{学年の目安}{高校1年〜2年　{\sffamily ★の範囲} \kgbst{2}{4}　19章・492ページ}
\kgbrow{当たる範囲}{数学Ⅱの微分・積分の考え、数学Ⅲの極限・微分法・積分法（体積・曲線の長さは一部）、数学Bの $\Sigma$}
\kgbrow{目安とのずれ}{数学Ⅲは、一般的な例では高校3年。目安の「高校1年〜2年」より早い}
\kgbrow{指導要領の外}{テイラー展開（5-17）、1次近似と誤差の伝わり（5-8）、微分方程式の入口（5-13・5-18 の初期条件・一般解・特殊解）}
\end{kgbcard}

\begin{kgbcard}{6}{未来を計算する（微分方程式と複素数）}
\kgbstrip{6}{12,13,14,15,16}{12}{15}
\kgbrow{学年の目安}{高校2年〜3年　{\sffamily ★の範囲} \kgbst{2}{5}　16章・600ページ}
\kgbrow{当たる範囲}{数学Ⅱの複素数（6-8）、数学Cの複素数平面とド・モアブルの定理（6-9）}
\kgbrow{目安とのずれ}{序章を除く15章のうち13章が指導要領の外（大学の内容）。それでも★★の章（6-1・6-2と序章6-0）がある}
\kgbrow{指導要領の外}{微分方程式（6-0〜6-7、6-10〜6-15）、オイラーの公式（6-10）、双曲線関数（6-11）、数値解法（6-2）、指数分布とポアソン分布（6-4）}
\end{kgbcard}

\begin{kgbcard}{7}{まとめて計算する（行列）}
\kgbstrip{7}{13,14,15,16}{14}{16}
\kgbrow{学年の目安}{高校3年（発展）〜大学1年　{\sffamily ★の範囲} \kgbst{3}{5}　12章・492ページ}
\kgbrow{当たる範囲}{数学Cの (3)「離散グラフや行列を用いて工夫して表現することの意義」に当たる所から入る。7-10 に数学Bの等比数列・漸化式、数学Ⅲの無限等比級数}
\kgbrow{目安とのずれ}{指導要領では行列は「表現の意義」だけで、計算はすべて指導要領の外}
\kgbrow{指導要領の外}{行列の演算、一次変換、行列式、逆行列と掃き出し法、回転行列、固有値と固有ベクトル、対角化、使いみち（7-11：逆運動学・四端子定数・分布定数回路など）}
\end{kgbcard}

\begin{kgbcard}{8}{空間に広がる量（ベクトル解析）}
\kgbstrip{8}{14,15,16}{16}{16}
\kgbrow{学年の目安}{大学1年〜2年　{\sffamily ★の範囲} \kgbst{3}{5}　15章・632ページ}
\kgbrow{当たる範囲}{ほぼすべてが指導要領の外。わずかに、だ円の式（8-1）、媒介変数と曲線の長さ（8-7）が数学Ⅲ・数学Cに当たる}
\kgbrow{指導要領の外}{2変数の関数、偏微分、全微分、重積分、勾配、ベクトル場、線積分、面積分、発散、回転、ガウスの定理・ストークスの定理、ラプラシアン、熱方程式・波動方程式、マクスウェル方程式を読む（8-14 読みもの）}
\end{kgbcard}

\begin{kgbcard}{9}{波と変化を読みとく（フーリエ解析とラプラス変換）}
\kgbstrip{9}{12,14,16}{16}{16}
\kgbrow{学年の目安}{大学2年　{\sffamily ★の範囲} \kgbst{3}{5}　14章・640ページ}
\kgbrow{当たる範囲}{ほぼすべてが指導要領の外。部分分数の前の多項式の計算（9-7）、3次の特性方程式（9-9）が数学Ⅱの一部に、無限級数としての $e^{At}$（9-11）が数学Ⅲの一部に当たる}
\kgbrow{指導要領の外}{フーリエ級数・フーリエ変換・FFT、ラプラス変換・伝達関数・極と安定性、状態方程式、行列の指数関数、使いみち（9-12 読みもの）、全9巻の地図（9-13 終章）}
\end{kgbcard}

%% ====================================================================
\kgbsec{6}{追加した内容と、その理由}\label{kgb-ch6}

このシリーズには、現行の指導要領にない内容や、指導要領よりずっと詳しく扱う内容があります。この章では、それを4つに分けて示し、入れた理由を添えます。理由は、企画の依頼文、編集方針、委員会の記録、各巻の「物理・工学ではこう使う」に書かれていることだけを書きました。記録に個別の理由がないものは、その章の題材（どの工学で使うか）を書いています。

\par\Needspace{80mm}% 見出しだけがページの下に残らないように
\kgbsub{6.1}{大学の内容（第6〜9巻と、第4巻の一部）}

\par\addvspace{4pt}
\kgbcap{追加した内容（1）大学の内容}{}
\begin{bktable}[long, theme=kgb]{hline{Z}={0pt,white}, vline{1-Z}={Z}{0pt,white}, row{Z}={abovesep=3pt, belowsep=1pt}, hspan=minimal, colspec={Q[l,m,wd=30mm] Q[l,m,wd=21mm] Q[l,m,wd=23mm] Q[l,m,wd=52mm]}, rowhead=1, rowsep=1.5pt, colsep=3pt, cells={font=\footnotesize}, row{1}={font=\sffamily\bfseries\footnotesize}}
内容 & 章 & 指導要領 & 入れた理由 \\
微分方程式（立式・オイラー法・変数分離・振動・特性方程式・共振・連立） & 第6巻、5-13・5-18（入口）、9-5〜9-11 & 数学Ⅲの本文に記載なし & 企画の依頼文に「古典制御や現代制御におけるラプラス変換を用いた物理現象の微分方程式の立式や解析」。編集方針4「微分方程式は『次の瞬間を決めるルール』。式で解く前に表計算で解く」 \\
行列（演算・一次変換・行列式・逆行列・回転行列・固有値・対角化） & 第7巻、9-10・9-11 & 数学C (3) は「表現の意義」だけ & 依頼文に「動力学や逆運動学などの行列や行列式による計算、回路網理論における分布定数回路の四端子定数モデル」。委員会（機械工学）「ロボットには3次元のベクトル、外積、回転行列、逆運動学が要ります」 \\
オイラーの公式・フェーザ & 6-10、9-3・9-5 & 数学Cは極形式・ド・モアブルの定理まで & 路線④（角→sin→複素数）の終点「音・交流」へ向かう道具。題材：家のコンセントと交流の計算、三相交流の発電機 \\
双曲線関数 & 6-11 & なし & 題材：送電線のたるみ、アーチ（さかさまの懸垂線） \\
偏微分・全微分・重積分 & 8-1〜8-3・8-6 & なし & 依頼文に「電磁気学や流体力学などの3次元のベクトル解析」。委員会（数学）「専門書を読むのに欠かせない数学を足してください。全微分、重積分、ラプラシアン」 \\
ベクトル解析（勾配・発散・回転・線積分・面積分・ガウスとストークスの定理） & 8-4〜8-12、8-14 & 数学Cのベクトルは内積まで & 依頼文（同上）。委員会（電気工学）「式を見て『わき出しとうずの式だ』と言えるようにはしたい。専門書で一番つまずくのは、記号の読み方です」 \\
外積、円柱座標・球座標 & 4-14・4-15、8-6 & 「外積」の語はない & 第1回の委員会で「足すもの」に外積といろいろな座標を挙げた。題材：電気モーター、静止衛星 \\
偏微分方程式（熱方程式・波動方程式・ラプラス方程式） & 8-12・8-13、9-1 & なし & 路線①（変化）の終点「流体・熱」。題材：スマートフォンの熱をにがす、地震の波、楽器の弦 \\
フーリエ級数・フーリエ変換（FFT・サンプリング定理） & 9-1〜9-4、9-12 & 数学Ⅲは無限等比級数まで & 依頼文に「フーリエ級数展開など機械工学、電気工学などで必須の知識」 \\
ラプラス変換・伝達関数・極と安定性 & 9-5〜9-9 & なし & 依頼文（微分方程式の行と同じ）。委員会（電気工学）「ラプラス変換は『伝達関数』と『極と安定性』までにします」 \\
状態方程式・行列の指数関数 & 9-10・9-11 & なし & 委員会（電気工学）「状態方程式は、行列と微分方程式の話として残します」 \\
数値解法（オイラー法・差分法・掃き出し法・べき乗法・DFT と FFT） & 6-2・7-5・7-8・8-13・9-4 & 数値解法の項目はない（コンピュータの活用は各所にある） & 編集方針4・5。委員会（機械工学）「行列のかけ算もフーリエ係数も、表計算なら中学生が自分で出せます」 \\*
\SetCell[c=4]{c} \kgbcapline & & & \\
\end{bktable}
\par\addvspace{4pt}

\kgbsub{6.2}{行列の扱いの移り変わり}

行列は、指導要領の改訂のたびに扱いが変わってきた内容です。現行の指導要領では、行列は数学Cの (3)「数学的な表現の工夫」の中で「表現の道具」として挙げられているだけで、行列の和・積、逆行列、行列式、一次変換、固有値の語は本文にありません。工学では、ロボットの動き、回路網、制御、振動の解析など、行列の計算を使わない分野はほとんどありません。このシリーズが第7巻を丸ごと行列に当てたのは、このためです。

\begin{bkfigure}
\begin{tikzpicture}[x=1mm, y=1mm]
  \draw[bkGray, line width=1.2pt, -{Stealth[length=3mm]}] (2,0) -- (138,0);
  \foreach \kgbx/\kgbt/\kgbb in {
    22/平成11年告示/行列は{\sffamily\bfseries 数学C}の内容,
    70/平成21年告示/数学Cをなくす。行列は「{\sffamily\bfseries 数学活用}」で「図，表，行列及び離散グラフなどを用いて，事象を数学的に表現し考察すること」,
    118/平成30年告示（現行）/数学Cを新しく置く。(3)「数学的な表現の工夫」で「離散グラフや{\sffamily\bfseries 行列}を用いて工夫して表現することの意義」} {
    \fill[bkVol0] (\kgbx,0) circle (1.8);
    \node[kgb lab, anchor=south, font=\sffamily\bfseries\footnotesize] at (\kgbx,3) {\kgbt};
    \node[kgb gbox, anchor=north, text width=40mm] at (\kgbx,-4) {\kgbb};}
\end{tikzpicture}
\bkfigcaption{高校の指導要領での行列の扱い（文部科学省の資料による）}
\end{bkfigure}

\kgbnote{※ 文部科学省は、学校設定科目などで使える高等学校数学科教材「行列入門」を公開しています。現行の指導要領について、この資料は行列の扱いが「従前と比べて大きな変更はありません」としています。}

\kgbsub{6.3}{高校の範囲を超える内容（主に第3〜6巻）}

\par\addvspace{4pt}
\kgbcap{追加した内容（2）高校の範囲を超える内容}{}
\begin{bktable}[long, theme=kgb]{hline{Z}={0pt,white}, vline{1-Z}={Z}{0pt,white}, row{Z}={abovesep=3pt, belowsep=1pt}, hspan=minimal, rowhead=1, colspec={Q[l,m,wd=30mm] Q[l,m,wd=21mm] Q[l,m,wd=23mm] Q[l,m,wd=52mm]}, rowsep=1.5pt, colsep=3pt, cells={font=\footnotesize}, row{1}={font=\sffamily\bfseries\footnotesize}}
内容 & 章 & 指導要領 & 入れた理由 \\
テイラー展開、$n$ 回微分 & 5-17 & 数学Ⅲは無限等比級数まで。第$n$次導関数は用語にない & 第1回の委員会で「足すもの」に挙げた。題材：振り子とばね、トランジスタと「小信号」 \\
最小二乗法・回帰直線 & 3-16、7-5 & 数学Bは入口だけ & 委員会（数学）「統計を入れるなら、平均だけでは足りない。標準偏差、正規分布、最小二乗法まで」。（電気工学）「計測とノイズを扱うには、確率と統計が欠かせません」 \\
指数分布・ポアソン分布 & 6-4 & なし & 題材：放射性炭素の年代 \\
1次近似・誤差の伝わり & 5-8・8-3 & 項目なし & 題材：カーナビの推測航法。5-8 は 6-2 のオイラー法の予告 \\
片対数グラフ・デシベル & 3-12 & なし & 題材：ダイオードの電流、ボード線図 \\
ブロック線図 & 3-3・9-8 & なし & 題材：温度センサー（温度を入れると電圧が出てくる箱） \\
sin$^{-1}$ などで角を求める & 4-6 & なし & 題材：すべり出す角で摩擦をはかる \\
順運動学・逆運動学 & 4-0・4-7・7-6・7-11 & なし（工学） & 委員会（機械工学）「余弦定理は逆運動学の準備として入れてください」 \\*
\SetCell[c=4]{c} \kgbcapline & & & \\
\end{bktable}
\par\addvspace{4pt}

\kgbsub{6.4}{算数・数学の外にある土台}

\par\addvspace{4pt}
\kgbcap{追加した内容（3）算数・数学の外にある土台}{}
\begin{bktable}[long, theme=kgb]{hline{Z}={0pt,white}, vline{1-Z}={Z}{0pt,white}, row{Z}={abovesep=3pt, belowsep=1pt}, hspan=minimal, rowhead=1, colspec={Q[l,m,wd=30mm] Q[l,m,wd=21mm] Q[l,m,wd=23mm] Q[l,m,wd=52mm]}, rowsep=1.5pt, colsep=3pt, cells={font=\footnotesize}, row{1}={font=\sffamily\bfseries\footnotesize}}
内容 & 章 & 指導要領 & 入れた理由 \\
量＝数値×単位、SI単位・組立単位・接頭語、単位つきの換算 & 2-1・2-3・2-10 & 小3で k・m に「触れる」程度 & 編集方針1「量＝数値×単位から出発する」。題材：単位のとりちがえで消えた火星探査機 \\
有効数字・誤差（±）・ずれの割合 & 2-6、3-14 & 中3の「誤差、近似値」 & 路線⑦（目もり→誤差→統計）。題材：部品の寸法と公差 \\
変化量 $\Delta$・1あたりの量・量の掛け算 & 2-3 & $\Delta$ はどの学年の用語・記号にもない & 編集方針3「大学まで続く種を早くまく（ひき算＝変化量、かけ算＝面積…）」 \\
表計算・プログラム & 2-18・2-19（以後の全巻で使う） & 「コンピュータなどを用いて」（各所）、小5のプログラミングの例 & 編集方針5。第1回の委員会で「表計算とプログラム」を足すと決めた \\
16進数（2進数は 6.5） & 2-17 & 16進数はない & 題材：センサーの値を数にする（分解能）、マイコンボード \\*
\SetCell[c=4]{c} \kgbcapline & & & \\
\end{bktable}
\par\addvspace{4pt}

\kgbsub{6.5}{指導要領と同じ項目を、詳しく扱う所}

\par\addvspace{4pt}
\kgbcap{追加した内容（4）同じ項目を、詳しく扱う所}{}
\begin{bktable}[long, theme=kgb]{hline{Z}={0pt,white}, vline{1-Z}={Z}{0pt,white}, row{Z}={abovesep=3pt, belowsep=1pt}, hspan=minimal, rowhead=1, colspec={Q[l,m,wd=24mm] Q[l,m,wd=14mm] Q[l,m,wd=42mm] Q[l,m,wd=46mm]}, rowsep=1.5pt, colsep=3pt, cells={font=\footnotesize}, row{1}={font=\sffamily\bfseries\footnotesize}}
項目 & 章 & 指導要領の扱い & このシリーズの扱い \\
反比例 & 2-15・3-7 & 小6は「知る」程度 & $xy$ が一定、グラフ、$pV$ が一定 \\
三平方の定理 & 4-4 & 中3は「証明できることを知る」 & 2通りの証明を本文で見せる \\
正弦定理・余弦定理 & 4-7 & 数学Ⅰ & ロボットの腕の順運動学・逆運動学、力の三角形 \\
加法定理 & 4-9 & 数学Ⅱ & 2通りで導く、積と和の式、うなり、sin$^2$ の平均（実効値 5-15） \\
内積 & 4-12 & 数学C & 仕事、正射影、相関係数は偏差ベクトルの向きのそろいぐあい \\
微分 & 5-4 & 数学Ⅱは三次までの関数が中心、極限は直観的に & $x^n$ の一般の形、表から規則を見つける \\
積分 & 5-10 & 数学Ⅱは二次までの関数が中心 & 区分求積・$\Sigma$ から定積分へ \\
標準偏差・相関 & 3-14・3-15 & 数学Ⅰ & 共分散、$\sigma/\sqrt{n}$、単位を変えたときの $\sigma$ \\
二次方程式 & 3-10 & 中3は実数解をもつものだけ & 判別式、2次不等式。虚数の解は 6-8 \\
指数・対数 & 3-11・3-12 & 数学Ⅱ & 片対数グラフ、デシベル、半減期 \\
2進数 & 2-17 & 数学Aの内容の取扱いで、数学と人間の活動の題材の例に「二進法」 & 2進数と10進数の行き来、ビットの数と何通り（$2^n$）、バイト、分解能。16進数も扱う（16進数は指導要領の外：6.4） \\*
\SetCell[c=4]{c} \kgbcapline & & & \\
\end{bktable}
\par\addvspace{4pt}

\par\Needspace{165mm}% 見出しから図までを同じページに
\kgbsub{6.6}{まとめ：路線ごとに見た、指導要領の先}

6.1〜6.5 の内容を、7本の路線ごとにまとめます。指導要領の中でどこまで進むかと、このシリーズでどこまで進むかを並べると、次のようになります。

\begin{bkfigure}
\begin{tikzpicture}[x=1mm, y=1mm]
  \node[kgb lab, font=\sffamily\bfseries\footnotesize] at (49,6) {指導要領の中};
  \node[kgb lab, font=\sffamily\bfseries\footnotesize] at (109,6) {このシリーズで、その先へ};
  \foreach \kgbr/\kgbn/\kgba/\kgbb/\kgbe [count=\kgbi from 0] in {
    1/変化/ひき算→変化の割合・傾き（中2）\\→微分（数Ⅱ・数Ⅲ）/微分方程式（第6巻）\\→偏微分・勾配・熱方程式（第8巻）/流体・熱,
    2/たし合わせ/かけ算→面積（小4〜）\\→積分（数Ⅱ・数Ⅲ）/重積分・線積分・面積分、\\ガウスの定理（第8巻）/電磁気,
    3/倍々/累乗→指数関数・対数関数（数Ⅱ）\\→$e$（数Ⅲ）/時定数（6-4）\\→ラプラス変換・伝達関数（第9巻）/制御,
    4/回転と波/角（小4）→三角関数（数Ⅱ）\\→複素数平面（数C）/オイラーの公式（6-10）\\→フーリエ解析（第9巻）/音・交流,
    5/表/表とグラフ（小2〜）\\→行列で表現する意義（数C）/行列の計算・固有値（第7巻）\\→状態方程式（第9巻）/回路網・制御,
    6/矢印/数直線・座標\\→ベクトルと内積（数C）/外積（4-14）\\→ベクトル解析（第8巻）/ロボット・電磁気,
    7/はかる/目もり→データの分析（数Ⅰ）\\→統計的な推測（数B）/最小二乗法（3-16）\\→スペクトル・FFT（9-4）/計測・診断} {
    \pgfmathsetmacro\kgby{-\kgbi*18.5}
    \node[anchor=west, inner sep=0pt] at (0,\kgby) {\bkrouteicon[1.3em]{\kgbr}};
    \node[kgb lab, anchor=west, font=\sffamily\bfseries\footnotesize] at (5.8,\kgby) {\kgbn};
    \node[kgb gbox, anchor=west, text width=44mm, minimum height=15.5mm] (g\kgbi) at (25,\kgby) {\kgba};
    \node[kgb tbox, anchor=west, text width=46mm, minimum height=15.5mm] (t\kgbi) at (82,\kgby) {\kgbb};
    \draw[kgb arr] (g\kgbi) -- (t\kgbi);
  }
\end{tikzpicture}
\bkfigcaption{路線ごとの、指導要領の中の到達点と、このシリーズの到達点}
\end{bkfigure}


%% ====================================================================
\kgbsec[含めていない・軽く扱う内容と、補い方]{7}{含めていない・軽く扱う内容と、\\補い方}\label{kgb-ch7}

このシリーズは、物理・工学で使う道筋を優先したため、指導要領にあっても含めていない内容、軽く扱っている内容があります。正直にお示しし、学校の教科書・問題集での補い方を添えます。

\kgbsub{7.1}{方針として軽くした内容}

次の内容は、編集委員会が構成を決めるときに、方針として軽くしたものです。「理由」は委員会の記録の言葉です。

\par\addvspace{4pt}
\kgbcap{方針として軽くした内容}{}
\begin{bktable}[long, theme=kgb]{hline{Z}={0pt,white}, vline{1-Z}={Z}{0pt,white}, row{Z}={abovesep=3pt, belowsep=1pt}, hspan=minimal, rowhead=1, colspec={Q[l,m,wd=27mm] Q[l,m,wd=25mm] Q[l,m,wd=42mm] Q[l,m,wd=32mm]}, rowsep=1.5pt, colsep=3pt, cells={font=\footnotesize}, row{1}={font=\sffamily\bfseries\footnotesize}}
内容 & 指導要領 & このシリーズでの扱い & 理由 \\
図形の証明・作図 & 小3（作図）、中1（作図）、中2・中3（証明）、数学A（作図） & 証明は 4-1 の三角形の内角の和、4-4 の三平方の定理（2通り）、4-9 の加法定理の導き方。作図は 4-2 で垂直二等分線を紙を折って見つける程度。作図の手順の節はない & 定理の意味が分かれば十分 \\
整数の性質（倍数判定・不定方程式・互除法） & 小5、中1、数学A & 約分に必要な分だけ（2-4：素数・公約数・公倍数）。偶数・奇数の定義は 2-9 の「もっと知りたい」 & 物理・工学でほとんど使わない \\
場合の数（順列・組合せ） & 数学A & 確率に必要な範囲だけ（3-17：樹形図・パスカルの三角形・組合せ） & 計算演習より、ばらつきの見方が大事 \\
積分の計算演習（回転体の体積・曲線の長さ） & 数学Ⅲ & 積分の意味（面積・仕事・流量）にしぼる。体積は 5-10 の例題、曲線の長さは 6-11・8-7 に式がある & 計算は数式処理ソフトで確かめられる \\
集合と論理 & 数学Ⅰ & 4-1 の「逆と反例」（仮定・結論）だけ。集合は扱わない & 必要最小限 \\
数列の一般論（等差数列・数学的帰納法） & 数学B & $\Sigma$ は 5-10、漸化式は 6-2（オイラー法の形）と 7-10。等差数列・帰納法はない & 微分方程式の数値解に吸収できる \\*
\SetCell[c=4]{c} \kgbcapline & & & \\
\end{bktable}
\par\addvspace{4pt}

\kgbsub{7.2}{今の版で扱っていない、または一部だけの内容}

このシリーズは、物理や工学で使わない単元は省くか軽くする方針（編集方針7）で作っており、次の内容は扱っていません（一部だけ扱うものを含みます）。補い方の欄の章を読んだ後に、学校の教科書・問題集で学んでください（図\ref{kgb-f-hokai}）。特に、掛け算・割り算の筆算は、今の版の第1巻では扱っていません。第1巻を使うときは、小学校の教科書で筆算の練習をしてください。

\par\addvspace{4pt}
\kgbcap{今の版で扱っていない、または一部だけの内容}{}
\begin{bktable}[long, theme=kgb]{hline{Z}={0pt,white}, vline{1-Z}={Z}{0pt,white}, row{Z}={abovesep=3pt, belowsep=1pt}, hspan=minimal, colspec={Q[l,m,wd=33mm] Q[l,m,wd=20mm] Q[l,m,wd=56mm] Q[l,m,wd=17mm]}, rowhead=1, rowsep=1.5pt, colsep=3pt, cells={font=\footnotesize}, row{1}={font=\sffamily\bfseries\footnotesize}}
内容 & 指導要領 & このシリーズでの扱い & 補い方 \\
掛け算・割り算の筆算 & 小3・小4 & 筆算は、足し算・引き算（1-7）だけ。掛け算は九九（1-12・1-13）と「3×10」程度、割り算は九九の範囲（1-14・1-15）。23×14 の筆算は 3-8 のコラムで「展開」として見せるだけ & 1-13・1-15 の後 \\
そろばん & 小3・小4 & 扱っていない & 1-15 の後 \\
円グラフ・帯グラフ & 小5 & 扱っていない & 2-5 の後 \\
独立な試行・条件付き確率 & 数学A & 扱っていない & 3-17 の後 \\
統計的な推測（確率分布・二項分布の名前、区間推定、信頼区間、有意水準） & 数学B & 正規分布と標本調査（3-17）、$\sigma/\sqrt{n}$（3-14）。仮説検定は 3-17 の「もっと知りたい」だけ。区間推定はない & 3-17 の後 \\
三次の因数分解・因数定理・高次方程式 & 数学Ⅱ & 3-8 で $(a+b)^3$ の展開、9-9 で3次の特性方程式を扱う程度 & 3-8 の後 \\
多項式の除法・分数式 & 数学Ⅱ & 9-7 で部分分数の前に1行使うだけ & 3-8 の後（9-7 の前に） \\
等式・不等式の証明 & 数学Ⅱ & 扱っていない（恒等式・係数比較は 3-8） & 3-8 の後 \\
二項定理 & 数学Ⅱ（用語） & 語はない。パスカルの三角形（3-8・5-4 の「もっと知りたい」、3-17） & 3-8 の後 \\
外分点・軌跡・不等式の表す領域 & 数学Ⅱ & 扱っていない（内分点は 4-11） & 4-11 の後 \\
分数関数・無理関数のグラフ & 数学Ⅲ & $y=a/x$（3-7）だけ & 3-7 の後 \\
凹凸・変曲点 & 数学Ⅲ & 5-7「2回微分で、山と谷を見分ける」まで。変曲点の語はない & 5-7 の後 \\
三角形・円の性質（内心・メネラウスの定理・チェバの定理・方べき・接弦定理など） & 数学A & 外心（4-2）、重心（4-11）、内接円（4-7 の「もっと知りたい」）だけ & 4-2 の後 \\
二次曲線（焦点・準線） & 数学C & 放物線（3-9）、だ円（4-8 の「もっと知りたい」・4-15・8-1）、双曲線は $y=a/x$ のグラフとして（3-7）。焦点・準線で3つをまとめる扱いはない & 4-15 の後 \\
数学と人間の活動・数学と社会生活 & 数学A・数学B の (3) & 2-17（2進数）、3-16（最小二乗法）が近い内容 & 学校の授業で \\
逆三角関数の微分・不定形の極限（ロピタルの定理） & 大学（微分積分学） & 扱っていない（sin$^{-1}$ は 4-6 で電卓を使って角を求めるだけ） & 5-15 の後に大学の教科書で \\*
\SetCell[c=4]{c} \kgbcapline & & & \\
\end{bktable}
\par\addvspace{4pt}

\par\Needspace{160mm}% 見出しと図を同じページに
\kgbsub{7.3}{補い方：どの章の後に学ぶとよいか}

7.1・7.2 の内容は、学校の教科書・問題集で補うことができます。次の図は、このシリーズのどの章を読んだ後に補うとつながりやすいかの例です。かっこの中は、指導要領で学ぶ学年・科目です。

\begin{bkfigure}
\begin{tikzpicture}[x=1mm, y=1mm]
  \draw[bkGray, line width=1.6pt] (11,2) -- (11,-133.4);
  \foreach \kgby/\kgbc/\kgbv/\kgbt [count=\kgbi from 0] in {
    0.0/1-13/1/掛け算の筆算（小3）,
    -8.6/1-15/1/割り算の筆算（小4）、そろばん（小3・小4）,
    -17.2/2-4/2/最大公約数・最小公倍数・素因数分解の名前、偶数と奇数（小5・中1）,
    -25.8/2-5/2/円グラフ・帯グラフ（小5）,
    -34.4/3-7/3/分数関数・無理関数のグラフ（数学Ⅲ）,
    -45.8/3-8/3/三次の因数分解・多項式の除法・分数式・因数定理・二項定理・等式と不等式の証明（数学Ⅱ）,
    -59.8/3-17/3/順列・和の法則と積の法則・独立な試行・条件付き確率（数学A）、確率分布・区間推定（数学B）,
    -71.2/4-1/4/図形の証明・作図（中1〜中3・数学A）、集合と命題（数学Ⅰ）,
    -79.8/4-2/4/三角形・円の性質（内心・方べきなど）（数学A）,
    -88.4/4-11/4/外分点・軌跡・不等式の表す領域（数学Ⅱ）,
    -97.0/4-15/4/二次曲線の焦点・準線（数学C）,
    -105.6/5-7/5/凹凸・変曲点（数学Ⅲ）,
    -114.2/5-10/5/等差数列・数学的帰納法（数学B）,
    -122.8/5-15/5/逆三角関数の微分・不定形の極限（大学の微分積分学）,
    -131.4/5-16/5/回転体の体積・曲線の長さの計算練習（数学Ⅲ）} {
    \node[fill=bkVol\kgbv, text=bkOn\kgbv, rounded corners=1mm, font=\sffamily\bfseries\footnotesize, minimum width=12mm, inner ysep=1mm] (c\kgbi) at (11,\kgby) {\kgbc};
    \node[kgb gbox, anchor=west, text width=106mm, minimum height=7mm, inner ysep=1.1mm] (b\kgbi) at (27,\kgby) {\kgbt};
    \draw[kgb arr] (c\kgbi) -- (b\kgbi);}
  \node[kgb lab, anchor=south, font=\sffamily\bfseries\footnotesize] at (11,4) {このシリーズの章};
  \node[kgb lab, anchor=south west, font=\sffamily\bfseries\footnotesize] at (27,4) {その後に、学校の教科書・問題集で補う内容};
\end{tikzpicture}
\bkfigcaption{補い方の例（章を読んだ後に、学校の教科書・問題集で学ぶ）}\label{kgb-f-hokai}
\end{bkfigure}

多項式の割り算は、第9巻 9-7 でラプラス変換の部分分数に使います。9-7 を読む前に、学校の教科書で学んでおくと安心です。


\kgbsub{7.4}{入試との関係}

このシリーズは、入試を解くための問題演習を目的にしていません（1.1、企画の依頼文）。高校入試・大学入試の準備には、この章の内容を含め、学校の教科書・問題集と併せて使ってください。

なお、令和7年度の大学入学共通テストの数学では、『数学Ⅰ，数学Ａ』の数学Aは「図形の性質」「場合の数と確率」の2項目から、『数学Ⅱ，数学Ｂ，数学Ｃ』は数学Bの「数列」「統計的な推測」と数学Cの「ベクトル」「平面上の曲線と複素数平面」の4項目から3項目を選んで解答する形で出題され、数学A・B・Cの (3) の項目は出題されません。7.1・7.2 の内容のうち、図形の証明・作図、三角形・円の性質、場合の数、独立な試行・条件付き確率、数列、統計的な推測、二次曲線は、これらの出題範囲の項目に含まれます。

%% ====================================================================
\kgbsec{8}{大学の工学部の数学との対応}\label{kgb-ch8}

このシリーズの第3巻の統計の章、第4巻の後半、第5〜9巻は、大学の工学部の1〜2年で普通学ぶ数学の土台になっています。科目の名前・中身・学年は大学や学科ごとに違うので、ここでは実際のシラバスの例から拾った「よくある範囲」と照らしました（表\ref{kgb-t-univ}・図\ref{kgb-f-univ}）。


\par\addvspace{4pt}
\kgbcap{大学の工学部の数学の科目と、このシリーズの章}{kgb-t-univ}
\begin{bktable}[long, theme=kgb]{hline{Z}={0pt,white}, vline{1-Z}={Z}{0pt,white}, row{Z}={abovesep=3pt, belowsep=1pt}, hspan=minimal, colspec={Q[l,m,wd=24mm] Q[l,m,wd=38mm] Q[l,m,wd=30mm] Q[l,m,wd=34mm]}, rowhead=1, rowsep=1.5pt, colsep=3pt, cells={font=\footnotesize}, row{1}={font=\sffamily\bfseries\footnotesize}}
大学の科目 & 代表的な内容（シラバスの例） & このシリーズの章 & このシリーズで扱わないこと \\
微分積分学（1変数） & 極限、逆三角関数とその微分、高次導関数、テイラーの定理、積分法 & 第5巻（5-3〜5-18）、テイラー展開 5-17 & $\varepsilon$-$\delta$ 論法は使わない（極限は、数値で近づくようすで示す）。逆三角関数の微分、不定形の極限（ロピタルの定理）は扱わない（7.2） \\
微分積分学（多変数） & 偏微分、全微分、ヤコビアン、極値問題、重積分 & 8-1〜8-4・8-6、ヤコビ行列 7-2 & ― \\
線形代数学Ⅰ & 行列の演算、行列式、連立一次方程式（掃き出し法）、逆行列 & 7-1・7-3〜7-5 & ― \\
線形代数学Ⅱ & ベクトル空間、線形写像、固有値・固有ベクトル、対角化、直交行列 & 7-2・7-6〜7-10、9-10・9-11 & 抽象的なベクトル空間（行列は2×2と3×3の具体的なもの） \\
常微分方程式 & 変数分離形、一階線形、高階線形・連立線形 & 第6巻、5-13・5-18（入口）、9-7・9-10・9-11 & ― \\
ベクトル解析 & 勾配・発散・回転、線積分・面積分、ガウス・ストークスの定理、直交曲線座標 & 4-14・4-15、8-4〜8-12・8-14 & ― \\
偏微分方程式 & 波動方程式・熱伝導方程式・ラプラス方程式（変数分離・フーリエ級数による解法） & 8-12・8-13、9-1 & ― \\
複素関数論 & 正則関数、コーシーの積分定理、留数定理 & 6-8〜6-10（複素数とオイラーの公式まで） & 複素関数論（留数定理など）。ラプラス逆変換は変換表と部分分数で行う \\
フーリエ解析 & フーリエ級数、複素フーリエ級数、フーリエ変換 & 9-1〜9-4・9-12 & ― \\
ラプラス変換 & ラプラス変換の性質、微分方程式への応用 & 9-5〜9-9 & ― \\
確率統計 & 確率変数、二項・ポアソン・正規分布、推定と検定 & 3-13〜3-17、6-4 & 区間推定（仮説検定は 3-17 の「もっと知りたい」だけ） \\*
\SetCell[c=4]{c} \kgbcapline & & & \\
\end{bktable}
\par\addvspace{4pt}

\kgbnote{※ 「代表的な内容」は、東北大学・京都工芸繊維大学・北海道大学・神戸大学・東京農工大学のシラバスなどから拾った例です（出典の U1〜U6・U8・U9）。確率統計は、工学部向けのシラバスの例を確かめられなかったため、数学科向けのシラバスを目安にしました。}
\kgbnote{※ 第9巻の終章（9-13）にも、電気回路・制御工学・信号処理・機械力学・流体力学・電磁気学などの専門の科目と、その土台になる巻・章の表があります。}

\begin{bkfigure}
\begin{tikzpicture}[x=1mm, y=1mm]
  \foreach \kgbv [count=\kgbk] in {3,...,9} {\node[inner sep=0pt] at ({57+(\kgbk-1)*12.6},4) {\kgbvol{\kgbv}};}
  \foreach \kgbs/\kgbl/\kgbf [count=\kgbi from 0] in {
    微分積分学（1変数）/1年前期/{5/5},
    微分積分学（多変数）/1年後期/{7/4,8/5},
    線形代数学Ⅰ/1年前期/{7/5},
    線形代数学Ⅱ/1年後期/{7/5,9/4},
    常微分方程式/1〜2年/{5/4,6/5,9/5},
    ベクトル解析/2年/{4/4,8/5},
    偏微分方程式/2年/{8/5,9/4},
    複素関数論/2年/{6/4},
    フーリエ解析/2年/{9/5},
    ラプラス変換/2年/{9/5},
    確率統計/1〜2年/{3/5,6/4}} {
    \pgfmathsetmacro\kgby{-2-\kgbi*7.6}
    \ifodd\kgbi\else\fill[bkPaper] (0,{\kgby-3.8}) rectangle (140,{\kgby+3.8});\fi
    \node[kgb lab, anchor=west, font=\sffamily\bfseries\footnotesize] at (1,\kgby) {\kgbs};
    \node[kgb lab, anchor=east, text=bkGray] at (51,\kgby) {\kgbl};
    \foreach \kgbv/\kgbm in \kgbf {
      \pgfmathsetmacro\kgbx{57+(\kgbv-3)*12.6}
      \ifnum\kgbm=5 \fill[bkVol\kgbv, draw=bkInk, line width=0.6pt] (\kgbx,\kgby) circle (2.4);
      \else \draw[bkVol\kgbv!80!black, line width=1.4pt, fill=white] (\kgbx,\kgby) circle (2.1);\fi}
  }
  \fill[bkVol0, draw=bkInk, line width=0.6pt] (60,-89) circle (2.4);
  \node[kgb lab, anchor=west] at (64,-89) {土台の中心};
  \draw[bkVol0, line width=1.4pt, fill=white] (92,-89) circle (2.1);
  \node[kgb lab, anchor=west] at (96,-89) {一部・入口};
\end{tikzpicture}
\bkfigcaption{大学の工学部の数学の科目と、その土台になる巻（学年はよくある例）}\label{kgb-f-univ}
\end{bkfigure}

%% ====================================================================
\kgbsec{9}{★と学年の目安の読み方}\label{kgb-ch9}

\kgbsub{9.1}{★は「このシリーズの中での難しさ」の目安}

各章と「できるようになること」の各項目には、★を1〜5個付けています。★は、このシリーズの中での難しさ、つまり、その章でどの道具を使うかの目安です。編集を始めるとき、★の数を次のように決めました。これは、おおよその目安です。

\begin{bkcenter}
\begin{bktable}{colspec={Q[l,m,wd=30mm] Q[l,m,wd=46mm] Q[l,m,wd=40mm]}, rowsep=1.5pt, colsep=4pt, cells={font=\small}, row{1}={font=\sffamily\bfseries\small}}
★ & 内容の高さ（おおよその目安） & 子どもへの言い方（山登り） \\
\bkstars{1} & 小4まで & はじめの一歩 \\
\bkstars{2} & 小5〜中1 & ふもとを歩く \\
\bkstars{3} & 中2〜高1 & 山を登る \\
\bkstars{4} & 高2〜高3 & 尾根を歩く \\
\bkstars{5} & 大学（工学部） & 頂上をめざす \\
\end{bktable}
\bktabcaption{★の数と、内容の高さのおおよその目安}
\end{bkcenter}

ただし、★は指導要領の学年と一対一には対応していません。同じ指導要領の内容でも、このシリーズの書き方でたどれる高さに合わせて★を付けたため、目安の表より低い★が付いている章があります。たとえば、数学Ⅱ（一般的な例では高校2年）の内容は、目安の表どおりなら★★★★ですが、次の図のように★★〜★★★の章がいくつもあります。反対に、目安の表より高い★が付いている章は見つかりませんでした。なお、章の★が低くても、「できるようになること」の項目の★は目安の表どおりの章があります（4-2 の弧度法の項目は★★★★。9.4）。

\begin{bkfigure}
\begin{tikzpicture}[x=1mm, y=1mm]
  \foreach \kgbs in {1,...,5} {\draw[bkLine, line width=0.4pt] (18,{\kgbs*9}) -- (138,{\kgbs*9});
    \node[anchor=east, inner sep=0pt, font=\sffamily\footnotesize] at (17,{\kgbs*9}) {\bkstars{\kgbs}};}
  \fill[bkVol0!10] (18,{4*9-3.5}) rectangle (138,{4*9+3.5});
  \node[kgb lab, anchor=south east, text=bkVol0, font=\sffamily\bfseries\footnotesize] at (137.5,{4*9+3.7}) {目安の表で、数学Ⅱ（高2の例）に当たる★};
  \foreach \kgbc/\kgbv/\kgba/\kgbb/\kgbt [count=\kgbi from 0] in {
    2-7/2/1/2/{0乗・負\\の指数},3-11/3/2/3/{指数\\関数},3-12/3/2/3/対数,4-8/4/3/4/{三角\\関数},4-9/4/4/4/{加法\\定理},
    5-1/5/2/2/{平均の\\速さ},5-2/5/2/2/{極限の\\入口},5-3/5/3/3/{微分\\係数},5-9/5/2/2/{面積と\\道のり},6-8/6/4/4/複素数} {
    \pgfmathsetmacro\kgbx{26+\kgbi*10.8}
    \draw[bkVol\kgbv!85!black, line width=4pt, line cap=round] (\kgbx,{\kgba*9}) -- (\kgbx,{\kgbb*9+0.01});
    \fill[bkVol\kgbv, draw=bkInk, line width=0.5pt] (\kgbx,{\kgba*9}) circle (1.9) (\kgbx,{\kgbb*9}) circle (1.9);
    \node[fill=bkVol\kgbv, text=bkOn\kgbv, rounded corners=0.8mm, font=\sffamily\bfseries\footnotesize, inner xsep=0.8mm, inner ysep=0.5mm] at (\kgbx,0) {\kgbc};
}
\end{tikzpicture}
\bkfigcaption{数学Ⅱの内容を中心に扱う章と、その章の★（線の長さは★の範囲）}
\end{bkfigure}
\kgbnote{※ 図の章と、数学Ⅱに当たる内容：2-7 0乗・負の指数／3-11 指数関数／3-12 対数／4-8 三角関数／4-9 加法定理／5-1 平均の速さ（平均変化率）／5-2 極限の入口／5-3 微分係数・導関数／5-9 面積と道のり／6-8 複素数}

このように★が低くなっているのは、編集方針8にあるように、図・実験・表計算でたどれる書き方を選んでいるためです。★は「学校の何年生の内容か」ではなく、「このシリーズの中で、どれくらいの道具を使うか」と読んでください。学校の学年との対応は、4章の図と表、付録A・付録Bで確かめてください。

\kgbsub{9.2}{3つの段の章の★}

6-10〜6-14、7-7〜7-10、8-1〜8-13、9-1〜9-11 の各章（第6〜9巻の★★★★★の章は、読みものと終章を除いてすべて含まれます）は「目でわかる→手でわかる→式でわかる」の3段で書かれていて（1.5、図\ref{kgd-f-dan}）、「できるようになること」の3つの項目の★は、普通、段ごとに★★★・★★★★・★★★★★と並んでいます。章の帯に★★★★★とあっても、前の2段は★★★以下の道具（図・実験・表計算・プログラム）だけでたどれます。大学の内容の章を中学生・高校生が読むときは、「手でわかる」の終わりの「ここまでで十分」の旗まで読めれば、その章のゴールに着いたと考えてください。

\kgbsub{9.3}{学年の目安の読み方}

4章の「学年の目安」は、編集のときに決めた目安です。巻の中には、その学年より前や後に学校で習う内容も含まれています（図\ref{kgb-f-grid}）。学年の目安は「その学年になったら、この巻を読む」という決まりではありません。お子さんが学校でいま何を習っているかと、付録Aの対応表から、読み始める巻と章を選んでください。生徒用の紙面には学年を書いていないので、学年の目安より早く読んでも、遅く読んでも、子ども自身は気にせずに読めるようになっています。

\kgbsub{9.4}{照らし合わせるときの注意}

\kgbitem{{\sffamily\bfseries 1つの章が、いくつかの学年にまたがります。}★が範囲になっている章が34あります（例：2-2 ★〜★★、3-8 ★★〜★★★）。★が1つの章でも、1-2（小1〜小3）、2-6（小4〜中3）、3-13（小6〜数学Ⅰ）、4-2（小3〜数学C）のように、いくつかの学年にまたがる章が多くあります。章の中の「できるようになること」の各項目にも★が付いているので、項目ごとに照らすと正確です。}
\kgbitem{{\sffamily\bfseries 章の★より高い項目があります。}3-5（項目に★★★）、4-2（項目に★★★★：ラジアン）、5-8（項目に★★★）、6-2（項目に★★★）。}
\kgbitem{{\sffamily\bfseries 「もっと知りたい」は発展です。}章より★が高いことが多く、章の本体とは分けて数えています。}
\kgbitem{{\sffamily\bfseries 物理・工学の題材は、算数・数学の単元そのものではありません。}「物理・工学ではこう使う」のオームの法則（2-15）やコンデンサーの回路（6-4）などは理科の内容に当たります。理科の指導要領との対応は、この本では扱っていません。}

%% ====================================================================
\kgbsec{10}{授業・家庭での使い方の例}\label{kgb-ch10}

\kgbsub{10.1}{ご家庭で使うときに}

ノート・電卓・表計算・巻末の使い方は、第1章（1.8・1.9）にまとめました。ここでは、そばで見守る方へのお願いを挙げます。

\kgbitem{例題と「じぶんでとこう」で、章末問題に出てくる問題の形を、すべて1回ずつ見せています。章末問題で手が止まったら、同じ形の例題に戻るよう声をかけてください。}
\kgbitem{1回で分からなくても大丈夫です。章の「ふりかえり」の各項目に、戻るページが書いてあります。}
\kgbitem{学年の目安より早く読んでも遅く読んでも構いません（9.3）。生徒用の紙面には学年を書いていないので、「何年生の内容か」を本人に伝える必要はありません。}
\kgbitem{7章の内容（掛け算・割り算の筆算など）は、学校の教科書・問題集で補ってください（図\ref{kgb-f-hokai}）。}

\kgbsub{10.2}{学校で使うときに}

\kgbitem{章の流れはどの章も同じです（はてな？→やってみよう→数学の道具→例題→じぶんでとこう→まちがえやすいところ→物理・工学ではこう使う→章末問題）。授業の1時間を、この流れの一部に当てることができます。}
\kgbitem{「まちがえやすいところ」は、×と○を1行ずつ並べ、なぜ間違いかを書いています。授業でのつまずきの確認に使えます。}
\kgbitem{用語・書き方が学校の教科書と違う所（表\ref{kgb-t-yougo}）があります。特に虚数単位を $j$ と書く点（6-8 から）は、生徒に伝えてください。}

\par\Needspace{100mm}% 見出しと表10-1を同じページに
\kgbsub{10.3}{目的に合わせた使い方}

\par\addvspace{4pt}
\kgbcap{目的に合わせた使い方の例}{}
\begin{bktable}[long, theme=kgb]{hline{Z}={0pt,white}, vline{1-Z}={Z}{0pt,white}, row{Z}={abovesep=3pt, belowsep=1pt}, hspan=minimal, rowhead=1, colspec={Q[l,m,wd=32mm] Q[l,m,wd=94mm]}, rowsep=2pt, colsep=4pt, cells={font=\footnotesize}, row{1}={font=\sffamily\bfseries\footnotesize}}
目的 & 使い方の例 \\
学校の授業と並べて使う & 付録Aで、いま学校で習っている単元の章を探します。章の初めの「はてな？」と、終わりの「物理・工学ではこう使う」を先に読むと、その単元が何に使われるかが分かります \\
先へ進む & 巻の順に読むのが基本です。章の帯の下の「道しるべ」に、前に読む章とこの先の章が書いてあります。自分の読める★まで読んで、先へ進んでもかまいません \\
学び直す & 第1章の「どこから始めるか」（1.7）で、始める巻を選びます。分からない所は、道しるべの「この章の前に」の章へ戻ります \\
大学の工学に備える & 8章の表で、学ぶ科目の土台になる巻を選びます。3つの段の章は「ここまでで十分」の旗まで読み、その後「式でわかる」で専門書の書き方に慣れます \\
理科の授業とつなぐ & 「物理・工学ではこう使う」や「やってみよう」の実験は、理科の授業の導入や発展に使えます（例：2-15 のばねとオームの法則、6-4 のコンデンサーの放電） \\*
\SetCell[c=2]{c} \kgbcapline & \\
\end{bktable}
\par\addvspace{4pt}

\kgbsub{10.4}{付録の使い方}

この本の巻末には、2つの付録があります。付録Aは学習指導要領の側から、付録Bはシリーズの側から引く表です。お子さんが学校でいま習っている単元から章を探すときは付録Aを、シリーズの章が学校のどの学年の内容に当たるかを知りたいときは付録Bを使ってください。

\begin{bkfigure}
\begin{tikzpicture}[x=1mm, y=1mm]
  \node[kgb box, minimum width=40mm, minimum height=11mm] (s) at (20.5,0) {学習指導要領の項目\\{\mdseries（学年・科目・番号）}};
  \node[kgb box, minimum width=40mm, minimum height=11mm, fill=bkVol0!8] (c) at (119.5,0) {シリーズの章\\{\mdseries（全152章）}};
  \draw[kgb arr] ([yshift=2.2mm]s.east) -- node[kgb lab, above=0.8mm, font=\sffamily\bfseries\footnotesize] {付録A：項目から章を引く} ([yshift=2.2mm]c.west);
  \draw[kgb arr] ([yshift=-2.2mm]c.west) -- node[kgb lab, below=0.8mm, font=\sffamily\bfseries\footnotesize] {付録B：章から学年・科目を引く} ([yshift=-2.2mm]s.east);
\end{tikzpicture}
\bkfigcaption{2つの付録の引き方}
\end{bkfigure}
}{\typeout{BKPAGE missing guide-p2a}}
\IfFileExists{guide-p2b.tex}{%% ===== 自動生成：tools/mkguidetables.py（手で直さない。材料を直したら動かし直す） =====
%% ガイドブック（保護者・先生のための本）の付録（組 G-C、頭文字 kgc）。材料：tools/guide-taiou.md・series.py
\makeatletter
%% G-B の見出し（guide-p2a.tex の写し。p2a が先に読まれていれば、そちらが使われる）
%% \kgbsec[柱・しおり用の題]{番号}{題}（改ページしない）、\kgbsub{番号}{題}
\providecommand\kgbsec[3][]{%
  \par\addvspace{18pt}\Needspace{8\baselineskip}\global\bk@decotrue
  \setcounter{bkfig}{0}\setcounter{bktab}{0}\gdef\kgbsecno{#2}%
  \if\relax\detokenize{#1}\relax\def\kgb@st{#3}\else\def\kgb@st{#1}\fi
  \bksetrunning{}{保護者・先生のための本}{\if\relax\detokenize{#2}\relax\else#2\hspace{0.6\zw}\fi\kgb@st}%
  \pdfbookmark[1]{#2　\kgb@st}{kgbsec-#2-\thepage}%
  \kgc@secband{#2}{#3}}
\newcommand\kgc@secband[2]{%
  \noindent\begin{tikzpicture}
    \if\relax\detokenize{#1}\relax
      \node[anchor=west, text=white, font=\sffamily\bfseries\LARGE, align=left, inner sep=0pt] (t) at (0.6,0) {#2};
    \else
      \node[anchor=west, circle, fill=white, text=bkVol0, minimum size=1.35cm, font=\sffamily\bfseries\Large, inner sep=0pt] (n) at (0.35,0) {#1};
      \node[anchor=west, text=white, font=\sffamily\bfseries\LARGE, text width=\linewidth-2.4cm, align=left, inner sep=0pt] (t) at ([xshift=0.4cm]n.east) {#2};
    \fi
    \begin{pgfonlayer}{background}
      \path let \p1=(t.north), \p2=(t.south) in coordinate (kgbT) at (0,{max(\y1+0.4cm,1cm)}) coordinate (kgbB) at (0,{min(\y2-0.4cm,-1cm)});
      \fill[bkVol0] (kgbB) rectangle (\linewidth,0 |- kgbT);
    \end{pgfonlayer}
  \end{tikzpicture}\par\vspace{10pt}}
\providecommand\kgbsub[2]{%
  \par\addvspace{12pt}\Needspace{8\baselineskip}\noindent
  \begin{tikzpicture}[baseline=(t.base)]
    \node[anchor=west, fill=bkVol0, text=white, rounded corners=1mm, font=\sffamily\bfseries\small, inner xsep=1.6mm, inner ysep=1mm] (i) at (0,0) {#1};
    \node[anchor=west, font=\sffamily\bfseries\large, text=bkVol0, inner sep=0pt] (t) at ([xshift=0.6\zw]i.east) {#2};
    \draw[bkVol0, line width=0.8pt] ([yshift=-4pt]i.south west) -- ++(\linewidth-1pt,0);
  \end{tikzpicture}\par\nobreak\vspace{6pt}}
\providecommand\kgbsecno{}
%% ---- G-C の部品（kgc） ----
% 表の中の章番号（太いゴシック）
\newcommand\kgcch[1]{\mbox{\sffamily\bfseries #1}}
% 「もっと知りたい」（発展の囲み）の札：白い地に枠と字（色だけにたよらない）
\newcommand\kgcmotto{\mbox{\tikz[baseline=(t.base)]\node[draw=bkVol0, line width=0.5pt, rounded corners=0.8mm, inner xsep=0.5mm, inner ysep=0.35mm, font=\sffamily\footnotesize, text=bkVol0](t){もっと};}}
% 扱い方の記号の札（枠つき・太字）
\newcommand\kgcsym[1]{\mbox{\tikz[baseline=(t.base)]\node[draw=bkInk, line width=0.5pt, inner xsep=0.6mm, inner ysep=0.4mm, font=\sffamily\bfseries\footnotesize](t){#1};}}
\newcommand\kgcsep{\hspace{0.2em}}
\newcommand\kgcor{\hspace{0.15em}{\footnotesize ／}\hspace{0.15em}}
% ★：塗りの★＝その章の★の下から、白抜きの☆＝上まで（★★☆＝★★〜★★★）
\newcommand\kgcstars[2]{\mbox{\footnotesize\@tempcnta=0\relax
  \loop\advance\@tempcnta\@ne
    \ifnum\@tempcnta>#1\relax\textcolor{bkStar!70!black}{☆}\else\textcolor{bkStar}{★}\fi
  \ifnum\@tempcnta<#2\repeat}}
% 路線：アイコン（形の記号）＋丸数字
\newcommand\kgcroute[1]{\mbox{\bkrouteicon[0.95em]{#1}\,{\footnotesize\ifcase#1\or①\or②\or③\or④\or⑤\or⑥\or⑦\fi}}}
\newcommand\kgcallroutes{{\footnotesize ①〜⑦\newline すべて}}
% 範囲外の印（塗りの丸＝中心、半分塗りの丸＝一部。形で区別する）
\newcommand\kgcout[1]{\tikz[baseline=-0.35em]{\ifx#1a\fill[bkInk] (0,0) circle (0.42em);\else
  \fill[bkInk] (0,-0.42em) arc[start angle=-90, end angle=90, radius=0.42em] -- cycle;\fi
  \draw[bkInk, line width=0.6pt] (0,0) circle (0.42em);}}

% 凡例の囲み
\newenvironment{kgclegend}[1]{\par\addvspace{6pt}\begin{tcolorbox}[enhanced, breakable, arc=2mm, boxrule=0.8pt,
  colframe=bkVol0, colback=white, left=3mm, right=3mm, top=2.5mm, bottom=2.5mm,
  title={#1}, fonttitle=\sffamily\bfseries\small, coltitle=white, colbacktitle=bkVol0]\small}{\end{tcolorbox}\par\addvspace{4pt}}
% 表の見出しの行の色（白い地に濃い字の行より、うすい紺灰色の地に濃い字）
\colorlet{kgcHead}{bkVol0!14}
\colorlet{kgcBand}{bkVol0!6}
% longtblr の「表 1」などの見出しと続きの注を出さない
\DefTblrTemplate{caption}{kgc}{}
\DefTblrTemplate{capcont}{kgc}{\noindent{\sffamily\bfseries\small\color{bkVol0}\InsertTblrText{caption}}\hspace{0.6\zw}{\sffamily\footnotesize（前のページからの続き）}\par\vspace{2pt}}
\DefTblrTemplate{contfoot}{kgc}{\raggedleft\footnotesize（次のページに続きます）\par}
\NewTblrTheme{kgc}{\SetTblrTemplate{caption}{kgc}\SetTblrTemplate{capcont}{kgc}\SetTblrTemplate{contfoot}{kgc}}
\makeatother

\bkevenpage\kgbsec[付録A　学習指導要領の項目からシリーズの章へ]{A}{付録A　学習指導要領の項目から\\シリーズの章へ}\label{kgc-a}\label{kgc-sec10}
この付録は、現行の学習指導要領（小学校・中学校は平成29年告示、高等学校は平成30年告示）の算数・数学の項目ごとに、このシリーズのどの章で扱っているか、どのくらい扱っているかを示した一覧です。小学校第1学年から順に、中学校、高等学校の数学Ⅰ・数学A・数学Ⅱ・数学B・数学Ⅲ・数学Cの順に並べています。

\begin{itemize}[leftmargin=1.2\zw, itemsep=1pt, topsep=3pt]
\item 「番号」は、学年（科目）・領域・項目の番号です。例えば「小3-A(4)ア(ｴ)」は、小学校第3学年のA（数と計算）の(4)のアの(ｴ)を表します。「語」は〔用語・記号〕、「扱」は「内容の取扱い」の番号です。小学校のCは、第1〜3学年が「測定」、第4〜6学年が「変化と関係」です。
\item 「項目（要点）」は、指導要領の言葉を短くしたものです。正確な文は、指導要領の本文をご覧ください。
\item 「扱っている章」の \kgcch{1-4} などは章の番号（巻-章）です。「」の中は、その章の小見出しなどの題です。\kgcmotto{} は、章の中の発展の囲み「もっと知りたい」を表します。
\item 高等学校の科目を履修する学年は学校によって異なります。一般的な例は、数学Ⅰ・数学Aが1年、数学Ⅱ・数学Bが2年、数学Ⅲ・数学Cが3年です（指導要領で決められているものではありません）。
\item 「なし」「一部」の項目を補う方法は、「含めていない・軽く扱う内容と、補い方」の章をご覧ください。
\end{itemize}

\begin{kgclegend}{「扱い」の記号}
\begin{tblr}{colspec={Q[c,wd=10mm] X[l] Q[c,wd=10mm] X[l]}, rowsep=1.5pt, colsep=3pt, cells={font=\small}}
\kgcsym{同} & 指導要領とほぼ同じ範囲・深さで扱う & \kgcsym{後} & 後回し（その巻の「学年の目安」が、指導要領の学年より上。例：小3の内容を第4巻で） \\
\kgcsym{詳} & 詳しく（指導要領より深い・広い。証明を2通りにする、一般の形まで進む など） & \kgcsym{一部} & 一部だけ（語・考えの一部、「もっと知りたい」だけ、コラムだけ など） \\
\kgcsym{先} & 先取り（その巻の「学年の目安」が、指導要領の学年より下） & \kgcsym{なし} & 扱わない \\
\end{tblr}\par\smallskip
「先」「後」は、その巻の学年の目安と指導要領の学年を比べたもので、内容の深さ（同・詳・一部）とは別に書いています（例：\kgcsym{同}\kgcsep\kgcsym{後}）。「\kgcsym{一部}\kgcor\kgcsym{なし}」は、項目の中の事項によって扱いが分かれることを表します。\kgcsym{範囲外}は、指導要領の本文にない内容（数学Ⅲの注の微分方程式）を、このシリーズで扱っていることを表します。
\end{kgclegend}
\Needspace{259pt}\kgbsub{A.1}{小学校　第1学年}
\begin{longtblr}[theme=kgc, caption={小学校　第1学年}, entry=none, label=none]{colspec={Q[l,wd=19mm] Q[l,wd=31mm] X[2.4,l] Q[l,wd=15mm] X[1.6,l]}, width=\linewidth,
  rowhead=1, row{1}={bg=kgcHead, font=\sffamily\bfseries\footnotesize}, row{even[2]}={bg=kgcBand},
  hlines={0.4pt,bkLine}, vline{2-5}={0.4pt,bkLine}, rowsep=2pt, colsep=3pt, cells={font=\footnotesize}}
番号 & 項目（要点） & 扱っている章 & 扱い & 補足 \\
{\sffamily 小1-\allowbreak{}A(1)} & 数の構成と表し方 & \kgcch{1-1}（1対1対応・数える・順序と大小・「どれも3」）、\kgcch{1-2}（10のまとまり・一の位・十の位・1〜100の表・100・1000・10000）、\kgcch{1-6}（数直線）、\kgcch{1-4}（いくつといくつ＝ア(ｴ)） & \kgcsym{同} & \kgcch{1-2} は 1000・10000 まで進む（小2〜小3の範囲も含む） \\
{\sffamily 小1-\allowbreak{}A(2)} & 加法及び減法 & \kgcch{1-4}（あわせる・ふえる・10をこえる足し算・0の足し算）、\kgcch{1-5}（のこり・ちがい・かわった量・くり下がり）、\kgcch{1-7}（何十の計算） & \kgcsym{同} & \kgcch{1-5} は「かわった量」（変化量）を足している（路線①の始発） \\
{\sffamily 小1-\allowbreak{}B(1)} & 図形（形） & \kgcch{1-19}（三角形・四角形・円、つつ・はこを転がす） & \kgcsym{一部} & ア(ｳ)の方向や位置を表す言葉（前後・左右・上下）を教える節はなく、右・左は \kgcch{1-6} の数直線（右へ進む・左へ戻る）や \kgcch{1-2} の位（左の位）で使うだけ \\
{\sffamily 小1-\allowbreak{}C(1)} & 量の大きさの比較 & \kgcch{1-1}（ペアにして比べる）、\kgcch{1-8}（1 cm のいくつ分）、\kgcch{1-9}（てんびんで重さを比べる）、\kgcch{1-21}（マスの数で広さを比べる） & \kgcsym{同} &  \\
{\sffamily 小1-\allowbreak{}C(2)} & 時刻の読み方 & \kgcch{1-10} & \kgcsym{同} &  \\
{\sffamily 小1-\allowbreak{}D(1)} & 数量の整理（絵や図） & \kgcch{1-17}（数えて表に・ぼうグラフ） & \kgcsym{同} & 「絵グラフ」の語はない \\
{\sffamily 小1-\allowbreak{}語} & 一の位、十の位、＋、－、＝ & \kgcch{1-2}・\kgcch{1-4}・\kgcch{1-5} & \kgcsym{同} & 記号カード：＋・＝（\kgcch{1-4}）、－（\kgcch{1-5}） \\
\end{longtblr}
\Needspace{229pt}\kgbsub{A.2}{小学校　第2学年}
\begin{longtblr}[theme=kgc, caption={小学校　第2学年}, entry=none, label=none]{colspec={Q[l,wd=19mm] Q[l,wd=31mm] X[2.4,l] Q[l,wd=15mm] X[1.6,l]}, width=\linewidth,
  rowhead=1, row{1}={bg=kgcHead, font=\sffamily\bfseries\footnotesize}, row{even[2]}={bg=kgcBand},
  hlines={0.4pt,bkLine}, vline{2-5}={0.4pt,bkLine}, rowsep=2pt, colsep=3pt, cells={font=\footnotesize}}
番号 & 項目（要点） & 扱っている章 & 扱い & 補足 \\
{\sffamily 小2-\allowbreak{}A(1)} & 数の構成と表し方（4位数まで・十や百を単位・積としてみる・簡単な分数） & \kgcch{1-2}（1000・10000）、\kgcch{1-6}（1目もり10・100の数直線）、\kgcch{1-12}・\kgcch{1-13}（積としてみる）、\kgcch{1-16}（1/2・1/4） & \kgcsym{同} &  \\
{\sffamily 小2-\allowbreak{}A(2)} & 加法及び減法（2位数・筆算・性質・相互関係） & \kgcch{1-7}「ひっ算と くり上がり」「ひっ算と くり下がり」「こたえを たしかめよう」 & \kgcsym{同} &  \\
{\sffamily 小2-\allowbreak{}A(3)} & 乗法（意味・九九・性質・2位数×1位数） & \kgcch{1-12}（意味・アレイ・2〜5のだん）、\kgcch{1-13}（九九の表・ふえ方・入れかえても同じ・だんをたしてつくる・3×10） & \kgcsym{同}\newline{}（ア(ｵ)は一部） & ア(ｵ)「簡単な場合の2位数と1位数の乗法」は、\kgcch{1-13} の「3×10」、\kgcch{1-20} の「30×4」程度 \\
{\sffamily 小2-\allowbreak{}B(1)} & 三角形や四角形などの図形（箱の形を含む） & \kgcch{1-19}「三角形と 四角形」「直角」「直角の ある 形」 & \kgcsym{同}\newline{}（ア(ｳ)は一部） & ア(ｳ)「箱の形」は、\kgcch{1-19} のはてな？で「はこ」を転がす程度 \\
{\sffamily 小2-\allowbreak{}C(1)} & 長さ（mm・cm・m）、かさ（mL・dL・L） & \kgcch{1-8}、\kgcch{1-9} & \kgcsym{同} &  \\
{\sffamily 小2-\allowbreak{}C(2)} & 時間の単位（日・時・分） & \kgcch{1-10}「60分で 1時間」「午前と 午後」 & \kgcsym{同} &  \\
{\sffamily 小2-\allowbreak{}D(1)} & 簡単な表やグラフ & \kgcch{1-17} & \kgcsym{同} &  \\
{\sffamily 小2-\allowbreak{}語} & 直線、直角、頂点、辺、面、単位、×、＞、＜ & \kgcch{1-19}（辺・ちょう点・直角）、\kgcch{1-12}（×）、\kgcch{2-2}（＜・＞） & \kgcsym{一部}\kgcsep{}\kgcsym{後} & ＞・＜は第1巻にない（第2巻 \kgcch{2-2} で不等号） \\
{\sffamily 小2-\allowbreak{}扱} & (2)（ ）や□の式 (3)交換・結合法則 (4)乗数が1ずつ増える (5)敷き詰め & (3)(4)＝\kgcch{1-13}「入れかえても同じ」「ふえかたの きまり」、\kgcch{2-11}（結合・分配）。(2)＝\kgcch{2-11} & \kgcsym{一部} &  \\
\end{longtblr}
\Needspace{184pt}\kgbsub{A.3}{小学校　第3学年}
\begin{longtblr}[theme=kgc, caption={小学校　第3学年}, entry=none, label=none]{colspec={Q[l,wd=19mm] Q[l,wd=31mm] X[2.4,l] Q[l,wd=15mm] X[1.6,l]}, width=\linewidth,
  rowhead=1, row{1}={bg=kgcHead, font=\sffamily\bfseries\footnotesize}, row{even[2]}={bg=kgcBand},
  hlines={0.4pt,bkLine}, vline{2-5}={0.4pt,bkLine}, rowsep=2pt, colsep=3pt, cells={font=\footnotesize}}
番号 & 項目（要点） & 扱っている章 & 扱い & 補足 \\
{\sffamily 小3-\allowbreak{}A(1)} & 整数の表し方（万の単位、10倍・1/10） & \kgcch{1-2}「1000と 10000」、\kgcch{2-1}「10倍すると、小数点が右へ1つ」、\kgcch{2-7}（万・億・兆の表） & \kgcsym{一部}\kgcsep{}\kgcsym{後} & 「万の単位」の読み書きの節は第1巻にない。万→\nobreak{}億→\nobreak{}兆は \kgcch{2-7} \\
{\sffamily 小3-\allowbreak{}A(2)} & 加法及び減法（3位数・4位数） & \kgcch{1-7}「100を こえる たし算」 & \kgcsym{同} &  \\
{\sffamily 小3-\allowbreak{}A(3)} & 乗法（2位数・3位数×1位数・2位数、筆算） & \kgcch{1-13}（3×10）、\kgcch{1-20}（30×2〜30×6）。筆算の形は \kgcch{3-8} のコラム「かけ算の筆算も、展開だった」（23×14）だけ & \kgcsym{一部} & \textbf{掛け算の筆算の節はない} \\
{\sffamily 小3-\allowbreak{}A(4)} & 除法（意味・余り・乗法との関係・九九で求める・商が2位数） & \kgcch{1-14}（等分除・包含除・九九で見つける・1でわる・0では分けられない）、\kgcch{1-15}（余り・確かめの式） & \kgcsym{同}\newline{}（ア(ｵ)はなし） & 除法の例は九九の範囲（10÷2〜32÷8 など）。ア(ｵ)「商が2位数」はない \\
{\sffamily 小3-\allowbreak{}A(5)} & 小数（1/10 の位、加減） & \kgcch{1-16}「1を 10に 分けた 1つ分：0.1」、\kgcch{2-1}（1/10・1/100 の位） & \kgcsym{同} & 計算は \kgcch{2-4} \\
{\sffamily 小3-\allowbreak{}A(6)} & 分数（等分・単位分数・簡単な加減） & \kgcch{1-16}（1/2・1/4、0.5と1/2は同じ）、\kgcch{2-4}（足し算・引き算） & \kgcsym{一部}\kgcsep{}\kgcsym{後} & 第1巻は 1/2・1/4 まで \\
{\sffamily 小3-\allowbreak{}A(7)} & 数量の関係を表す式（□） & \kgcch{2-11}「ことばの式を、□で短く書く」 & \kgcsym{同}\kgcsep{}\kgcsym{後} &  \\
{\sffamily 小3-\allowbreak{}A(8)} & そろばん & なし & \kgcsym{なし} & どの巻にも語がない \\
{\sffamily 小3-\allowbreak{}B(1)} & 二等辺三角形・正三角形、作図、角、円・球 & 円・中心・半径・直径・コンパス＝\kgcch{1-19}「円と タイヤ」。角＝\kgcch{1-20}。二等辺三角形・正三角形＝\kgcch{4-1}（対称で性質を説明）。球＝\kgcch{4-5} & \kgcsym{一部}\kgcsep{}\kgcsym{後} & 定規とコンパスで二等辺三角形をかく作図はない \\
{\sffamily 小3-\allowbreak{}C(1)} & 長さ（km）、重さ（g・kg）、測定 & \kgcch{1-9}「おもさの たんい：g と kg」、\kgcch{1-11}（目もり・はかり）、km は \kgcch{2-10}（接頭語 k） & \kgcsym{一部}\kgcsep{}\kgcsym{後} & 第1巻に km はない \\
{\sffamily 小3-\allowbreak{}C(2)} & 時刻と時間（秒、時間を求める） & \kgcch{1-10}「時こくと 時間」「\kgcmotto{} みじかい 時間：秒」 & \kgcsym{同} &  \\
{\sffamily 小3-\allowbreak{}D(1)} & 表と棒グラフ & \kgcch{1-17}「1目もりの 大きさ」「1目もりの きめかた」 & \kgcsym{同} &  \\
{\sffamily 小3-\allowbreak{}語} & 等号、不等号、小数点、1/10の位、数直線、分母、分子、÷ & \kgcch{1-4}（＝）、\kgcch{1-6}（数直線）、\kgcch{1-14}（÷）、\kgcch{1-16}（小数点）、\kgcch{2-1}（1/10の位）、\kgcch{2-2}（不等号）、\kgcch{2-4}（分母・分子） & \kgcsym{同}\newline{}（不等号は後） &  \\
{\sffamily 小3-\allowbreak{}扱} & (7)トン（t）、接頭語（k・m） & \kgcch{2-10}「大きさを表す記号：k・m・$\symup{\mu}$」 & \kgcsym{詳}\kgcsep{}\kgcsym{後} &  \\
\end{longtblr}
\Needspace{154pt}\kgbsub{A.4}{小学校　第4学年}
\begin{longtblr}[theme=kgc, caption={小学校　第4学年}, entry=none, label=none]{colspec={Q[l,wd=19mm] Q[l,wd=31mm] X[2.4,l] Q[l,wd=15mm] X[1.6,l]}, width=\linewidth,
  rowhead=1, row{1}={bg=kgcHead, font=\sffamily\bfseries\footnotesize}, row{even[2]}={bg=kgcBand},
  hlines={0.4pt,bkLine}, vline{2-5}={0.4pt,bkLine}, rowsep=2pt, colsep=3pt, cells={font=\footnotesize}}
番号 & 項目（要点） & 扱っている章 & 扱い & 補足 \\
{\sffamily 小4-\allowbreak{}A(1)} & 整数の表し方（億・兆） & \kgcch{2-7}「大きな数は、10の累乗で」 & \kgcsym{同} &  \\
{\sffamily 小4-\allowbreak{}A(2)} & 概数（四捨五入・見積り） & \kgcch{2-6}「がい数と四捨五入」 & \kgcsym{同} & 「概数」の語はなく「がい数」 \\
{\sffamily 小4-\allowbreak{}A(3)} & 整数の除法（筆算、被除数＝除数×商＋余り） & ア(ｳ)＝\kgcch{1-15}「たしかめの しき」 & \kgcsym{一部} & \textbf{割り算の筆算の節はない} \\
{\sffamily 小4-\allowbreak{}A(4)} & 小数の仕組みと計算 & \kgcch{2-1}、\kgcch{2-4} & \kgcsym{同} &  \\
{\sffamily 小4-\allowbreak{}A(5)} & 分数（等しい分数・同分母の加減） & \kgcch{2-4}「大きさの同じ分数：約分と通分」 & \kgcsym{同} & 真分数・仮分数・帯分数の語も \kgcch{2-4} \\
{\sffamily 小4-\allowbreak{}A(6)} & 数量の関係を表す式（四則混合・（ ）・公式・□△） & \kgcch{2-11}、\kgcch{2-13} & \kgcsym{同} &  \\
{\sffamily 小4-\allowbreak{}A(7)} & 計算に関して成り立つ性質 & \kgcch{2-11}「文字で、いつでも成り立つきまりを書く」、\kgcch{1-13} & \kgcsym{同} &  \\
{\sffamily 小4-\allowbreak{}A(8)} & そろばん & なし & \kgcsym{なし} &  \\
{\sffamily 小4-\allowbreak{}B(1)} & 平面図形（平行・垂直、平行四辺形・ひし形・台形） & \kgcch{4-1}（⊥・∥、平行線と角）、\kgcch{4-5}（平行四辺形・台形・ひし形の面積） & \kgcsym{同}\kgcsep{}\kgcsym{後} &  \\
{\sffamily 小4-\allowbreak{}B(2)} & 立体図形（立方体・直方体、見取図・展開図） & \kgcch{2-7}（見取図）、\kgcch{4-5}「展開図と表面積・回転体」、\kgcch{4-13}（投影図・空間の平行垂直） & \kgcsym{同}\kgcsep{}\kgcsym{後} &  \\
{\sffamily 小4-\allowbreak{}B(3)} & ものの位置の表し方 & \kgcch{3-2}（座標）、\kgcch{3-2} \kgcmotto{}「部屋の中を飛ぶハエは、3つの数」、\kgcch{4-13} & \kgcsym{詳}\kgcsep{}\kgcsym{後} &  \\
{\sffamily 小4-\allowbreak{}B(4)} & 平面図形の面積（cm$^{2}$・m$^{2}$・km$^{2}$、正方形・長方形） & \kgcch{1-21}「1cm$^{2}$」「長方形の 面積は かけ算で」「大きな 広さ：1m$^{2}$」、\kgcch{2-10}「面積・体積の換算」 & \kgcsym{同}\kgcsep{}\kgcsym{先} & 小4-扱(8) アール・ヘクタールはない \\
{\sffamily 小4-\allowbreak{}B(5)} & 角の大きさ（度・測定） & \kgcch{1-20} & \kgcsym{同}\kgcsep{}\kgcsym{先} &  \\
{\sffamily 小4-\allowbreak{}C(1)} & 伴って変わる二つの数量（折れ線グラフ） & \kgcch{1-18} & \kgcsym{同}\kgcsep{}\kgcsym{先} & 「かたむきが 急なほど、はやく かわる」は傾き（中2）の種 \\
{\sffamily 小4-\allowbreak{}C(2)} & 簡単な場合についての割合 & \kgcch{2-5}「割合：もとにする量を1とみる」 & \kgcsym{同} &  \\
{\sffamily 小4-\allowbreak{}D(1)} & データの分類整理（二つの観点・折れ線グラフ） & ア(ｲ)＝\kgcch{1-18} & \kgcsym{一部} &  \\
{\sffamily 小4-\allowbreak{}語} & 和・差・積・商、以上・以下・未満、真分数・仮分数・帯分数、平行・垂直・対角線・平面 & \kgcch{2-2}（以上・以下・未満）、\kgcch{2-4}（真分数など）、\kgcch{4-1}（平行・垂直）、\kgcch{4-13}（平面） & \kgcsym{同} &  \\
\end{longtblr}
\Needspace{259pt}\kgbsub{A.5}{小学校　第5学年}
\begin{longtblr}[theme=kgc, caption={小学校　第5学年}, entry=none, label=none]{colspec={Q[l,wd=19mm] Q[l,wd=31mm] X[2.4,l] Q[l,wd=15mm] X[1.6,l]}, width=\linewidth,
  rowhead=1, row{1}={bg=kgcHead, font=\sffamily\bfseries\footnotesize}, row{even[2]}={bg=kgcBand},
  hlines={0.4pt,bkLine}, vline{2-5}={0.4pt,bkLine}, rowsep=2pt, colsep=3pt, cells={font=\footnotesize}}
番号 & 項目（要点） & 扱っている章 & 扱い & 補足 \\
{\sffamily 小5-\allowbreak{}A(1)} & 整数の性質（偶数・奇数、約数・倍数） & \kgcch{2-4}（公約数・公倍数・素数）、\kgcch{2-9} \kgcmotto{}「$\sqrt{2}$ が分数で書けないわけ」（偶数・奇数の定義） & \kgcsym{一部} & 偶数・奇数を定義するのは \kgcch{2-9} の「もっと知りたい」（\mbox{★★★}）だけ。「最大公約数」「最小公倍数」の語はない（「公倍数でいちばん小さい数」と言う） \\
{\sffamily 小5-\allowbreak{}A(2)} & 整数及び小数の表し方（小数点の移動） & \kgcch{2-1} & \kgcsym{同} &  \\
{\sffamily 小5-\allowbreak{}A(3)} & 小数の乗法及び除法 & \kgcch{2-4}「0と1のあいだの数をかける・わる」 & \kgcsym{同} &  \\
{\sffamily 小5-\allowbreak{}A(4)} & 分数の意味と表し方 & \kgcch{2-4}「分数は、わり算の答え」 & \kgcsym{同} &  \\
{\sffamily 小5-\allowbreak{}A(5)} & 分数の加法及び減法（異分母） & \kgcch{2-4}「たし算・ひき算は、通分してから」 & \kgcsym{同} &  \\
{\sffamily 小5-\allowbreak{}A(6)} & 数量の関係を表す式 & \kgcch{2-11}、\kgcch{2-15} & \kgcsym{同} &  \\
{\sffamily 小5-\allowbreak{}B(1)} & 平面図形（合同・多角形・正多角形・円周率） & \kgcch{4-1}（合同と対称・多角形の内角と外角・正多角形）、\kgcch{2-9}「円周率 $\pi$」 & \kgcsym{同}\kgcsep{}\kgcsym{後}\newline{}（$\pi$ は同） & プログラミングで正多角形をかく例（小5-プログラミング）はない \\
{\sffamily 小5-\allowbreak{}B(2)} & 立体図形（角柱・円柱） & \kgcch{4-5}「柱の体積：底面積×高さ」 & \kgcsym{同}\kgcsep{}\kgcsym{後} &  \\
{\sffamily 小5-\allowbreak{}B(3)} & 平面図形の面積（三角形・平行四辺形・ひし形・台形） & \kgcch{4-5} & \kgcsym{同}\kgcsep{}\kgcsym{後} &  \\
{\sffamily 小5-\allowbreak{}B(4)} & 立体図形の体積（cm$^{3}$・m$^{3}$、直方体・立方体） & \kgcch{2-7}「面積は2乗、体積は3乗」、\kgcch{2-10}、\kgcch{4-5} & \kgcsym{同} &  \\
{\sffamily 小5-\allowbreak{}C(1)} & 比例（簡単な場合） & \kgcch{2-15} & \kgcsym{同} &  \\
{\sffamily 小5-\allowbreak{}C(2)} & 速さなど単位量当たりの大きさ & \kgcch{2-3}「わり算で『1あたりの量』」、\kgcch{2-10}「時速を秒速に直す」 & \kgcsym{同} &  \\
{\sffamily 小5-\allowbreak{}C(3)} & 割合・百分率 & \kgcch{2-5}「百分率：もとを100とみる」、\kgcch{2-5} で歩合 & \kgcsym{同} &  \\
{\sffamily 小5-\allowbreak{}D(1)} & 円グラフ・帯グラフ、統計的な問題解決 & なし & \kgcsym{なし} & 円グラフ・帯グラフはどの巻にもない。統計的な問題解決の手順（問い→\nobreak{}計画→\nobreak{}データ→\nobreak{}分析→\nobreak{}結論）としての節はない。データの集め方・表・グラフは \kgcch{1-17}・\kgcch{1-18}・\kgcch{3-14}〜\kgcch{3-16} などで扱う \\
{\sffamily 小5-\allowbreak{}D(2)} & 平均 & \kgcch{2-6}「平均：でこぼこをならして1つの数にする」、\kgcch{3-13} & \kgcsym{同} &  \\
{\sffamily 小5-\allowbreak{}語} & 最大公約数、最小公倍数、通分、約分、底面、側面、比例、％ & \kgcch{2-4}（通分・約分）、\kgcch{2-5}（％）、\kgcch{2-15}（比例）、\kgcch{4-5}（底面） & \kgcsym{一部} & 最大公約数・最小公倍数の語はない \\
\end{longtblr}
\Needspace{184pt}\kgbsub{A.6}{小学校　第6学年}
\begin{longtblr}[theme=kgc, caption={小学校　第6学年}, entry=none, label=none]{colspec={Q[l,wd=19mm] Q[l,wd=31mm] X[2.4,l] Q[l,wd=15mm] X[1.6,l]}, width=\linewidth,
  rowhead=1, row{1}={bg=kgcHead, font=\sffamily\bfseries\footnotesize}, row{even[2]}={bg=kgcBand},
  hlines={0.4pt,bkLine}, vline{2-5}={0.4pt,bkLine}, rowsep=2pt, colsep=3pt, cells={font=\footnotesize}}
番号 & 項目（要点） & 扱っている章 & 扱い & 補足 \\
{\sffamily 小6-\allowbreak{}A(1)} & 分数の乗法及び除法 & \kgcch{2-4}「かけ算は、面積図で」「わり算（その2）：なぜ、上と下を入れかえた数をかけるの？」 & \kgcsym{同} &  \\
{\sffamily 小6-\allowbreak{}A(2)} & 文字を用いた式 & \kgcch{2-11}「□のかわりに、文字を使う」 & \kgcsym{同} &  \\
{\sffamily 小6-\allowbreak{}B(1)} & 縮図・拡大図、対称な図形 & \kgcch{4-3}「拡大図・縮図と縮尺」、\kgcch{4-1}「ずらす・回す・裏返す：合同と対称」 & \kgcsym{同}\kgcsep{}\kgcsym{後} &  \\
{\sffamily 小6-\allowbreak{}B(2)} & 概形とおよその面積 & \kgcch{4-5}「方眼で数えて、細かくするほど正確に」、\kgcch{5-9} \kgcmotto{}「グラフが曲がっていても、マスを数えればおよそわかる」 & \kgcsym{同}\kgcsep{}\kgcsym{後} &  \\
{\sffamily 小6-\allowbreak{}B(3)} & 円の面積 & \kgcch{4-5}「円の面積：細いおうぎ形に切って並べかえる」 & \kgcsym{同}\kgcsep{}\kgcsym{後} & \kgcch{4-2} は「\kgcch{4-5} でくわしく」と先に使う \\
{\sffamily 小6-\allowbreak{}B(4)} & 角柱及び円柱の体積 & \kgcch{4-5} & \kgcsym{同}\kgcsep{}\kgcsym{後} &  \\
{\sffamily 小6-\allowbreak{}C(1)} & 比例（反比例は「知る」程度） & \kgcch{2-15}「かけると一定：反比例」、\kgcch{3-7}（反比例のグラフ） & \kgcsym{詳} & 反比例を xy 一定・グラフ・pV 一定まで扱う \\
{\sffamily 小6-\allowbreak{}C(2)} & 比 & \kgcch{2-5}「比と比の値」「比はかけ算でつなげる」 & \kgcsym{同} &  \\
{\sffamily 小6-\allowbreak{}D(1)} & データの考察（代表値・度数分布） & \kgcch{3-13}「1つの数で代表させる：平均・中央値・最頻値」「数を区切って数える：度数分布表」、\kgcch{2-6}（点を並べた図） & \kgcsym{同}\kgcsep{}\kgcsym{後} & \textbf{「ドットプロット」の語は紙面にない}（\kgcch{2-6} は点を並べた図で、名前は付けていない） \\
{\sffamily 小6-\allowbreak{}D(2)} & 起こり得る場合 & \kgcch{3-17}「場合を数えて、確率を求める」（表・樹形図） & \kgcsym{同}\kgcsep{}\kgcsym{後} &  \\
{\sffamily 小6-\allowbreak{}語} & 線対称・点対称・対称の軸・対称の中心、比の値、ドットプロット、平均値・中央値・最頻値・階級、： & \kgcch{4-1}（線対称・点対称）、\kgcch{2-5}（比の値・：）、\kgcch{3-13}（中央値・最頻値・階級） & \kgcsym{一部} & ドットプロットの語なし（上） \\
\end{longtblr}
\Needspace{169pt}\kgbsub{A.7}{中学校　第1学年}
\begin{longtblr}[theme=kgc, caption={中学校　第1学年}, entry=none, label=none]{colspec={Q[l,wd=19mm] Q[l,wd=31mm] X[2.4,l] Q[l,wd=15mm] X[1.6,l]}, width=\linewidth,
  rowhead=1, row{1}={bg=kgcHead, font=\sffamily\bfseries\footnotesize}, row{even[2]}={bg=kgcBand},
  hlines={0.4pt,bkLine}, vline{2-5}={0.4pt,bkLine}, rowsep=2pt, colsep=3pt, cells={font=\footnotesize}}
番号 & 項目（要点） & 扱っている章 & 扱い & 補足 \\
{\sffamily 中1-\allowbreak{}A(1)} & 正の数・負の数（扱(1)素数の積） & \kgcch{2-2}（正負の数の四則）、\kgcch{2-4}（素数の掛け算に分ける） & \kgcsym{同} & 「素因数分解」の語はない \\
{\sffamily 中1-\allowbreak{}A(2)} & 文字を用いた式 & \kgcch{2-11}、\kgcch{2-12}、\kgcch{2-13} & \kgcsym{同} &  \\
{\sffamily 中1-\allowbreak{}A(3)} & 一元一次方程式（扱(4)比例式） & \kgcch{2-14}「等式は天びん」「移項」、\kgcch{2-16}「1次方程式の解き方」、比例式＝\kgcch{3-3}「いろいろな関数」の中の段 & \kgcsym{同} &  \\
{\sffamily 中1-\allowbreak{}B(1)} & 平面図形（基本的な作図、移動） & 移動＝\kgcch{4-1}「ずらす・回す・裏返す」、\kgcch{4-10}「平行移動」。作図＝\kgcch{4-2}「垂直二等分線と、円の中心」（紙を折って見つける） & \kgcsym{一部}\kgcsep{}\kgcsym{後} & 角の二等分線・垂線の作図はない \\
{\sffamily 中1-\allowbreak{}B(2)} & 空間図形（位置関係、扇形、表面積と体積） & \kgcch{4-13}（ねじれの位置・投影図）、\kgcch{4-2}「おうぎ形の弧と面積は、中心角に比例する」、\kgcch{4-5}「すいの体積」「球の体積と表面積」 & \kgcsym{同}\kgcsep{}\kgcsym{後} &  \\
{\sffamily 中1-\allowbreak{}C(1)} & 比例・反比例（関数・座標・表式グラフ） & \kgcch{2-15}、\kgcch{3-1}（変数・変域）、\kgcch{3-2}（座標）、\kgcch{3-3}（関数）、\kgcch{3-7} & \kgcsym{同}\kgcsep{}\kgcsym{詳} &  \\
{\sffamily 中1-\allowbreak{}D(1)} & データの分布（ヒストグラム・相対度数） & \kgcch{3-13}「柱のグラフにする：ヒストグラム」 & \kgcsym{同}\kgcsep{}\kgcsym{後} & 累積度数も \kgcch{3-13} \\
{\sffamily 中1-\allowbreak{}D(2)} & 多数回の試行による確率 & \kgcch{3-17}「たくさんくり返すと、割合が落ち着く」 & \kgcsym{同}\kgcsep{}\kgcsym{後} & 表計算の RANDBETWEEN で試す \\
{\sffamily 中1-\allowbreak{}語} & 自然数・素数・符号・絶対値・項・係数・移項・≦≧、弧・弦・回転体・ねじれの位置・$\pi$・∥・⊥・∠・△、関数・変数・変域、範囲・累積度数 & \kgcch{2-2}・\kgcch{2-4}・\kgcch{2-9}・\kgcch{2-12}・\kgcch{2-14}・\kgcch{3-1}・\kgcch{3-2}・\kgcch{3-13}・\kgcch{4-1}・\kgcch{4-2}・\kgcch{4-5}・\kgcch{4-13} & \kgcsym{同} &  \\
{\sffamily 中1-\allowbreak{}扱} & (5)円の接線は半径に垂直 (6)投影図 & \kgcch{4-2}「円の接線は、半径に垂直」、\kgcch{4-13}「空間の形を紙にかく：投影図」 & \kgcsym{同}\kgcsep{}\kgcsym{後} &  \\
\end{longtblr}
\Needspace{184pt}\kgbsub{A.8}{中学校　第2学年}
\begin{longtblr}[theme=kgc, caption={中学校　第2学年}, entry=none, label=none]{colspec={Q[l,wd=19mm] Q[l,wd=31mm] X[2.4,l] Q[l,wd=15mm] X[1.6,l]}, width=\linewidth,
  rowhead=1, row{1}={bg=kgcHead, font=\sffamily\bfseries\footnotesize}, row{even[2]}={bg=kgcBand},
  hlines={0.4pt,bkLine}, vline{2-5}={0.4pt,bkLine}, rowsep=2pt, colsep=3pt, cells={font=\footnotesize}}
番号 & 項目（要点） & 扱っている章 & 扱い & 補足 \\
{\sffamily 中2-\allowbreak{}A(1)} & 文字を用いた式（整式の加減・単項式の乗除・式の変形） & \kgcch{2-12}、\kgcch{2-14}「x＝vt から、時間を求める式を作る」、\kgcch{3-8} & \kgcsym{同}\kgcsep{}\kgcsym{先} &  \\
{\sffamily 中2-\allowbreak{}A(2)} & 連立二元一次方程式 & \kgcch{3-6} & \kgcsym{同} & \kgcmotto{}「3つの文字の連立方程式」（\mbox{★★★}） \\
{\sffamily 中2-\allowbreak{}B(1)} & 平行線や角、多角形の角 & \kgcch{4-1}「平行線と角：対頂角・同位角・錯角」 & \kgcsym{同} &  \\
{\sffamily 中2-\allowbreak{}B(2)} & 図形の合同（合同条件・証明） & \kgcch{4-1}「三角形が1つに決まる条件：合同条件」「三角形の内角の和は180°」 & \kgcsym{一部} & 証明は「内角の和」の1題。二等辺三角形・平行四辺形の性質は対称で説明し、証明しない \\
{\sffamily 中2-\allowbreak{}C(1)} & 一次関数（二元一次方程式を関数とみる） & \kgcch{3-5}、\kgcch{3-6}「交わる点は、2つの式を同時にみたす」 & \kgcsym{同} &  \\
{\sffamily 中2-\allowbreak{}D(1)} & 四分位範囲・箱ひげ図 & \kgcch{3-13}「4つに分けて見る：四分位数と箱ひげ図」 & \kgcsym{同} &  \\
{\sffamily 中2-\allowbreak{}D(2)} & 場合の数を基にして得られる確率 & \kgcch{3-17} & \kgcsym{同} &  \\
{\sffamily 中2-\allowbreak{}語} & 同類項、対頂角・内角・外角・定義・証明・逆・反例・≡、変化の割合・傾き & \kgcch{2-12}、\kgcch{4-1}「逆と反例」、\kgcch{4-3}（≡ の記号カード）、\kgcch{3-5} & \kgcsym{同} &  \\
{\sffamily 中2-\allowbreak{}扱} & (1)正方形・ひし形・長方形は平行四辺形の特別な形 & \kgcch{4-1}（「ひし形は、正方形ではない」：逆と反例）、\kgcch{4-5}「台形・ひし形・多角形」（ひし形の定義と面積） & \kgcsym{一部} &  \\
\end{longtblr}
\Needspace{184pt}\kgbsub{A.9}{中学校　第3学年}
\begin{longtblr}[theme=kgc, caption={中学校　第3学年}, entry=none, label=none]{colspec={Q[l,wd=19mm] Q[l,wd=31mm] X[2.4,l] Q[l,wd=15mm] X[1.6,l]}, width=\linewidth,
  rowhead=1, row{1}={bg=kgcHead, font=\sffamily\bfseries\footnotesize}, row{even[2]}={bg=kgcBand},
  hlines={0.4pt,bkLine}, vline{2-5}={0.4pt,bkLine}, rowsep=2pt, colsep=3pt, cells={font=\footnotesize}}
番号 & 項目（要点） & 扱っている章 & 扱い & 補足 \\
{\sffamily 中3-\allowbreak{}A(1)} & 平方根 & \kgcch{2-8}（√・√どうしの掛け算）、\kgcch{2-9}（無理数）、\kgcch{3-10}「√ のついた数の計算」（有理化） & \kgcsym{同}\kgcsep{}\kgcsym{先} & \kgcch{2-8}・\kgcch{2-9} は\mbox{★★}の章（第2巻） \\
{\sffamily 中3-\allowbreak{}A(2)} & 式の展開と因数分解 & \kgcch{3-8}、\kgcch{3-10}「かけ算の形に直す：因数分解」 & \kgcsym{同} &  \\
{\sffamily 中3-\allowbreak{}A(3)} & 二次方程式（解の公式） & \kgcch{3-10} & \kgcsym{詳} & 判別式・2次不等式まで（数Ⅰ・数Ⅱの内容） \\
{\sffamily 中3-\allowbreak{}B(1)} & 図形の相似（相似条件・面積比・体積比、平行線と線分の比） & \kgcch{4-3} & \kgcsym{同} & 相似の証明の場面はない \\
{\sffamily 中3-\allowbreak{}B(2)} & 円周角と中心角 & \kgcch{4-2}「円周角は、中心角の半分」 & \kgcsym{同} &  \\
{\sffamily 中3-\allowbreak{}B(3)} & 三平方の定理 & \kgcch{4-4}「どんな直角三角形でも成り立つわけ：2通りの証明」 & \kgcsym{詳} & 指導要領は「証明できることを知る」。シリーズは2通りの証明を本文で見せる \\
{\sffamily 中3-\allowbreak{}C(1)} & 関数 y＝ax$^{2}$ & \kgcch{3-9} & \kgcsym{同} &  \\
{\sffamily 中3-\allowbreak{}D(1)} & 標本調査 & \kgcch{3-17}「全部を調べずに、一部から全体を知る：標本調査」（母集団・標本・無作為・全数調査） & \kgcsym{同} &  \\
{\sffamily 中3-\allowbreak{}語} & 根号・有理数・無理数・因数・√、∽、全数調査 & \kgcch{2-8}・\kgcch{2-9}・\kgcch{3-10}・\kgcch{4-3}・\kgcch{3-17} & \kgcsym{同} &  \\
{\sffamily 中3-\allowbreak{}扱} & (1)誤差・近似値・$a\times10^{n}$ (2)実数解をもつ二次方程式 (4)円周角の定理の逆 & (1)＝\kgcch{2-6}（誤差・有効数字）、\kgcch{2-7}（10の累乗） & \kgcsym{同}\kgcsep{}\kgcsym{先} & (2)：シリーズは \kgcch{6-8} で虚数解まで進む \\
\end{longtblr}
\Needspace{169pt}\kgbsub{A.10}{高等学校　数学Ⅰ}
\begin{longtblr}[theme=kgc, caption={高等学校　数学Ⅰ}, entry=none, label=none]{colspec={Q[l,wd=19mm] Q[l,wd=31mm] X[2.4,l] Q[l,wd=15mm] X[1.6,l]}, width=\linewidth,
  rowhead=1, row{1}={bg=kgcHead, font=\sffamily\bfseries\footnotesize}, row{even[2]}={bg=kgcBand},
  hlines={0.4pt,bkLine}, vline{2-5}={0.4pt,bkLine}, rowsep=2pt, colsep=3pt, cells={font=\footnotesize}}
番号 & 項目（要点） & 扱っている章 & 扱い & 補足 \\
{\sffamily 数Ⅰ-\allowbreak{}(1)ア(ｱ)} & 実数、無理数の四則 & \kgcch{2-9}、\kgcch{3-10} & \kgcsym{同}\kgcsep{}\kgcsym{先} &  \\
{\sffamily 数Ⅰ-\allowbreak{}(1)ア(ｲ)} & 集合と命題 & \kgcch{4-1}「逆と反例」（仮定・結論） & \kgcsym{一部} & \textbf{集合はなし}（「集合」の語は \kgcch{4-0} の「集合場所」だけ）。「命題」の語も \kgcch{4-1} にない。「必要十分条件」の語は \kgcch{7-5} の専門書の囲みに出る \\
{\sffamily 数Ⅰ-\allowbreak{}(1)ア(ｳ)} & 二次の乗法公式・因数分解 & \kgcch{3-8}、\kgcch{3-10} & \kgcsym{同}\kgcsep{}\kgcsym{先} &  \\
{\sffamily 数Ⅰ-\allowbreak{}(1)ア(ｴ)} & 一次不等式 & \kgcch{3-5}「どこから上になる？：1次不等式の入り口」、\kgcch{2-2}（不等号） & \kgcsym{一部}\kgcsep{}\kgcsym{先} & 「入り口」の扱い \\
{\sffamily 数Ⅰ-\allowbreak{}(2)ア(ｱ)(ｲ)} & 三角比（鋭角・鈍角） & \kgcch{4-6}、\kgcch{4-7}「三角形の名前のつけ方と、鈍角の三角比」 & \kgcsym{同} & 数Ⅰ-扱(3) 0°・90°・180° は \kgcch{4-6}「0°から180°までの三角比」 \\
{\sffamily 数Ⅰ-\allowbreak{}(2)ア(ｳ)} & 正弦定理・余弦定理 & \kgcch{4-7} & \kgcsym{詳} & ロボットの腕の順運動学・逆運動学、力の三角形に使う \\
{\sffamily 数Ⅰ-\allowbreak{}(3)ア(ｱ)(ｲ)} & 二次関数のグラフ・最大最小 & \kgcch{3-9} & \kgcsym{同}\kgcsep{}\kgcsym{先} &  \\
{\sffamily 数Ⅰ-\allowbreak{}(3)ア(ｳ)} & 二次方程式・二次不等式とグラフ & \kgcch{3-10}「式がプラスになるのは、どこ？」 & \kgcsym{同}\kgcsep{}\kgcsym{先} & 本文に「高校では『2次不等式』として学びます」とある \\
{\sffamily 数Ⅰ-\allowbreak{}(4)ア(ｱ)} & 分散・標準偏差・散布図・相関係数 & \kgcch{3-14}、\kgcch{3-15} & \kgcsym{詳}\kgcsep{}\kgcsym{先} & 共分散、$\sigma/\sqrt{n}$（\kgcch{3-14}）、相関係数＝偏差ベクトルの向き（\kgcch{4-12}） \\
{\sffamily 数Ⅰ-\allowbreak{}(4)ア(ｲ)} & コンピュータでの整理 & \kgcch{2-18}（表計算）、\kgcch{3-14}・\kgcch{3-15} & \kgcsym{同} &  \\
{\sffamily 数Ⅰ-\allowbreak{}(4)ア(ｳ)} & 仮説検定の考え方 & \kgcch{3-17} \kgcmotto{}「このコインは、いかさま？（仮説検定）」（\mbox{★★★}） & \kgcsym{一部} & 「もっと知りたい」だけ \\
{\sffamily 数Ⅰ-\allowbreak{}(4)語} & 外れ値 & \kgcch{3-13}、\kgcch{3-15} & \kgcsym{同} &  \\
{\sffamily 数Ⅰ-\allowbreak{}扱} & (2)有限小数・循環小数 & \kgcch{2-9}「分数を小数にすると、終わるかくり返す」 & \kgcsym{同}\kgcsep{}\kgcsym{先} &  \\
\end{longtblr}
\Needspace{199pt}\kgbsub{A.11}{高等学校　数学A}
\begin{longtblr}[theme=kgc, caption={高等学校　数学A}, entry=none, label=none]{colspec={Q[l,wd=19mm] Q[l,wd=31mm] X[2.4,l] Q[l,wd=15mm] X[1.6,l]}, width=\linewidth,
  rowhead=1, row{1}={bg=kgcHead, font=\sffamily\bfseries\footnotesize}, row{even[2]}={bg=kgcBand},
  hlines={0.4pt,bkLine}, vline{2-5}={0.4pt,bkLine}, rowsep=2pt, colsep=3pt, cells={font=\footnotesize}}
番号 & 項目（要点） & 扱っている章 & 扱い & 補足 \\
{\sffamily 数A-\allowbreak{}(1)ア(ｱ)} & 三角形の性質 & \kgcch{4-1}（内角の和）、\kgcch{4-2}（外心）、\kgcch{4-11}（重心）、\kgcch{4-7} \kgcmotto{}「3辺から面積：ヘロンの公式と内接円」 & \kgcsym{一部} & 内心・メネラウス・チェバの語はない \\
{\sffamily 数A-\allowbreak{}(1)ア(ｲ)} & 円の性質 & \kgcch{4-2}（円周角・接線） & \kgcsym{一部} & 方べき・接弦定理の語はない \\
{\sffamily 数A-\allowbreak{}(1)ア(ｳ)} & 空間図形 & \kgcch{4-13} & \kgcsym{一部} & 多面体の語はない \\
{\sffamily 数A-\allowbreak{}(1)イ} & 作図 & なし & \kgcsym{なし} &  \\
{\sffamily 数A-\allowbreak{}(2)ア(ｱ)} & 数え上げの原則 & \kgcch{3-17}（表・樹形図で数える） & \kgcsym{一部} & 和の法則・積の法則の名前はない \\
{\sffamily 数A-\allowbreak{}(2)ア(ｲ)} & 順列・組合せ & \kgcch{3-17}「コインの表の枚数を数える：パスカルの三角形」（組合せ・${}_{n}\mathrm{C}_{r}$）、\kgcch{3-17} \kgcmotto{}「数の爆発：トランプの並べ方」 & \kgcsym{一部} & 「順列」の語・${}_{n}\mathrm{P}_{r}$ はない \\
{\sffamily 数A-\allowbreak{}(2)ア(ｳ)} & 確率・期待値 & \kgcch{3-17}「たくさんやったときの平均：期待値」 & \kgcsym{同} &  \\
{\sffamily 数A-\allowbreak{}(2)ア(ｴ)} & 独立な試行 & なし & \kgcsym{なし} & \kgcch{3-17} \kgcmotto{}で「べつべつにふったサイコロの和では、分散がたし算になる」と言うだけ \\
{\sffamily 数A-\allowbreak{}(2)ア(ｵ)} & 条件付き確率 & なし & \kgcsym{なし} &  \\
{\sffamily 数A-\allowbreak{}(2)語} & ${}_{n}\mathrm{P}_{r}$・${}_{n}\mathrm{C}_{r}$・階乗・n!・排反 & ${}_{n}\mathrm{C}_{r}$＝\kgcch{3-17}（記号カード）、階乗・n!＝\kgcch{5-17} & \kgcsym{一部} & ${}_{n}\mathrm{P}_{r}$・排反はない \\
{\sffamily 数A-\allowbreak{}(3)} & 数学と人間の活動（整数の性質を含む） & \kgcch{2-17}（2進数・16進数）、\kgcch{2-4}（素数・公約数・公倍数）、\kgcch{2-9} \kgcmotto{}（偶数・奇数） & \kgcsym{一部} & 内容の取扱い(4)に約数・倍数、ユークリッドの互除法、二進法、座標の考え方が挙がる（互除法はシリーズにない）。シリーズの \kgcch{2-17} は「2進数」「16進数」の語で扱う \\
\end{longtblr}
\Needspace{169pt}\kgbsub{A.12}{高等学校　数学Ⅱ}
\begin{longtblr}[theme=kgc, caption={高等学校　数学Ⅱ}, entry=none, label=none]{colspec={Q[l,wd=19mm] Q[l,wd=31mm] X[2.4,l] Q[l,wd=15mm] X[1.6,l]}, width=\linewidth,
  rowhead=1, row{1}={bg=kgcHead, font=\sffamily\bfseries\footnotesize}, row{even[2]}={bg=kgcBand},
  hlines={0.4pt,bkLine}, vline{2-5}={0.4pt,bkLine}, rowsep=2pt, colsep=3pt, cells={font=\footnotesize}}
番号 & 項目（要点） & 扱っている章 & 扱い & 補足 \\
{\sffamily 数Ⅱ-\allowbreak{}(1)ア(ｱ)} & 三次の乗法公式・因数分解 & \kgcch{3-8}「立方体で見る (a＋b)$^{3}$」 & \kgcsym{一部}\kgcsep{}\kgcsym{先} & 三次の因数分解の公式はない \\
{\sffamily 数Ⅱ-\allowbreak{}(1)ア(ｲ)} & 多項式の除法・分数式 & \kgcch{9-7}（部分分数の前に「2s＋5＝2(s＋2)＋1」の1行） & \kgcsym{一部} & 除法の節はない \\
{\sffamily 数Ⅱ-\allowbreak{}(1)ア(ｳ)} & 複素数の四則 & \kgcch{6-8} & \kgcsym{同}\kgcsep{}\kgcsym{後} & 虚数単位は \textbf{j}（i は電流に使う） \\
{\sffamily 数Ⅱ-\allowbreak{}(1)ア(ｴ)} & 判別、解と係数の関係 & \kgcch{3-10}（判別式）、\kgcch{6-8}「2次方程式の、虚数の解」、\kgcch{6-12}（解と係数の関係） & \kgcsym{同} &  \\
{\sffamily 数Ⅱ-\allowbreak{}(1)ア(ｵ)} & 因数定理・高次方程式 & \kgcch{9-9}「3次の特性方程式：強くしすぎると、こわれる」、\kgcch{6-8} \kgcmotto{}「複素数まで広げれば、…」 & \kgcsym{一部} & 因数定理・剰余の定理の語はない \\
{\sffamily 数Ⅱ-\allowbreak{}(1)イ} & 等式・不等式の証明 & なし & \kgcsym{なし} & 恒等式・係数比較は \kgcch{3-8} にある \\
{\sffamily 数Ⅱ-\allowbreak{}(1)語} & 二項定理・虚数・i & 二項定理＝\kgcch{3-8} \kgcmotto{}「係数のならび：1・2・1、1・3・3・1」（パスカルの三角形）、\kgcch{5-4} \kgcmotto{}「どんな n でも成り立つわけ」（(x＋h)$^{n}$ の係数）、虚数＝\kgcch{6-8} & \kgcsym{一部} & 「二項定理」の語はない。i ではなく j \\
{\sffamily 数Ⅱ-\allowbreak{}(2)ア(ｱ)} & 内分点・外分点・2点間の距離 & \kgcch{4-11}（内分点）、\kgcch{4-4}「座標の上の2点の距離」 & \kgcsym{一部} & 外分点はない \\
{\sffamily 数Ⅱ-\allowbreak{}(2)ア(ｲ)} & 直線と円の方程式 & \kgcch{3-5}（2点・1点と傾きから直線の式）、\kgcch{4-4}（$(x-a)^{2}+(y-b)^{2}$\allowbreak$=r^{2}$）、\kgcch{4-12} \kgcmotto{}「点と直線の距離、直交する2直線」 & \kgcsym{同}\kgcsep{}\kgcsym{先} &  \\
{\sffamily 数Ⅱ-\allowbreak{}(2)ア(ｳ)} & 軌跡 & なし & \kgcsym{なし} & \kgcch{9-9} の「根軌跡」は制御の用語で別もの \\
{\sffamily 数Ⅱ-\allowbreak{}(2)ア(ｴ)} & 不等式の表す領域 & なし & \kgcsym{なし} &  \\
{\sffamily 数Ⅱ-\allowbreak{}(3)ア(ｱ)} & 指数の拡張（整数→\nobreak{}有理数） & \kgcch{2-7}「0乗とマイナスの指数」、\kgcch{2-12}、\kgcch{3-11}「分数の指数」 & \kgcsym{同}\kgcsep{}\kgcsym{先} & 0乗・負の指数は第2巻（\mbox{★〜★★}） \\
{\sffamily 数Ⅱ-\allowbreak{}(3)ア(ｲ)〜(ｴ)} & 指数関数・対数・対数関数 & \kgcch{3-11}、\kgcch{3-12} & \kgcsym{詳}\kgcsep{}\kgcsym{先} & 片対数グラフ・デシベル・底の変換（電卓）まで \\
{\sffamily 数Ⅱ-\allowbreak{}(3)語} & 累乗根・$\log_{a}x$・常用対数 & \kgcch{3-11}、\kgcch{3-12} & \kgcsym{同} &  \\
{\sffamily 数Ⅱ-\allowbreak{}(4)ア(ｱ)} & 一般角・弧度法 & \kgcch{4-8}「何周しても使える角：一般角」、\kgcch{4-2}「角を長さではかる：弧度法とラジアン」 & \kgcsym{同} &  \\
{\sffamily 数Ⅱ-\allowbreak{}(4)ア(ｲ)(ｳ)} & 三角関数のグラフ・相互関係 & \kgcch{4-8} & \kgcsym{同} &  \\
{\sffamily 数Ⅱ-\allowbreak{}(4)ア(ｴ)} & 加法定理・2倍角・合成 & \kgcch{4-9} & \kgcsym{詳} & 加法定理を2通りで導く、積と和の式、うなり、sin$^{2}$ の平均 \\
{\sffamily 数Ⅱ-\allowbreak{}(5)ア(ｱ)} & 微分係数・導関数 & \kgcch{5-1}〜\kgcch{5-4} & \kgcsym{詳} & 数Ⅱ-扱(1)は「三次までの関数を中心」。シリーズは \kgcch{5-4} で x$^{n}$ 一般 \\
{\sffamily 数Ⅱ-\allowbreak{}(5)ア(ｲ)} & 増減・極大極小 & \kgcch{5-7} & \kgcsym{同} &  \\
{\sffamily 数Ⅱ-\allowbreak{}(5)ア(ｳ)} & 不定積分・定積分 & \kgcch{5-9}〜\kgcch{5-11} & \kgcsym{詳} & 数Ⅱ-扱(1)は「二次までの関数を中心」。区分求積（\kgcch{5-10}）から入る \\
{\sffamily 数Ⅱ-\allowbreak{}(5)イ} & 面積 & \kgcch{5-9}、\kgcch{5-11}「横軸より下の部分」 & \kgcsym{同} &  \\
{\sffamily 数Ⅱ-\allowbreak{}(5)語} & 極限値・lim & \kgcch{5-3}（記号カード lim） & \kgcsym{同} & 極限は数値で近づく様子で（$\varepsilon$-$\delta$は使わない） \\
\end{longtblr}
\Needspace{199pt}\kgbsub{A.13}{高等学校　数学B}
\begin{longtblr}[theme=kgc, caption={高等学校　数学B}, entry=none, label=none]{colspec={Q[l,wd=19mm] Q[l,wd=31mm] X[2.4,l] Q[l,wd=15mm] X[1.6,l]}, width=\linewidth,
  rowhead=1, row{1}={bg=kgcHead, font=\sffamily\bfseries\footnotesize}, row{even[2]}={bg=kgcBand},
  hlines={0.4pt,bkLine}, vline{2-5}={0.4pt,bkLine}, rowsep=2pt, colsep=3pt, cells={font=\footnotesize}}
番号 & 項目（要点） & 扱っている章 & 扱い & 補足 \\
{\sffamily 数B-\allowbreak{}(1)ア(ｱ)} & 等差数列・等比数列 & \kgcch{7-10}（等比数列・等比級数） & \kgcsym{一部} & \textbf{等差数列はない}（\kgcch{6-7} で「等差」の語が1回出るだけ） \\
{\sffamily 数B-\allowbreak{}(1)ア(ｲ)} & いろいろな数列の和 & \kgcch{5-10}「たし算をまとめて書く：$\Sigma$」（$1^{2}+2^{2}+\dots+n^{2}$\allowbreak$=n(n+1)(2n+1)/6$） & \kgcsym{一部} &  \\
{\sffamily 数B-\allowbreak{}(1)ア(ｳ)} & 漸化式 & \kgcch{6-2}（オイラー法 $y_{n+1}$＝…の形。「漸化式」の語はない）、\kgcch{7-10}「行列のくり返し：漸化式とオイラー法を行列で」 & \kgcsym{一部} &  \\
{\sffamily 数B-\allowbreak{}(1)ア(ｴ)} & 数学的帰納法 & なし & \kgcsym{なし} & \kgcch{5-4} \kgcmotto{}「どんな n でも成り立つわけ」は (x＋h)$^{n}$ の係数（パスカルの三角形）による説明で、帰納法ではない \\
{\sffamily 数B-\allowbreak{}(1)語} & $\Sigma$ & \kgcch{5-10} & \kgcsym{同} &  \\
{\sffamily 数B-\allowbreak{}(2)ア(ｱ)} & 標本調査の考え方 & \kgcch{3-17}、\kgcch{3-14}「何回もはかって平均すると、ずれは $\sigma/\sqrt{n}$」 & \kgcsym{一部} &  \\
{\sffamily 数B-\allowbreak{}(2)ア(ｲ)} & 確率変数と確率分布 & \kgcch{3-17}（コインの表の枚数の分布、\kgcmotto{}でサイコロの和の分散）、\kgcch{6-4}「めったに起きないことの待ち時間」（指数分布・ポアソン分布） & \kgcsym{一部} & 「確率変数」「確率分布」の語はない \\
{\sffamily 数B-\allowbreak{}(2)ア(ｳ)} & 二項分布・正規分布 & \kgcch{3-17}「たくさん足すと、つりがね形：正規分布」（±$\sigma$ に約68 \%、±2$\sigma$ に約95 \%）、\kgcch{3-17} のパスカルの三角形 & \kgcsym{一部} & 「二項分布」の語はない \\
{\sffamily 数B-\allowbreak{}(2)ア(ｴ)} & 区間推定・仮説検定 & 仮説検定＝\kgcch{3-17} \kgcmotto{}（\mbox{★★★}）。区間推定＝なし & \kgcsym{一部}\kgcor{}\kgcsym{なし} &  \\
{\sffamily 数B-\allowbreak{}(2)語} & 信頼区間・有意水準 & なし & \kgcsym{なし} &  \\
{\sffamily 数B-\allowbreak{}(3)} & 数学と社会生活（扱(3)散布図のデータを関数とみなす） & \kgcch{3-16}（最小二乗法・回帰直線） & \kgcsym{詳} &  \\
\end{longtblr}
\Needspace{229pt}\kgbsub{A.14}{高等学校　数学Ⅲ}
\begin{longtblr}[theme=kgc, caption={高等学校　数学Ⅲ}, entry=none, label=none]{colspec={Q[l,wd=19mm] Q[l,wd=31mm] X[2.4,l] Q[l,wd=15mm] X[1.6,l]}, width=\linewidth,
  rowhead=1, row{1}={bg=kgcHead, font=\sffamily\bfseries\footnotesize}, row{even[2]}={bg=kgcBand},
  hlines={0.4pt,bkLine}, vline{2-5}={0.4pt,bkLine}, rowsep=2pt, colsep=3pt, cells={font=\footnotesize}}
番号 & 項目（要点） & 扱っている章 & 扱い & 補足 \\
{\sffamily 数Ⅲ-\allowbreak{}(1)ア(ｱ)} & 数列の極限・$\{r^{n}\}$ & \kgcch{5-10}「n を大きくした先を、式で求める」、\kgcch{5-14}「$(1+1/n)^{n}$」、\kgcch{7-10}（$\lambda^{n}$\allowbreak$=r^{n}e^{jn\theta}$、固有値の大きさで「ふえる・消える」） & \kgcsym{一部} & 数列の極限の節はない \\
{\sffamily 数Ⅲ-\allowbreak{}(1)ア(ｲ)} & 無限級数・無限等比級数 & \kgcch{7-10}「くり返しの合計：等比級数」（$1+x+x^{2}+\dots$\allowbreak$=1/(1-x)$）、\kgcch{5-17}（テイラー展開）、\kgcch{9-11}「無限に続くたし算で書く」 & \kgcsym{一部} &  \\
{\sffamily 数Ⅲ-\allowbreak{}(1)ア(ｳ)} & 分数関数・無理関数 & \kgcch{3-7}（y＝a/x）、\kgcch{3-7} \kgcmotto{}（1/x$^{2}$） & \kgcsym{一部} & 平行移動した分数関数、無理関数のグラフはない \\
{\sffamily 数Ⅲ-\allowbreak{}(1)ア(ｴ)} & 合成関数・逆関数 & \kgcch{3-3}「箱をつなぐ：合成関数」、\kgcch{3-12}（逆関数：本文の囲み）、\kgcch{5-14}（ln） & \kgcsym{同}\kgcsep{}\kgcsym{先} &  \\
{\sffamily 数Ⅲ-\allowbreak{}(1)ア(ｵ)} & 関数の値の極限 & \kgcch{5-3}「かぎりなく近づく：極限と lim」 & \kgcsym{同} &  \\
{\sffamily 数Ⅲ-\allowbreak{}(2)ア(ｱ)} & 微分可能性・積・商 & \kgcch{5-6}（積の微分）、\kgcch{5-6} \kgcmotto{}「商の微分」、\kgcch{5-3} \kgcmotto{}「とがった所には、接線が引けない」 & \kgcsym{同}\newline{}（商は一部） &  \\
{\sffamily 数Ⅲ-\allowbreak{}(2)ア(ｲ)} & 合成関数の導関数 & \kgcch{5-6} & \kgcsym{同} &  \\
{\sffamily 数Ⅲ-\allowbreak{}(2)ア(ｳ)} & 三角・指数・対数関数の導関数 & \kgcch{5-14}、\kgcch{5-15} & \kgcsym{同} &  \\
{\sffamily 数Ⅲ-\allowbreak{}(2)ア(ｴ)} & 接線・増減・凹凸・グラフ & \kgcch{5-3}「接線の式」、\kgcch{5-7}「2回微分で、山と谷を見分ける」 & \kgcsym{一部} & 「変曲点」の語はない。「下に凸」は \kgcch{5-17} に出る \\
{\sffamily 数Ⅲ-\allowbreak{}(2)語} & 自然対数・e・変曲点 & \kgcch{5-14}（e・ln）、\kgcch{5-5}「2回微分する」 & \kgcsym{一部} & 変曲点なし。第二次導関数は〔用語・記号〕ではなく(2)イ(ｲ)の本文にある \\
{\sffamily 数Ⅲ-\allowbreak{}(3)ア(ｱ)} & 不定積分・定積分の性質 & \kgcch{5-11} & \kgcsym{同} &  \\
{\sffamily 数Ⅲ-\allowbreak{}(3)ア(ｲ)} & 置換積分・部分積分 & \kgcch{5-16} & \kgcsym{同} &  \\
{\sffamily 数Ⅲ-\allowbreak{}(3)ア(ｳ)} & 面積・体積・曲線の長さ & 面積＝\kgcch{5-11}。体積＝\kgcch{5-10} の例題「うすい板に分けて体積を求める」（円すい）と章末問題（四角すい）。曲線の長さ＝\kgcch{6-11} のせいかくに言うと（$s=\int\!\sqrt{1+y'^{2}}\,dx$）、\kgcch{8-7}「長さでたす線積分」 & \kgcsym{一部} & 回転体の体積の計算演習・曲線の長さの節はない \\
{\sffamily 数Ⅲ-\allowbreak{}扱(1)} & 直線上・平面上の点の速度・加速度 & \kgcch{5-5}、\kgcch{5-15}「円の上を回る点の速さと向き」、\kgcch{4-11} & \kgcsym{詳} &  \\
{\sffamily 数Ⅲ-\allowbreak{}注} & 微分方程式（本文になし） & 第6巻全体 & \kgcsym{範囲外} &  \\
\end{longtblr}
\Needspace{169pt}\kgbsub{A.15}{高等学校　数学C}
\begin{longtblr}[theme=kgc, caption={高等学校　数学C}, entry=none, label=none]{colspec={Q[l,wd=19mm] Q[l,wd=31mm] X[2.4,l] Q[l,wd=15mm] X[1.6,l]}, width=\linewidth,
  rowhead=1, row{1}={bg=kgcHead, font=\sffamily\bfseries\footnotesize}, row{even[2]}={bg=kgcBand},
  hlines={0.4pt,bkLine}, vline{2-5}={0.4pt,bkLine}, rowsep=2pt, colsep=3pt, cells={font=\footnotesize}}
番号 & 項目（要点） & 扱っている章 & 扱い & 補足 \\
{\sffamily 数C-\allowbreak{}(1)ア(ｱ)} & ベクトル・位置ベクトル・成分 & \kgcch{4-10}、\kgcch{4-11} & \kgcsym{同} & \kgcch{4-10} は\mbox{★★}、\kgcch{4-11} は\mbox{★★〜★★★} \\
{\sffamily 数C-\allowbreak{}(1)ア(ｲ)} & 内積 & \kgcch{4-12} & \kgcsym{詳} & 仕事・正射影・相関係数 \\
{\sffamily 数C-\allowbreak{}(1)ア(ｳ)} & 空間への拡張 & \kgcch{4-13} & \kgcsym{詳} & 法線ベクトル・空間の直線と平面・右手系。外積は \kgcch{4-14}（範囲外） \\
{\sffamily 数C-\allowbreak{}(2)ア(ｱ)} & 放物線・楕円・双曲線（二次曲線） & 放物線＝\kgcch{3-9}（パラボラアンテナの焦点）。だ円＝\kgcch{4-8} \kgcmotto{}「円をつぶすと、だ円」（\mbox{★★★★}）、\kgcch{4-15}（極方程式）、\kgcch{8-1}「等高線の式は、だ円」、\kgcch{7-9} \kgcmotto{}。双曲線＝\kgcch{3-7}（y＝a/x のグラフとして） & \kgcsym{一部} & 焦点・準線で3つをまとめる扱いはない。語は「楕円」でなく「だ円」 \\
{\sffamily 数C-\allowbreak{}(2)ア(ｲ)} & 媒介変数表示 & \kgcch{4-8}（だ円の媒介変数）、\kgcch{4-11}（$\vec{p}=\vec{p}_{0}+t\vec{v}$）、\kgcch{8-7}「道を1つの数 t で表す」 & \kgcsym{同} &  \\
{\sffamily 数C-\allowbreak{}(2)ア(ｳ)} & 極座標・極方程式 & \kgcch{4-2}「距離と向きで場所を言う：極座標」、\kgcch{4-15}「曲線を、極座標の式で表す」 & \kgcsym{同} &  \\
{\sffamily 数C-\allowbreak{}(2)ア(ｴ)} & 複素数平面・極形式・四則の図形的意味 & \kgcch{6-9}「複素数を点で表す：複素平面」「向きの角：偏角と極形式」 & \kgcsym{同} & 語は「複素平面」 \\
{\sffamily 数C-\allowbreak{}(2)ア(ｵ)} & ド・モアブルの定理 & \kgcch{6-9}「何回もかける：ド・モアブルの定理」 & \kgcsym{同} & \kgcmotto{}「1の n 乗根」 \\
{\sffamily 数C-\allowbreak{}(3)ア(ｱ)} & 図・表・統計グラフ & \kgcch{1-17}、\kgcch{1-18}、\kgcch{3-4}、\kgcch{3-13} & \kgcsym{同} &  \\
{\sffamily 数C-\allowbreak{}(3)ア(ｲ)} & 離散グラフや行列で表現する意義 & 行列＝第7巻全体（\kgcch{3-6}「係数を表にならべる：行列の種」、\kgcch{1-13} が種） & \kgcsym{詳}\newline{}（行列） & 行列は演算・行列式・逆行列・固有値・対角化まで（指導要領の範囲外） \\
\end{longtblr}
\bkevenpage\kgbsec{B}{付録B　全152章の一覧}\label{kgc-b}\label{kgc-sec11}
この付録は、全9巻・152章（序章9・章139・読みもの3・終章1）を、巻と章の順に並べた一覧です。章名・★・路線は各巻のもくじと同じです。「指導要領の学年・科目」は、その章の内容が当たる学習指導要領の学年や科目で、1つの章が複数の学年にまたがるときは「入口〜出口」（最も低い学年〜最も高い学年・科目）で示しています。

\begin{kgclegend}{表の読み方}
\begin{tblr}{colspec={Q[l,wd=24mm] X[l]}, rowsep=2pt, colsep=3pt, cells={font=\small, valign=m}}
★ & 章の内容の高さの目安です。塗りの★はその章の下の高さ、白抜きの☆は上の高さまでを表します（例：\kgcstars{2}{3}＝★★〜★★★）。★は指導要領の学年と一対一ではありません。 \\
路線 & 7本の路線のどの駅に当たるかを、形の記号と番号で示しています。\par\smallskip{\raggedright \mbox{\kgcroute{1}\,変化（ひき算→傾き→微分）}\hspace{1.2em plus 1fill}\mbox{\kgcroute{2}\,たし合わせ（かけ算→面積→積分）}\hspace{1.2em plus 1fill}\mbox{\kgcroute{3}\,倍々（累乗→指数→e）}\hspace{1.2em plus 1fill}\mbox{\kgcroute{4}\,回転と波（角→sin→複素数）}\hspace{1.2em plus 1fill}\mbox{\kgcroute{5}\,表（九九の表→行列）}\hspace{1.2em plus 1fill}\mbox{\kgcroute{6}\,矢印（向き→ベクトル→場）}\hspace{1.2em plus 1fill}\mbox{\kgcroute{7}\,はかる（目もり→誤差→統計）}\par} \\
指導要領の学年・科目 & 小＝小学校、中＝中学校、数Ⅰ〜数C＝高等学校の科目（数学Ⅰ〜数学C）、大学＝指導要領の範囲外で、大学の工学部で一般に学ぶ内容です。「範囲外」は、指導要領にも大学の科目の一覧にも項目がない内容です。「―」は、特定の学年に当てはめていない章です（1-0・2-0・2-18・2-19・4-0・5-0・9-13）。 \\
範囲外 & \kgcout{a}＝指導要領の範囲外の内容が中心の章、\kgcout{p}＝一部に指導要領の範囲外の内容を含む章、印なし＝指導要領の範囲の中の章です。 \\
\end{tblr}
\end{kgclegend}

\Needspace{8\baselineskip}\kgbsub{第1巻}{かぞえる・はかる}\par\nobreak{\small\noindent 学年の目安：小学1年〜3年　章数：22}\par\nobreak\vspace{3pt}
\begin{longtblr}[theme=kgc, caption={第1巻　かぞえる・はかる}, entry=none, label=none]{colspec={Q[l,wd=14mm] X[l] Q[l,wd=17mm] Q[l,wd=12.5mm] Q[l,wd=25mm] Q[c,wd=10mm]}, width=\linewidth,
  rowhead=1, row{1}={bg=kgcHead, font=\sffamily\bfseries\footnotesize}, row{even[2]}={bg=kgcBand},
  hlines={0.4pt,bkLine}, vline{2-6}={0.4pt,bkLine}, rowsep=1.5pt, colsep=3pt, cells={font=\footnotesize, valign=m}, row{1}={valign=m}}
章 & 章名 & ★ & 路線 & 指導要領の学年・科目 & 範囲外 \\
{\sffamily\bfseries 1-0}\newline{\footnotesize\sffamily 序章} & かずを つかうと、なにが いいの？ & \kgcstars{1}{1} & \kgcroute{7} & ― &  \\
{\sffamily\bfseries 1-1} & いくつ あるかな？ & \kgcstars{1}{1} & \kgcroute{2} & 小1 &  \\
{\sffamily\bfseries 1-2} & 10の まとまりを つくろう & \kgcstars{1}{1} & \kgcroute{3} & 小1〜小3 &  \\
{\sffamily\bfseries 1-3} & なにも ないときの「0」 & \kgcstars{1}{1} & \kgcroute{6} & 小1〜小2 &  \\
{\sffamily\bfseries 1-4} & たし算は「あわせる」「ふえる」 & \kgcstars{1}{1} & \kgcroute{2} & 小1 &  \\
{\sffamily\bfseries 1-5} & ひき算は「のこり」「ちがい」「かわった 量」 & \kgcstars{1}{1} & \kgcroute{1} & 小1 &  \\
{\sffamily\bfseries 1-6} & かずの 線：数直線 & \kgcstars{1}{1} & \kgcroute{6} & 小1〜小3 &  \\
{\sffamily\bfseries 1-7} & 大きな かずの たし算・ひき算 & \kgcstars{1}{1} & \kgcroute{2} & 小1〜小3 &  \\
{\sffamily\bfseries 1-8} & ながさを はかろう & \kgcstars{1}{1} & \kgcroute{7} & 小1〜小2 &  \\
{\sffamily\bfseries 1-9} & かさと おもさ & \kgcstars{1}{1} & \kgcroute{7} & 小1〜小3 &  \\
{\sffamily\bfseries 1-10} & 時こくと 時間 & \kgcstars{1}{1} & \kgcroute{1} & 小1〜小3 &  \\
{\sffamily\bfseries 1-11} & 目もりを よもう & \kgcstars{1}{1} & \kgcroute{7} & 小2〜小3 &  \\
{\sffamily\bfseries 1-12} & かけ算は「同じ かずの まとまり」 & \kgcstars{1}{1} & \kgcroute{2} & 小2 &  \\
{\sffamily\bfseries 1-13} & 九九の 表を たんけんしよう & \kgcstars{1}{1} & \kgcroute{5} & 小2〜小3 &  \\
{\sffamily\bfseries 1-14} & わり算は「同じ かずずつ 分ける」 & \kgcstars{1}{1} & \kgcroute{1} & 小3 &  \\
{\sffamily\bfseries 1-15} & あまりの ある わり算 & \kgcstars{1}{1} & \kgcroute{4} & 小3〜小4 &  \\
{\sffamily\bfseries 1-16} & 1より 小さい かず：小数と 分数 & \kgcstars{1}{1} & \kgcroute{7} & 小2〜小5 &  \\
{\sffamily\bfseries 1-17} & 表と ぼうグラフ & \kgcstars{1}{1} & \kgcroute{5} & 小1〜小3 &  \\
{\sffamily\bfseries 1-18} & おれ線グラフ：かわり方を 見る & \kgcstars{1}{1} & \kgcroute{1} & 小4 &  \\
{\sffamily\bfseries 1-19} & かたち：三角・四角・円 & \kgcstars{1}{1} & \kgcroute{4} & 小1〜小3 &  \\
{\sffamily\bfseries 1-20} & 角の 大きさ & \kgcstars{1}{1} & \kgcroute{4} & 小3〜小4 &  \\
{\sffamily\bfseries 1-21} & 広さを くらべよう：面積 & \kgcstars{1}{1} & \kgcroute{2} & 小1〜小4 &  \\
\end{longtblr}
\Needspace{8\baselineskip}\kgbsub{第2巻}{数と量のことば}\par\nobreak{\small\noindent 学年の目安：小学4年〜中学1年　章数：20}\par\nobreak\vspace{3pt}
\begin{longtblr}[theme=kgc, caption={第2巻　数と量のことば}, entry=none, label=none]{colspec={Q[l,wd=14mm] X[l] Q[l,wd=17mm] Q[l,wd=12.5mm] Q[l,wd=25mm] Q[c,wd=10mm]}, width=\linewidth,
  rowhead=1, row{1}={bg=kgcHead, font=\sffamily\bfseries\footnotesize}, row{even[2]}={bg=kgcBand},
  hlines={0.4pt,bkLine}, vline{2-6}={0.4pt,bkLine}, rowsep=1.5pt, colsep=3pt, cells={font=\footnotesize, valign=m}, row{1}={valign=m}}
章 & 章名 & ★ & 路線 & 指導要領の学年・科目 & 範囲外 \\
{\sffamily\bfseries 2-0}\newline{\footnotesize\sffamily 序章} & 「だんだん速く」を数で言うと？ & \kgcstars{1}{1} & \kgcroute{1} & ― &  \\
{\sffamily\bfseries 2-1} & 「かぞえる数」と「はかる数」 & \kgcstars{1}{1} & \kgcroute{7} & 小3〜小5 &  \\
{\sffamily\bfseries 2-2} & マイナスの数と数直線 & \kgcstars{1}{2} & \kgcroute{6} & 小2〜中1 &  \\
{\sffamily\bfseries 2-3} & ＋ － × ÷ は、何を表しているの？ & \kgcstars{1}{1} & \kgcroute{1}\par \kgcroute{2} & 小1〜小5 &  \\
{\sffamily\bfseries 2-4} & 分数と小数の計算 & \kgcstars{1}{2} & \kgcroute{1} & 小4〜中1 &  \\
{\sffamily\bfseries 2-5} & 比・割合・百分率 & \kgcstars{1}{2} & \kgcroute{7} & 小4〜小6 &  \\
{\sffamily\bfseries 2-6} & はかった数には「ずれ」がある & \kgcstars{1}{2} & \kgcroute{7} & 小4〜中3 &  \\
{\sffamily\bfseries 2-7} & 何回もかける：累乗と指数 & \kgcstars{1}{2} & \kgcroute{3} & 小4〜数Ⅱ &  \\
{\sffamily\bfseries 2-8} & 2回かけると□になる数：平方根 & \kgcstars{2}{2} & \kgcroute{3} & 中3 &  \\
{\sffamily\bfseries 2-9} & 分数で書けない数がある？ & \kgcstars{2}{2} & \kgcroute{6} & 小5〜数Ⅰ &  \\
{\sffamily\bfseries 2-10} & 単位の世界ルール：SI単位 & \kgcstars{1}{2} & \kgcroute{7} & 小2〜小5 & \kgcout{p} \\
{\sffamily\bfseries 2-11} & なぜ文字を使うの？ & \kgcstars{2}{2} & \kgcroute{5} & 小2〜中1 &  \\
{\sffamily\bfseries 2-12} & 文字式のやくそく & \kgcstars{2}{2} & \kgcroute{5} & 中1〜数Ⅱ &  \\
{\sffamily\bfseries 2-13} & 公式に数を入れてみよう：代入 & \kgcstars{2}{2} & \kgcroute{5} & 小4〜中1 &  \\
{\sffamily\bfseries 2-14} & 式を変形して、ほしい量をとり出そう & \kgcstars{2}{2} & \kgcroute{5} & 中1〜中2 &  \\
{\sffamily\bfseries 2-15} & 比例と反比例 & \kgcstars{1}{2} & \kgcroute{1} & 小5〜中1 &  \\
{\sffamily\bfseries 2-16} & 何秒後に追いつく？1次方程式 & \kgcstars{2}{2} & \kgcroute{1} & 中1 &  \\
{\sffamily\bfseries 2-17} & コンピュータの数：2進数と16進数 & \kgcstars{2}{2} & \kgcroute{3} & 数A & \kgcout{p} \\
{\sffamily\bfseries 2-18} & 表計算で計算しよう & \kgcstars{2}{2} & \kgcroute{5} & ― & \kgcout{p} \\
{\sffamily\bfseries 2-19} & プログラムで計算しよう & \kgcstars{2}{2} & \kgcroute{1} & ― & \kgcout{p} \\
\end{longtblr}
\Needspace{8\baselineskip}\kgbsub{第3巻}{関数とデータ　きまりを見つける}\par\nobreak{\small\noindent 学年の目安：中学1年〜3年　章数：18}\par\nobreak\vspace{3pt}
\begin{longtblr}[theme=kgc, caption={第3巻　関数とデータ}, entry=none, label=none]{colspec={Q[l,wd=14mm] X[l] Q[l,wd=17mm] Q[l,wd=12.5mm] Q[l,wd=25mm] Q[c,wd=10mm]}, width=\linewidth,
  rowhead=1, row{1}={bg=kgcHead, font=\sffamily\bfseries\footnotesize}, row{even[2]}={bg=kgcBand},
  hlines={0.4pt,bkLine}, vline{2-6}={0.4pt,bkLine}, rowsep=1.5pt, colsep=3pt, cells={font=\footnotesize, valign=m}, row{1}={valign=m}}
章 & 章名 & ★ & 路線 & 指導要領の学年・科目 & 範囲外 \\
{\sffamily\bfseries 3-0}\newline{\footnotesize\sffamily 序章} & はかったデータから、きまりを見つけよう & \kgcstars{2}{2} & \kgcroute{7} & 中1 &  \\
{\sffamily\bfseries 3-1} & 変わる量と変わらない量 & \kgcstars{1}{1} & \kgcroute{7} & 中1 &  \\
{\sffamily\bfseries 3-2} & 場所を数で表す：座標 & \kgcstars{1}{2} & \kgcroute{6} & 小4〜中1 &  \\
{\sffamily\bfseries 3-3} & 関数は「入れると出てくる」ルール & \kgcstars{1}{2} & \kgcroute{5} & 中1〜数Ⅲ & \kgcout{p} \\
{\sffamily\bfseries 3-4} & グラフをかこう・読もう & \kgcstars{1}{1} & \kgcroute{7} & 中1 &  \\
{\sffamily\bfseries 3-5} & グラフの傾きは「変化の速さ」 & \kgcstars{2}{2} & \kgcroute{1} & 中2〜数Ⅱ &  \\
{\sffamily\bfseries 3-6} & 2本の直線が交わるところ：連立方程式 & \kgcstars{2}{2} & \kgcroute{5} & 中2 &  \\
{\sffamily\bfseries 3-7} & 反比例のグラフ & \kgcstars{2}{2} & \kgcroute{1} & 小6〜数Ⅲ &  \\
{\sffamily\bfseries 3-8} & かっこをはずす：展開と平方完成 & \kgcstars{2}{3} & \kgcroute{1} & 中3〜数Ⅱ &  \\
{\sffamily\bfseries 3-9} & 放物線と $y=ax^2$ & \kgcstars{2}{3} & \kgcroute{1} & 中3〜数Ⅰ &  \\
{\sffamily\bfseries 3-10} & ボールはいつその高さに届く？：2次方程式 & \kgcstars{2}{3} & \kgcroute{1} & 中3〜数Ⅱ &  \\
{\sffamily\bfseries 3-11} & 指数関数：いまある量に比例してふえる・へる & \kgcstars{2}{3} & \kgcroute{3} & 数Ⅱ &  \\
{\sffamily\bfseries 3-12} & 対数：けたを数える道具 & \kgcstars{2}{3} & \kgcroute{3} & 数Ⅱ〜数Ⅲ & \kgcout{p} \\
{\sffamily\bfseries 3-13} & データのちらばり：ヒストグラムと平均 & \kgcstars{2}{2} & \kgcroute{7} & 小6〜数Ⅰ &  \\
{\sffamily\bfseries 3-14} & ばらつきを1つの数で：標準偏差 & \kgcstars{3}{3} & \kgcroute{7} & 数Ⅰ〜数B &  \\
{\sffamily\bfseries 3-15} & 2つの量の関係：散布図と相関 & \kgcstars{3}{3} & \kgcroute{7} & 数Ⅰ &  \\
{\sffamily\bfseries 3-16} & データに直線を引く：最小二乗法 & \kgcstars{3}{3} & \kgcroute{7} & 数B〜大学 & \kgcout{p} \\
{\sffamily\bfseries 3-17} & 偶然のばらつき：確率と正規分布 & \kgcstars{3}{3} & \kgcroute{7} & 小6〜数B &  \\
\end{longtblr}
\Needspace{8\baselineskip}\kgbsub{第4巻}{形と向き　図形・三角関数・ベクトル}\par\nobreak{\small\noindent 学年の目安：中学2年〜高校2年　章数：16}\par\nobreak\vspace{3pt}
\begin{longtblr}[theme=kgc, caption={第4巻　形と向き}, entry=none, label=none]{colspec={Q[l,wd=14mm] X[l] Q[l,wd=17mm] Q[l,wd=12.5mm] Q[l,wd=25mm] Q[c,wd=10mm]}, width=\linewidth,
  rowhead=1, row{1}={bg=kgcHead, font=\sffamily\bfseries\footnotesize}, row{even[2]}={bg=kgcBand},
  hlines={0.4pt,bkLine}, vline{2-6}={0.4pt,bkLine}, rowsep=1.5pt, colsep=3pt, cells={font=\footnotesize, valign=m}, row{1}={valign=m}}
章 & 章名 & ★ & 路線 & 指導要領の学年・科目 & 範囲外 \\
{\sffamily\bfseries 4-0}\newline{\footnotesize\sffamily 序章} & ロボットの手は、どこにある？ & \kgcstars{2}{2} & \kgcroute{4} & ― & \kgcout{p} \\
{\sffamily\bfseries 4-1} & きまりから考えを積み上げる：定義・公理・定理 & \kgcstars{2}{3} & \kgcroute{6} & 小3〜数Ⅰ &  \\
{\sffamily\bfseries 4-2} & 角度と円、ラジアン & \kgcstars{1}{2} & \kgcroute{4} & 小3〜数C &  \\
{\sffamily\bfseries 4-3} & にた形と比：相似 & \kgcstars{2}{3} & \kgcroute{6} & 小6〜中3 &  \\
{\sffamily\bfseries 4-4} & 三平方の定理 & \kgcstars{3}{3} & \kgcroute{6} & 中3〜数Ⅱ &  \\
{\sffamily\bfseries 4-5} & 面積と体積：細かく分けてたす & \kgcstars{1}{2} & \kgcroute{2} & 小4〜中3 &  \\
{\sffamily\bfseries 4-6} & 坂道と三角比 & \kgcstars{2}{3} & \kgcroute{4} & 数Ⅰ & \kgcout{p} \\
{\sffamily\bfseries 4-7} & 三角形を解く：正弦定理と余弦定理 & \kgcstars{3}{3} & \kgcroute{4} & 数Ⅰ〜数A & \kgcout{p} \\
{\sffamily\bfseries 4-8} & 三角関数：円運動の影 & \kgcstars{3}{4} & \kgcroute{4} & 数Ⅱ〜数C &  \\
{\sffamily\bfseries 4-9} & 波をたし合わせる：加法定理とうなり & \kgcstars{4}{4} & \kgcroute{4} & 数Ⅱ & \kgcout{p} \\
{\sffamily\bfseries 4-10} & 向きと大きさをもつ量：ベクトル & \kgcstars{2}{2} & \kgcroute{6} & 数C &  \\
{\sffamily\bfseries 4-11} & ベクトルを成分で計算しよう & \kgcstars{2}{3} & \kgcroute{6} & 数Ⅱ〜数C &  \\
{\sffamily\bfseries 4-12} & 内積：どれだけ同じ向きか & \kgcstars{3}{4} & \kgcroute{6} & 数Ⅰ〜数C &  \\
{\sffamily\bfseries 4-13} & 空間の座標と3次元のベクトル & \kgcstars{4}{4} & \kgcroute{6} & 小4〜数C &  \\
{\sffamily\bfseries 4-14} & 外積：回す力と右ねじ & \kgcstars{4}{5} & \kgcroute{6} & 大学 & \kgcout{a} \\
{\sffamily\bfseries 4-15} & いろいろな座標：極座標・円柱座標・球座標 & \kgcstars{4}{5} & \kgcroute{6} & 数C〜大学 & \kgcout{p} \\
\end{longtblr}
\Needspace{8\baselineskip}\kgbsub{第5巻}{変化をつかむ　微分と積分}\par\nobreak{\small\noindent 学年の目安：高校1年〜2年　章数：19}\par\nobreak\vspace{3pt}
\begin{longtblr}[theme=kgc, caption={第5巻　変化をつかむ}, entry=none, label=none]{colspec={Q[l,wd=14mm] X[l] Q[l,wd=17mm] Q[l,wd=12.5mm] Q[l,wd=25mm] Q[c,wd=10mm]}, width=\linewidth,
  rowhead=1, row{1}={bg=kgcHead, font=\sffamily\bfseries\footnotesize}, row{even[2]}={bg=kgcBand},
  hlines={0.4pt,bkLine}, vline{2-6}={0.4pt,bkLine}, rowsep=1.5pt, colsep=3pt, cells={font=\footnotesize, valign=m}, row{1}={valign=m}}
章 & 章名 & ★ & 路線 & 指導要領の学年・科目 & 範囲外 \\
{\sffamily\bfseries 5-0}\newline{\footnotesize\sffamily 序章} & スピードメーターは何を計算している？ & \kgcstars{2}{2} & \kgcroute{1} & ― &  \\
{\sffamily\bfseries 5-1} & 平均の速さと瞬間の速さ & \kgcstars{2}{2} & \kgcroute{1} & 小5〜数Ⅱ &  \\
{\sffamily\bfseries 5-2} & 時間をどんどん短くしてみよう & \kgcstars{2}{2} & \kgcroute{1} & 数Ⅱ &  \\
{\sffamily\bfseries 5-3} & かぎりなく近づく：極限と接線 & \kgcstars{3}{3} & \kgcroute{1} & 数Ⅱ〜数Ⅲ &  \\
{\sffamily\bfseries 5-4} & 表から導関数を見つけよう & \kgcstars{3}{3} & \kgcroute{1} & 数Ⅱ〜数Ⅲ &  \\
{\sffamily\bfseries 5-5} & 位置→速度→加速度 & \kgcstars{3}{3} & \kgcroute{1} & 数Ⅲ & \kgcout{p} \\
{\sffamily\bfseries 5-6} & 微分の道具箱：積と合成関数 & \kgcstars{4}{4} & \kgcroute{1} & 数Ⅲ &  \\
{\sffamily\bfseries 5-7} & いちばん高いのはどこ？　最大と最小 & \kgcstars{3}{3} & \kgcroute{1} & 数Ⅱ〜数Ⅲ &  \\
{\sffamily\bfseries 5-8} & 短い時間なら、かけ算で近似できる & \kgcstars{2}{2} & \kgcroute{1} & 範囲外 & \kgcout{a} \\
{\sffamily\bfseries 5-9} & グラフの面積は進んだ距離 & \kgcstars{2}{2} & \kgcroute{2} & 数Ⅱ &  \\
{\sffamily\bfseries 5-10} & 細い長方形に分けてたそう & \kgcstars{2}{4} & \kgcroute{2} & 数Ⅱ〜数Ⅲ &  \\
{\sffamily\bfseries 5-11} & 積分は微分の逆 & \kgcstars{3}{3} & \kgcroute{2} & 数Ⅱ〜数Ⅲ &  \\
{\sffamily\bfseries 5-12} & 物理での積分：仕事と流量 & \kgcstars{3}{3} & \kgcroute{2} & 数Ⅱ〜数Ⅲ &  \\
{\sffamily\bfseries 5-13} & 積分定数は「はじめの状態」 & \kgcstars{3}{3} & \kgcroute{2} & 数Ⅱ〜数Ⅲ & \kgcout{p} \\
{\sffamily\bfseries 5-14} & ふしぎな数 e & \kgcstars{3}{4} & \kgcroute{3} & 数Ⅲ &  \\
{\sffamily\bfseries 5-15} & sin・cosを微分すると？ & \kgcstars{4}{4} & \kgcroute{4} & 数Ⅲ & \kgcout{p} \\
{\sffamily\bfseries 5-16} & 積分の道具箱：置換積分と部分積分 & \kgcstars{4}{4} & \kgcroute{2} & 数Ⅲ &  \\
{\sffamily\bfseries 5-17} & 近くでは多項式で近似できる：テイラー展開 & \kgcstars{4}{4} & \kgcroute{1} & 大学 & \kgcout{a} \\
{\sffamily\bfseries 5-18} & まとめ：微分したときのふるまい & \kgcstars{3}{3} & \kgcroute{1} & 数Ⅱ〜大学 & \kgcout{p} \\
\end{longtblr}
\Needspace{8\baselineskip}\kgbsub{第6巻}{未来を計算する　微分方程式と複素数}\par\nobreak{\small\noindent 学年の目安：高校2年〜3年　章数：16}\par\nobreak\vspace{3pt}
\begin{longtblr}[theme=kgc, caption={第6巻　未来を計算する}, entry=none, label=none]{colspec={Q[l,wd=14mm] X[l] Q[l,wd=17mm] Q[l,wd=12.5mm] Q[l,wd=25mm] Q[c,wd=10mm]}, width=\linewidth,
  rowhead=1, row{1}={bg=kgcHead, font=\sffamily\bfseries\footnotesize}, row{even[2]}={bg=kgcBand},
  hlines={0.4pt,bkLine}, vline{2-6}={0.4pt,bkLine}, rowsep=1.5pt, colsep=3pt, cells={font=\footnotesize, valign=m}, row{1}={valign=m}}
章 & 章名 & ★ & 路線 & 指導要領の学年・科目 & 範囲外 \\
{\sffamily\bfseries 6-0}\newline{\footnotesize\sffamily 序章} & 「次の瞬間を決めるルール」で未来を予想する & \kgcstars{2}{2} & \kgcroute{1} & 大学 & \kgcout{a} \\
{\sffamily\bfseries 6-1} & 微分方程式ってなに？ & \kgcstars{2}{2} & \kgcroute{1} & 大学 & \kgcout{a} \\
{\sffamily\bfseries 6-2} & 表計算で解こう：オイラー法 & \kgcstars{2}{2} & \kgcroute{1} & 数B〜大学 & \kgcout{p} \\
{\sffamily\bfseries 6-3} & 積分するだけで解ける形 & \kgcstars{3}{3} & \kgcroute{2} & 数Ⅲ〜大学 & \kgcout{p} \\
{\sffamily\bfseries 6-4} & 自分に比例してへる：冷めるお湯 & \kgcstars{3}{3} & \kgcroute{3} & 大学 & \kgcout{a} \\
{\sffamily\bfseries 6-5} & 空気抵抗と終端速度 & \kgcstars{2}{3} & \kgcroute{1} & 大学 & \kgcout{a} \\
{\sffamily\bfseries 6-6} & ばねと振り子 & \kgcstars{3}{3} & \kgcroute{4} & 大学 & \kgcout{a} \\
{\sffamily\bfseries 6-7} & ゆれながら止まる：減衰振動 & \kgcstars{3}{4} & \kgcroute{4} & 大学 & \kgcout{a} \\
{\sffamily\bfseries 6-8} & 2乗してマイナスになる数：虚数 & \kgcstars{4}{4} & \kgcroute{4} & 数Ⅱ & \kgcout{p} \\
{\sffamily\bfseries 6-9} & 複素平面：かけ算は回転と拡大 & \kgcstars{4}{4} & \kgcroute{4} & 数C &  \\
{\sffamily\bfseries 6-10} & オイラーの公式 & \kgcstars{4}{5} & \kgcroute{4} & 大学 & \kgcout{a} \\
{\sffamily\bfseries 6-11} & 双曲線関数 cosh・sinh & \kgcstars{4}{5} & \kgcroute{3} & 大学 & \kgcout{a} \\
{\sffamily\bfseries 6-12} & 減衰振動を式で解く：特性方程式 & \kgcstars{5}{5} & \kgcroute{4} & 大学 & \kgcout{a} \\
{\sffamily\bfseries 6-13} & ゆらし続けると：強制振動と共振 & \kgcstars{5}{5} & \kgcroute{4} & 大学 & \kgcout{a} \\
{\sffamily\bfseries 6-14} & 微分方程式がいくつも：連立微分方程式 & \kgcstars{5}{5} & \kgcroute{5} & 大学 & \kgcout{a} \\
{\sffamily\bfseries 6-15} & まとめ：物理は「微分方程式＋はじめの状態」で書ける & \kgcstars{2}{3} & \kgcroute{1} & 大学 & \kgcout{a} \\
\end{longtblr}
\Needspace{8\baselineskip}\kgbsub{第7巻}{まとめて計算する　行列}\par\nobreak{\small\noindent 学年の目安：高校3年（発展）〜大学1年　章数：12}\par\nobreak\vspace{3pt}
\begin{longtblr}[theme=kgc, caption={第7巻　まとめて計算する}, entry=none, label=none]{colspec={Q[l,wd=14mm] X[l] Q[l,wd=17mm] Q[l,wd=12.5mm] Q[l,wd=25mm] Q[c,wd=10mm]}, width=\linewidth,
  rowhead=1, row{1}={bg=kgcHead, font=\sffamily\bfseries\footnotesize}, row{even[2]}={bg=kgcBand},
  hlines={0.4pt,bkLine}, vline{2-6}={0.4pt,bkLine}, rowsep=1.5pt, colsep=3pt, cells={font=\footnotesize, valign=m}, row{1}={valign=m}}
章 & 章名 & ★ & 路線 & 指導要領の学年・科目 & 範囲外 \\
{\sffamily\bfseries 7-0}\newline{\footnotesize\sffamily 序章} & 表のままで計算できたら？ & \kgcstars{3}{3} & \kgcroute{5} & 数C〜大学 & \kgcout{p} \\
{\sffamily\bfseries 7-1} & 数の表で計算する：行列 & \kgcstars{3}{3} & \kgcroute{5} & 数C〜大学 & \kgcout{p} \\
{\sffamily\bfseries 7-2} & 行列×ベクトル：形を変える機械 & \kgcstars{4}{4} & \kgcroute{5} & 大学 & \kgcout{a} \\
{\sffamily\bfseries 7-3} & 行列×行列：変換をつなぐ & \kgcstars{4}{4} & \kgcroute{5} & 大学 & \kgcout{a} \\
{\sffamily\bfseries 7-4} & 行列式：面積が何倍になるか & \kgcstars{4}{4} & \kgcroute{5} & 大学 & \kgcout{a} \\
{\sffamily\bfseries 7-5} & 逆行列で連立方程式を解く & \kgcstars{4}{4} & \kgcroute{5} & 大学 & \kgcout{a} \\
{\sffamily\bfseries 7-6} & 回転行列：くるっと回す & \kgcstars{4}{4} & \kgcroute{4} & 数Ⅱ〜大学 & \kgcout{p} \\
{\sffamily\bfseries 7-7} & 3次元の回転と座標の変換 & \kgcstars{5}{5} & \kgcroute{6} & 大学 & \kgcout{a} \\
{\sffamily\bfseries 7-8} & 固有値と固有ベクトル：向きの変わらない方向 & \kgcstars{5}{5} & \kgcroute{5} & 大学 & \kgcout{a} \\
{\sffamily\bfseries 7-9} & 対角化：まっすぐな方向に分けて考える & \kgcstars{5}{5} & \kgcroute{5} & 大学 & \kgcout{a} \\
{\sffamily\bfseries 7-10} & 行列のくり返し：漸化式とオイラー法を行列で & \kgcstars{5}{5} & \kgcroute{5} & 数B〜大学 & \kgcout{p} \\
{\sffamily\bfseries 7-11}\newline{\footnotesize\sffamily 読みもの} & 行列の使いみち & \kgcstars{3}{5} & \kgcroute{5} & 大学 & \kgcout{a} \\
\end{longtblr}
\Needspace{8\baselineskip}\kgbsub{第8巻}{空間に広がる量　ベクトル解析}\par\nobreak{\small\noindent 学年の目安：大学1年〜2年　章数：15}\par\nobreak\vspace{3pt}
\begin{longtblr}[theme=kgc, caption={第8巻　空間に広がる量}, entry=none, label=none]{colspec={Q[l,wd=14mm] X[l] Q[l,wd=17mm] Q[l,wd=12.5mm] Q[l,wd=25mm] Q[c,wd=10mm]}, width=\linewidth,
  rowhead=1, row{1}={bg=kgcHead, font=\sffamily\bfseries\footnotesize}, row{even[2]}={bg=kgcBand},
  hlines={0.4pt,bkLine}, vline{2-6}={0.4pt,bkLine}, rowsep=1.5pt, colsep=3pt, cells={font=\footnotesize, valign=m}, row{1}={valign=m}}
章 & 章名 & ★ & 路線 & 指導要領の学年・科目 & 範囲外 \\
{\sffamily\bfseries 8-0}\newline{\footnotesize\sffamily 序章} & 空間のどこにも矢印がある & \kgcstars{3}{3} & \kgcroute{6} & 大学 & \kgcout{a} \\
{\sffamily\bfseries 8-1} & 地図と等高線：2つ以上の変数の関数 & \kgcstars{4}{4} & \kgcroute{1} & 大学 & \kgcout{a} \\
{\sffamily\bfseries 8-2} & 偏微分：ほかを止めて1つだけ動かす & \kgcstars{4}{4} & \kgcroute{1} & 大学 & \kgcout{a} \\
{\sffamily\bfseries 8-3} & 全微分：少し動いたときの変化 & \kgcstars{5}{5} & \kgcroute{1} & 大学 & \kgcout{a} \\
{\sffamily\bfseries 8-4} & いちばん急な向き：勾配 grad & \kgcstars{5}{5} & \kgcroute{1} & 大学 & \kgcout{a} \\
{\sffamily\bfseries 8-5} & ベクトル場 & \kgcstars{4}{4} & \kgcroute{6} & 大学 & \kgcout{a} \\
{\sffamily\bfseries 8-6} & 面や立体でたし合わせる：重積分 & \kgcstars{5}{5} & \kgcroute{2} & 大学 & \kgcout{a} \\
{\sffamily\bfseries 8-7} & 道にそってたす：線積分 & \kgcstars{5}{5} & \kgcroute{2} & 数Ⅲ〜大学 & \kgcout{p} \\
{\sffamily\bfseries 8-8} & 面を通りぬける量：面積分と流束 & \kgcstars{5}{5} & \kgcroute{2} & 大学 & \kgcout{a} \\
{\sffamily\bfseries 8-9} & わき出しと吸いこみ：発散 div & \kgcstars{5}{5} & \kgcroute{6} & 大学 & \kgcout{a} \\
{\sffamily\bfseries 8-10} & うず：回転 rot & \kgcstars{5}{5} & \kgcroute{6} & 大学 & \kgcout{a} \\
{\sffamily\bfseries 8-11} & 中と外をつなぐ：ガウスの定理とストークスの定理 & \kgcstars{5}{5} & \kgcroute{2} & 大学 & \kgcout{a} \\
{\sffamily\bfseries 8-12} & ならす働き：ラプラシアン & \kgcstars{5}{5} & \kgcroute{1} & 大学 & \kgcout{a} \\
{\sffamily\bfseries 8-13} & 熱と波の伝わり方：偏微分方程式 & \kgcstars{5}{5} & \kgcroute{1} & 大学 & \kgcout{a} \\
{\sffamily\bfseries 8-14}\newline{\footnotesize\sffamily 読みもの} & div と rot で書かれた式を読もう & \kgcstars{5}{5} & \kgcroute{6} & 大学 & \kgcout{a} \\
\end{longtblr}
\Needspace{8\baselineskip}\kgbsub{第9巻}{波と変化を読みとく　フーリエ解析とラプラス変換}\par\nobreak{\small\noindent 学年の目安：大学2年　章数：14}\par\nobreak\vspace{3pt}
\begin{longtblr}[theme=kgc, caption={第9巻　波と変化を読みとく}, entry=none, label=none]{colspec={Q[l,wd=14mm] X[l] Q[l,wd=17mm] Q[l,wd=12.5mm] Q[l,wd=25mm] Q[c,wd=10mm]}, width=\linewidth,
  rowhead=1, row{1}={bg=kgcHead, font=\sffamily\bfseries\footnotesize}, row{even[2]}={bg=kgcBand},
  hlines={0.4pt,bkLine}, vline{2-6}={0.4pt,bkLine}, rowsep=1.5pt, colsep=3pt, cells={font=\footnotesize, valign=m}, row{1}={valign=m}}
章 & 章名 & ★ & 路線 & 指導要領の学年・科目 & 範囲外 \\
{\sffamily\bfseries 9-0}\newline{\footnotesize\sffamily 序章} & 音を分ける、変化を予想する & \kgcstars{3}{3} & \kgcroute{4} & 大学 & \kgcout{a} \\
{\sffamily\bfseries 9-1} & 波をたし合わせて形を作る：フーリエ級数 & \kgcstars{4}{4} & \kgcroute{4} & 大学 & \kgcout{a} \\
{\sffamily\bfseries 9-2} & 成分の取り出し方：直交性 & \kgcstars{5}{5} & \kgcroute{4} & 大学 & \kgcout{a} \\
{\sffamily\bfseries 9-3} & 複素数で書くフーリエ級数 & \kgcstars{5}{5} & \kgcroute{4} & 大学 & \kgcout{a} \\
{\sffamily\bfseries 9-4} & スペクトル：フーリエ変換 & \kgcstars{5}{5} & \kgcroute{7} & 大学 & \kgcout{a} \\
{\sffamily\bfseries 9-5} & 微分をかけ算に変える翻訳機：ラプラス変換 & \kgcstars{5}{5} & \kgcroute{3} & 大学 & \kgcout{a} \\
{\sffamily\bfseries 9-6} & 変換表と性質 & \kgcstars{5}{5} & \kgcroute{3} & 大学 & \kgcout{a} \\
{\sffamily\bfseries 9-7} & ラプラス変換で微分方程式を解く & \kgcstars{5}{5} & \kgcroute{3} & 数Ⅱ〜大学 & \kgcout{p} \\
{\sffamily\bfseries 9-8} & 伝達関数：入れると出てくるを s で表す & \kgcstars{5}{5} & \kgcroute{3} & 大学 & \kgcout{a} \\
{\sffamily\bfseries 9-9} & 極と安定性：s の値でふるまいがわかる & \kgcstars{5}{5} & \kgcroute{3} & 数Ⅱ〜大学 & \kgcout{p} \\
{\sffamily\bfseries 9-10} & 連立微分方程式を行列で：状態方程式 & \kgcstars{5}{5} & \kgcroute{5} & 大学 & \kgcout{a} \\
{\sffamily\bfseries 9-11} & 行列の指数関数と固有値 & \kgcstars{5}{5} & \kgcroute{5} & 数Ⅲ〜大学 & \kgcout{p} \\
{\sffamily\bfseries 9-12}\newline{\footnotesize\sffamily 読みもの} & フーリエとラプラスの使いみち & \kgcstars{3}{5} & \kgcroute{7} & 大学 & \kgcout{a} \\
{\sffamily\bfseries 9-13}\newline{\footnotesize\sffamily 終章} & 工学の数学の地図 & \kgcstars{3}{5} & \kgcallroutes{} & ― &  \\
\end{longtblr}
}{\typeout{BKPAGE missing guide-p2b}}
\IfFileExists{guide-p2a-src.tex}{%% guide-p2a-src.tex　ガイドブック　保護者・先生のための本（大人用）の「出典」（組 G-B、頭文字 kgb）。親ファイルで付録（guide-p2b）のあとに読む。
%% \kgbsec・\kgbsub は guide-p2a.tex で定義（\providecommand）。材料：tools/guide-shidou.md 1節。
%% P1・P2 の URL のドメインは、材料では「co.jP」と書かれているが、ドメイン名は大文字・小文字を区別しないので「co.jp」と書いた。
\clearpage\kgbsec{}{出典}\label{kgb-src}
\begingroup
\urlstyle{same}\def\UrlFont{\sffamily\footnotesize}%
\Urlmuskip=0mu plus 1mu\relax
\expandafter\def\expandafter\UrlBreaks\expandafter{\UrlBreaks\do\a\do\b\do\c\do\d\do\e\do\f\do\g\do\h\do\i\do\j\do\k\do\l\do\m\do\n\do\o\do\p\do\q\do\r\do\s\do\t\do\u\do\v\do\w\do\x\do\y\do\z\do\0\do\1\do\2\do\3\do\4\do\5\do\6\do\7\do\8\do\9\do\_\do\-}%
この本で照らし合わせた資料です（2026年10月30日に調べました）。小学校・中学校・高等学校の指導要領の本文は、文部科学省の「学習指導要領コード」（告示本文を項目ごとに収めたデータ。M13）で、項目ごとに確かめました。小学校・中学校の本文は、教科書会社の新旧対照表（P1〜P3）と解説（M4・M5）とも照らしました。
\def\kgbsrc#1#2#3{\par\addvspace{3pt}\noindent\hangindent=3.2\zw\hangafter=1\makebox[3.2\zw][l]{\sffamily\bfseries\small #1}{\small #2}\par\noindent\hspace{3.2\zw}\parbox[t]{\dimexpr\linewidth-3.2\zw\relax}{\url{#3}}\par}%
\kgbsub{1}{文部科学省の資料（学習指導要領の本文・解説など）}
\kgbsrc{M1}{平成29・30・31年改訂学習指導要領（本文、解説）}{https://www.mext.go.jp/a_menu/shotou/new-cs/1384661.htm}
\kgbsrc{M2}{小学校学習指導要領（平成29年告示）}{https://www.mext.go.jp/content/1413522_001.pdf}
\kgbsrc{M3}{中学校学習指導要領（平成29年告示）}{https://www.mext.go.jp/content/1413522_002.pdf}
\kgbsrc{M4}{小学校学習指導要領（平成29年告示）解説 算数編}{https://www.mext.go.jp/content/20211102-mxt_kyoiku02-100002607_04.pdf}
\kgbsrc{M5}{中学校学習指導要領（平成29年告示）解説 数学編}{https://www.mext.go.jp/component/a_menu/education/micro_detail/__icsFiles/afieldfile/2019/03/18/1387018_004.pdf}
\kgbsrc{M6}{高等学校学習指導要領 比較対照表（数学）}{https://www.mext.go.jp/component/a_menu/education/micro_detail/__icsFiles/afieldfile/2018/07/13/1407085_5.pdf}
\kgbsrc{M7}{高等学校学習指導要領（平成30年告示）解説 数学編・理数編}{https://www.mext.go.jp/content/20230217-mxt_kyoiku02-100002620_05.pdf}
\kgbsrc{M8}{平成29年改訂の小・中学校学習指導要領に関するQ\&A（算数、数学に関すること）}{https://www.mext.go.jp/content/1422350_001.pdf}
\kgbsrc{M9}{中学校学習指導要領の特例を定める件（平成29年文部科学省告示第94号）}{https://www.mext.go.jp/a_menu/shotou/new-cs/__icsFiles/afieldfile/2017/07/11/1387780_003_1.pdf}
\kgbsrc{M10}{高等学校数学科教材「行列入門」（文部科学省）}{https://www.mext.go.jp/content/20230828-mxt-kyoiku01_000250597_1.pdf}
\kgbsrc{M11}{中央教育審議会資料「標準授業時数の在り方について」}{https://www.mext.go.jp/content/20200205-mxt_kyoiku01-000004708_9.pdf}
\kgbsrc{M12}{令和7年度大学入学共通テストにおける新教育課程対応について（大学入試センター提出資料）}{https://www.mext.go.jp/content/20211014-mxt_daigakuc02-000018441_3.pdf}
\kgbsrc{M13}{学習指導要領コード（文部科学省。小学校・中学校は平成29年告示、高等学校は平成30年告示の本文を項目ごとに表示したもの）}{https://jp-cos.github.io/}
\kgbsrc{M14}{大学入試センター「出題方法等の予告」（令和7年度大学入学共通テスト、数学）}{https://www.dnc.ac.jp/albums/abm.php?d=743&f=abm00003201.pdf}
\kgbsub{2}{教科書会社の新旧対照表（告示の本文を写したもの）}
\kgbsrc{P1}{大日本図書「小学校算数 現行・新学習指導要領 教授用資料 新旧対照表」}{https://www.dainippon-tosho.co.jp/sansu/files/cos2017e_sansu_f.pdf}
\kgbsrc{P2}{大日本図書「中学校数学 現行・新学習指導要領 教授用資料 新旧対照表」}{https://www.dainippon-tosho.co.jp/math/files/cos2017j_sugaku_f.pdf}
\kgbsrc{P3}{日本文教出版「小学校算数 新旧対照表」}{https://www.nichibun-g.co.jp/textbooks/sansu/download/new_cs/1705_newcs_sansu.pdf}
\kgbsub{3}{大学のシラバス・資料（8章）}
\kgbsrc{U1}{東北大学 工学部「工業数学Ⅱ」（2024）}{https://qsl.cds.tohoku.ac.jp/qsl/syllabus/display/tb44011}
\kgbsrc{U2}{東北大学 工学部「数学Ⅰ」（2024）}{https://qsl.cds.tohoku.ac.jp/qsl/syllabus/display/tb13011}
\kgbsrc{U3}{東北大学「線形代数学B」（2024）}{https://qsl.cds.tohoku.ac.jp/qsl/syllabus/display/cb42201}
\kgbsrc{U4}{京都工芸繊維大学 設計工学域「線形代数学Ⅱ」（2026）}{https://www.kit.ac.jp/wp/wp-content/uploads/sy2026/12061030.pdf}
\kgbsrc{U5}{北海道大学「微分積分学Ⅰ」（2026）}{https://www.math.sci.hokudai.ac.jp/~furuhata/ed/syllabus/26Calculus1.pdf}
\kgbsrc{U6}{神戸大学 工学部 前期シラバス}{https://www.office.kobe-u.ac.jp/eng-ofc/kym/syllabus/09_bz_a.pdf}
\kgbsrc{U8}{東京農工大学 工学部 数学科目一覧}{https://web.tuat.ac.jp/~mathsci/syllabus.html}
\kgbsrc{U9}{北海道大学「確率・統計入門」（2018）}{https://www.math.sci.hokudai.ac.jp/~furuhata/ed/syllabus/18IntroductionToStatistics.pdf}
\endgroup
}{\typeout{BKPAGE missing guide-p2a-src}}

%% ===================== 奥付・表4 =====================
\makeatletter
\kgaColophon{物理を理解するための数学　ガイドブック}{保護者・先生のための本（大人用）}{第1版（第16回編集委員会）　2026年10月}%
  {シリーズは全9巻です。学習指導要領との対応は第3章〜第9章と付録A・Bにまとめました。各巻の問題の詳しい解き方は、今後刊行する問題集と指導書に掲載する予定です。お子さん向けには、子ども用の『ガイドブック　はじめに読む本』があります。}
\bkartpage{表4}{\kgaBackArt{本書の内容}{%
  \kgaBackItem{197}{1}{シリーズの概要と使い方}
  \kgaBackItem{188}{2}{編集方針と書き方の原則}
  \kgaBackItem{179}{3}{学習指導要領との関係・学年と巻の対応}
  \kgaBackItem{170}{6・7}{追加した内容と、含めていない内容・補い方}
  \kgaBackItem{161}{A}{付録：指導要領の項目からの索引、全152章の一覧}}%
  {学年・単元から\\巻と章が引けます}}
\makeatother
\end{document}
